\documentclass[11pt,a4paper]{article}
\usepackage[T1]{fontenc}
\usepackage[utf8]{inputenc}
\usepackage{lmodern,amsmath,amssymb,amsthm,mathtools,microtype,mathrsfs}
\usepackage[margin=25mm,top=23mm,bottom=25mm]{geometry}
\usepackage{booktabs,array,longtable,enumitem,xcolor,fancyhdr}
\usepackage[colorlinks=true,linkcolor=blue!45!black,urlcolor=blue!45!black,citecolor=blue!45!black]{hyperref}
\hypersetup{pdftitle={Reading the Alpoge--Claude RH and Dirichlet-GRH proof: working version 0.53},pdfauthor={ChatGPT},pdfsubject={Partial human-readable reconstruction with explicit arithmetic inputs},pdfkeywords={Riemann hypothesis, Lean, theta lifts, prime sums, Arakelov slopes}}
\setlength{\parindent}{0pt}
\setlength{\parskip}{5pt plus 1pt}
\setlist{topsep=3pt,itemsep=3pt,leftmargin=*}
\setlength{\headheight}{14pt}
\pagestyle{fancy}\fancyhf{}
\fancyhead[L]{\small RH and Dirichlet GRH: a reconstruction}
\fancyhead[R]{\small Working version 0.53}
\fancyfoot[C]{\thepage}
\renewcommand{\headrulewidth}{0.3pt}
\numberwithin{equation}{section}
\newtheorem{theorem}{Theorem}[section]
\newtheorem{proposition}[theorem]{Proposition}
\newtheorem{lemma}[theorem]{Lemma}
\newtheorem{corollary}[theorem]{Corollary}
\theoremstyle{remark}\newtheorem{remark}[theorem]{Remark}
\newcommand{\Rea}{\operatorname{Re}}
\newcommand{\HS}{\mathrm{HS}}
\newcommand{\tr}{\operatorname{tr}}
\newcommand{\codim}{\operatorname{codim}}
\newcommand{\covol}{\operatorname{covol}}
\newcommand{\OK}{\mathcal O_K}
\newcommand{\src}[2]{\href{https://github.com/CocoSgt/riemann-hypothesis-lean/blob/7cbf7a0c32fca1badf05931879a37c3f8a2b03b4/src/#1}{\texttt{#2}}}
\newcommand{\status}[1]{\textbf{Status:} #1}
\begin{document}
\begin{center}
{\LARGE\bfseries RH and Dirichlet GRH: new human proofs}\\[7pt]
{\large Version 0.53 --- 23 November 2026}\\[5pt]
{\large Noetherian models, finite points, compatible measures and Wach lattices}
\end{center}
\textbf{Scope.} This edition contains the current audit and all new chapters.
It is a partial human reconstruction of the accepted formal proof.
Earlier named claims are in the cumulative manuscript and preserved components.
Each result below states its original inputs and exact remaining producers.
\setcounter{tocdepth}{1}
\tableofcontents
\clearpage
\begingroup\emergencystretch=3em
\section{Current proof obligations and completed suppliers: version 0.53}
\label{sec:dependency-audit-v53}

This is a continuation of the human reconstruction, with the formal
proof accepted as correct. The current selected sufficient route has
121 nodes and 161 edges. There are 24 open whole obligations:
17 class C and 7 class B. Of these, 17 are shared by ordinary
RH and unrestricted Dirichlet GRH; 7 are additional Dirichlet-GRH
inputs, and none are RH-only. The earlier complete
\texttt{GLOBAL\_\allowbreak LOCAL\_\allowbreak GAIN} closure remains in force.

The four new chapters contain 82 numbered claims and bring the
cumulative inventory to 3361. Their endpoints are original
mathematical subtheorems with their hypotheses printed in full:
the algebraic Noetherian section-cube model, the finite-point and
abstract propagation implications, the global and local compatible
measure constructors with \texttt{Rescale\allowbreak Step3}, and the generic
Wach-lattice constructors. These results supply specific inputs to
the next consumers. They do not prove their own remaining antecedents.

In particular, \texttt{Fin\allowbreak Pt\allowbreak Of} consumes a supplied space and
\texttt{Off\allowbreak Fin}; the actual rank-specific propagation must produce
that latter assertion. The graded cube now has its Noetherian model,
while the comparison with the actual geometric section cube still
requires localization and gluing of sections. The measure construction
keeps the actual setting and one character before every word and
place, while arithmetic quotient-volume identities and transport
remain distinct. The Wach construction is inside the supplied
\'etale module; its all-layer return retains three precise period-ring
suppliers, followed by the separate Herr and quantitative comparison
work. The original \texttt{Layers} constructor adds no unspecified
transition condition between outer indices.

A class-B claim below is a written human deduction at its stated scope.
Classifications and counts are an inventory of this selected route.
They are not a claim of minimality or a measure of the percentage
completed. All earlier proof records and historical text are retained;
the following register governs the current remaining scope.

\subsection{Complete current whole-obligation register}


\paragraph{1. GLOBAL\_\allowbreak DENOMINATOR — C — Actual LinkN local table and binding producer}

Closed(Des $\mathcal C$ 3) LinkN: the actual etaSum is bounded with exact dZ0 and same-host vvTerm, subject only to the literal ShaFinOdd and summability guards.

\emph{Shared objects and choice order.} one FamP; one class $\mathcal A$ and its target-compatible subclass; one coefficient field per family datum; one x and i at each application; original selected rows, words and levels; thresholds chosen before the common family; actual D.nd; actual etaJ and Aj; sha\_\allowbreak eq

\emph{Completed inputs (version 0.53).} WLD1-WLD26 (full generic Wach lattice constructors; exact conditional three-supplier all-layer return), SSC1-SSC17 (full original Selmer SignatureChoice and family setup retaining only coherent Wach comparison), DFD1-DFD15, DNE1-DNE14 (full original ConclF-to-LinkN consumer)

\emph{Remaining producer scope.} WLD completes the full original adic-coordinate, denominator-matrix, base-change lattice, restricted-projectivity and local-prism value producers, and both generic WachLattice.build/build\_\allowbreak prism implications. The actual lattice is the inverse image of one prismatic image inside the supplied etale module; its original Frobenius, group action, trivial-line and all lattice-map fields are proved together. Full all-layer WachEx still needs the actual Wach prism, ioE injectivity/unit reflection and homogeneous cyclotomic-kernel-fixed Frobenius matrix descent at the original scopes. After those, faithful invariants, etale bases, Herr class maps and quantitative comparison still supply the complete HW/Fon/Cmp triple. No additional inter-layer transition system is required by the original Layers assembly. SSC already supplies the full Selmer SignatureChoice and DFD12 the exact Bhatt-Scholze coefficient conversion. Then prove every original local row/table bound on one common family for ConclF, retaining ShaFinOdd, exceptional sets, vvTerm summability and uniform thresholds. DNE supplies ConclF-to-AdmF-to-LinkN with K=2000(|Ca|+1)\textasciicircum{}3 fixed before the family; GLOBAL\_\allowbreak DENOMINATOR remains open.

\emph{Source.} \href{https://github.com/CocoSgt/riemann-hypothesis-lean/blob/7cbf7a0c32fca1badf05931879a37c3f8a2b03b4/src/Formalism/SharpRestriction/GlobalTheoremInputs.lean#L45-L48}{LinkN}.

\emph{Source.} \href{https://github.com/CocoSgt/riemann-hypothesis-lean/blob/7cbf7a0c32fca1badf05931879a37c3f8a2b03b4/src/Formalism/SharpRestriction/GlobalTheoremInputs.lean#L40-L44}{LinkN\_\allowbreak at}.

\emph{Source.} \href{https://github.com/CocoSgt/riemann-hypothesis-lean/blob/7cbf7a0c32fca1badf05931879a37c3f8a2b03b4/src/Development/AutomorphicArithmetic/CharacterBoxes/DivinNOfHypV.lean#L117-L134}{divinN\_\allowbreak of\_\allowbreak hypV}.



\paragraph{2. GLOBAL\_\allowbreak RETAINED\_\allowbreak DEFECT — C — Original retained-row defect bound LinkIRR}

Rows.DelpNonneg and delp$\le$exp(-x\textasciicircum{}(ez/7)) uniformly for every late actual row datum. CURRENT VERSION0.47 QUALIFICATION: Finish the same-D/common-family adelic setting after ASG: normalized nonzero Bessel W and decay, measured-field block, original Siegel-Weil rescaling and measure/character ties, Bessel uniqueness/local ties, then the original AutTie/RowDictTie/LabTie. RowsFinDim is now proved universally by RFD. Supply the actual CountIdempotentsClass theorem or its three sufficient arithmetic cylinder counts, with one simultaneous family threshold; reuse the completed ARO structure. Completed new original producers: ARO1-ARO30, RFD1-RFD9, ASG1-ASG18. Any earlier narrative naming these completed producers as outstanding is historical; this qualification is the current scope authority. CURRENT VERSION0.48 QUALIFICATION: Finish the same-D/common-family adelic setting after ASG and AMD: normalized nonzero Bessel W and decay, original orbital/fiber and Siegel-Weil measure/character ties, Bessel uniqueness/local ties, and the full original transport theorem, then AutTie/RowDictTie/LabTie. The complete MExists measured-data block and unipotent covolume normalization are now supplied, and universal RowsFinDim remains available. Supply the actual CountIdempotentsClass theorem or its three sufficient arithmetic cylinder counts with one simultaneous family threshold. Completed current producers or conditional consumers: ARO1-ARO30, RFD1-RFD9, ASG1-ASG18, AMD1-AMD17. Earlier frontier descriptions are historical; this qualification is the current scope authority. Original whole-contract premises and quantifiers are retained. CURRENT VERSION0.49 QUALIFICATION: Prove the exact original BesselNVGap nonzero integral on its supplied discrete-series, compact-representation and character data. CoefWeightedLoc, WlocExists with its original nonvanishing antecedent, WFromLoc, the normalized global product and all-later-measured-block decay are complete. On the inherited Forms0 scope this leaves BesselNVGap for ConcreteSettingExists. Continue original orbital/fiber and Siegel-Weil measure/character ties, the required Bessel uniqueness/local ties and the full original transport theorem, then AutTie/RowDictTie/LabTie. Supply the actual CountIdempotentsClass theorem or its three sufficient arithmetic cylinder counts with one simultaneous family threshold. Completed current producers or conditional consumers: ARO1-ARO30, RFD1-RFD9, ASG1-ASG18, AMD1-AMD17, BWP1-BWP24 at their exact BesselNVGap/Forms0 scope. Earlier frontier descriptions are historical; this qualification is the current scope authority. Original whole-contract premises and quantifiers are retained.

\emph{Shared objects and choice order.} one FamP; one class $\mathcal A$ and its target-compatible subclass; one coefficient field per family datum; one x and i at each application; original selected rows, words and levels; thresholds chosen before the common family; AllRows A; LabelCountGoodInjective; AllJoinA; CountAllRows

\emph{Completed inputs (version 0.53).} SSC full Selmer-signature comparison, DFD common-family subresults, DNE full separate denominator consumer, WRS1-WRS10 full selected-class local gain, Inherited first DivDef3 and row-pair integrability

\emph{Remaining producer scope.} Complete the original F3 local coefficient bridges and F4 common-overfield readings, with valuation/norm/GoodCoef/iq and the entire coefficient field, and actual complex Hecke standard/no-Jordan plus arithmetic character data on the same form datum. Construct the actual base rows, host/frame lengths, nonzero chains and group ties, MultEq support, eigensystem/carrier/parts, all-embedding logLam, HD/d3b and NEb. Keep the stronger second DivRows3 scope distinct from the completed first DivDef3 and row-pair integrability. The full arithmetic entry and conditional common-frame construction do not by themselves install HOST\_\allowbreak NUMBER\_\allowbreak S, HOST\_\allowbreak FRAME\_\allowbreak LENGTH, TASK\_\allowbreak HOST\_\allowbreak VAN\_\allowbreak E\_\allowbreak INSTALL, TASK\_\allowbreak HOST\_\allowbreak XPP\_\allowbreak INSTALL or TASK\_\allowbreak COUNT\_\allowbreak V\_\allowbreak INSTALL. Continue the distinct retained-defect, period-height and normalized spectral producers, and the original compatible second-layer selection on one common arithmetic/geometric family. GLOBAL\_\allowbreak LOCAL\_\allowbreak GAIN remains fully closed at its original v0.49 scope. Complete local comparison theorems do not themselves construct the whole common family. Common design existence, NBUILT compatibility, original true-column records, shared fields/embeddings, Fits and eta agreement, simultaneous rows and eventual-family thresholds remain explicit. A stronger historical WPt/TSel construction may serve other producers but is not needed to reopen the proved selected-class local-gain estimate.

\emph{Source.} \href{https://github.com/CocoSgt/riemann-hypothesis-lean/blob/7cbf7a0c32fca1badf05931879a37c3f8a2b03b4/src/Formalism/SharpRestriction/GlobalTheoremInputs.lean#L55-L58}{LinkIRR}.

\emph{Source.} \href{https://github.com/CocoSgt/riemann-hypothesis-lean/blob/7cbf7a0c32fca1badf05931879a37c3f8a2b03b4/src/Development/AbelianVarietyTheory/Proofs/GlobalInterpolation/LinkIRROfJoinEx.lean#L17-L38}{linkIRR\_\allowbreak of\_\allowbreak join\_\allowbreak ex}.

\emph{Source.} \href{https://github.com/CocoSgt/riemann-hypothesis-lean/blob/7cbf7a0c32fca1badf05931879a37c3f8a2b03b4/src/Development/AbelianVarietyTheory/Proofs/GlobalInterpolation/LinkIRROfJoinEx.lean#L143-L166}{link\_\allowbreak irr}.



\paragraph{3. GLOBAL\_\allowbreak PERIOD\_\allowbreak HEIGHT — C — Full actual period and height comparison LinkDIVt}

Closed conditional TuSize -> DGamSize -> (LinkEta -> DIVinN -> LinkDIV), with every original LinkDIV\_\allowbreak at clause.

\emph{Shared objects and choice order.} one FamP; one class $\mathcal A$ and its target-compatible subclass; one coefficient field per family datum; one x and i at each application; original selected rows, words and levels; thresholds chosen before the common family; same RR target; same Bu/besS; same Nu and Bk

\emph{Source.} \href{https://github.com/CocoSgt/riemann-hypothesis-lean/blob/7cbf7a0c32fca1badf05931879a37c3f8a2b03b4/src/Formalism/SharpRestriction/GlobalTheoremInputs.lean#L77-L90}{LinkDIV\_\allowbreak at}.

\emph{Source.} \href{https://github.com/CocoSgt/riemann-hypothesis-lean/blob/7cbf7a0c32fca1badf05931879a37c3f8a2b03b4/src/Formalism/SharpRestriction/GlobalTheoremInputs.lean#L94-L96}{LinkDIVt}.

\emph{Source.} \href{https://github.com/CocoSgt/riemann-hypothesis-lean/blob/7cbf7a0c32fca1badf05931879a37c3f8a2b03b4/src/Development/AbelianVarietyTheory/Proofs/GlobalInterpolation/LinkBrest.lean#L7-L25}{GlobalTheory.Top.link\_\allowbreak brest}.

\emph{Source.} \href{https://github.com/CocoSgt/riemann-hypothesis-lean/blob/7cbf7a0c32fca1badf05931879a37c3f8a2b03b4/src/Development/AbelianVarietyTheory/Proofs/GlobalInterpolation/WireSel.lean#L11-L41}{GlobalTheory.Top.wire\_\allowbreak sel}.



\paragraph{4. GLOBAL\_\allowbreak SPECTRAL — C — Original normalized kernel/Gram estimates LinkS}

Closed TuSize -> LinkS with exact Gamma normalization, off-diagonal Hilbert-Schmidt bound, diagonal error and Dall/Dstar estimate, errors exp(-x\textasciicircum{}(epsilon/8)). CURRENT VERSION0.48 QUALIFICATION: Complete the original same-family spectral/Bessel and Siegel-Weil inputs with all measure, character, transport and common-object ties. The complete MExists measured block and prescribed-measure covolume normalization are supplied; these are not the full normalized spectral theorem. Completed current producers or conditional consumers: ASG1-ASG18, AMD1-AMD17, GCF1-GCF18 at their stated scopes. Earlier frontier descriptions are historical; this qualification is the current scope authority. Original whole-contract premises and quantifiers are retained. CURRENT VERSION0.49 QUALIFICATION: Complete the original same-family spectral/Bessel and Siegel-Weil inputs with every measure, character, transport and common-object tie. The normalized packet and decay are now constructed conditional on the exact BesselNVGap, and the analytic setting uses the original supplied Forms0 scope. These do not imply the full normalized spectral theorem. Completed current producers or conditional consumers: ASG1-ASG18, AMD1-AMD17, GCF1-GCF18 at their stated scopes, BWP1-BWP24 at their exact BesselNVGap/Forms0 scope, WRS1-WRS10 original local gain. Earlier frontier descriptions are historical; this qualification is the current scope authority. Original whole-contract premises and quantifiers are retained.

\emph{Shared objects and choice order.} one FamP; one class $\mathcal A$ and its target-compatible subclass; one coefficient field per family datum; one x and i at each application; original selected rows, words and levels; thresholds chosen before the common family; actual KD; same S.Cu spaces; GamOf; Dstar; Dall

\emph{Completed inputs (version 0.53).} GMR1-GMR19 (full global/local measure suppliers and RescaleStep3 on the actual setting), BQS1-BQS23 (full QuotP/QuotM, StepStatement and BesselNVGap; setting at Forms0), BST1-BST11, BSE1-BSE9, BNV1-BNV23, BSP1-BSP10, BWP1-BWP24 at original Forms0 scope, AMD1-AMD24

\emph{Remaining producer scope.} GMR completes the full original global stabilizer-measure constructor, normalized local frame-stabilizer constructor and RescaleStep3. The proof constructs continuous Gram fiber families, proves uniqueness at every parameter including singular ones, computes exact good-place masses and proves a strictly positive convergent product. One standard character is fixed before all words, places and auxiliary tuples; the actual setting, supplied bridge and infinity measures are preserved. The original SiegelWeilUnitaryConvergent premise is retained; its arithmetic theta/quotient-volume identities are not claimed proved by the measure construction. BQS and the inherited chain already give full BesselNVGap. Forms0 and the all-family Witt/genus/level/intertwiner Transport/LevelMatch suppliers remain, as do Bessel uniqueness, local volume identities and product ties, AutTie/RowDictTie/LabTie, simultaneous arithmetic cylinder counts and the normalized GLOBAL\_\allowbreak SPECTRAL estimate.

\emph{Source.} \href{https://github.com/CocoSgt/riemann-hypothesis-lean/blob/7cbf7a0c32fca1badf05931879a37c3f8a2b03b4/src/Formalism/SharpRestriction/GlobalTheoremInputs.lean#L145-L149}{LinkS\_\allowbreak at}.

\emph{Source.} \href{https://github.com/CocoSgt/riemann-hypothesis-lean/blob/7cbf7a0c32fca1badf05931879a37c3f8a2b03b4/src/Formalism/SharpRestriction/GlobalTheoremInputs.lean#L152-L157}{LinkS}.

\emph{Source.} \href{https://github.com/CocoSgt/riemann-hypothesis-lean/blob/7cbf7a0c32fca1badf05931879a37c3f8a2b03b4/src/Development/AbelianVarietyTheory/Proofs/GlobalInterpolation/LinkAlt.lean#L16-L50}{GlobalTheory.Top.link\_\allowbreak alt}.



\paragraph{5. AUT\_\allowbreak UNRAM\_\allowbreak COEF — C — Actual unramified-coefficient automorphy}

For each actual Dat with coefficient ramification one and fixed iota, choose one marked g1 supplying FinA, AutB and LemW together. CURRENT VERSION0.48 QUALIFICATION: Construct the actual AutomorphyUnramifiedCoefficients package: for every original Dat with coefficient ramification one and fixed iota, choose one marked g1 that supplies FinA, AutB and LemW together. The original InpZ record, full RigII, CNR3, finite PropGK consumer and now the entire OneChar implication are supplied. Continue RH Cases I/V and the independent additional Dirichlet-GRH CaseIII/promodularity, SeedIV, tensor-factor and rank-four-chain tasks. Completed current producers or conditional consumers: OCH1-OCH30 (the full consumer implication), KIZ original InpZ, KTR1-KTR13, RNC16, PGK1-PGK6. Earlier frontier descriptions are historical; this qualification is the current scope authority. Original whole-contract premises and quantifiers are retained. CURRENT VERSION0.49 QUALIFICATION: Construct one common GdEnd for each original datum and fixed iota: the same marked g1 must supply FinEnd, AutEnd3 and AutEnd2. UCE now proves all ordinary local/curve/coefficient consumers and the exact implication to FinA/AutB/LemW, including arbitrary entry maps and every original outer-corner extension. The producer still needs the actual common abelian-CM GdFrameEx and LinkedEndE with both rank-two and rank-three runs, full ModUAlt/PfromAut and the two actual seeds, followed by the original local/automorphic descent. At this pin the alternative upstream route gives GdEndAlt; its additional local/descent fields must be matched to the actual original points, so no name-only GdEndAlt-to-GdEnd identification is allowed. Keep the supplied InpZ, RigII, CNR3, finite PropGK and whole OneChar implication. Continue RH Cases I/V and the independent additional Dirichlet-GRH CaseIII/promodularity, SeedIV, tensor-factor and rank-four-chain tasks. Completed current producers or conditional consumers: UCE1-UCE19 (complete GdEnd consumer), OCH1-OCH30, KIZ original InpZ, KTR1-KTR13, RNC16, PGK1-PGK6. Earlier frontier descriptions are historical; this qualification is the current scope authority. Original whole-contract premises and quantifiers are retained.

\emph{Completed inputs (version 0.53).} ACP1-ACP19 (full FinPtOf and abstract propagation principles; rank-specific suppliers retained), AFE1-AFE21 (full Gen entries, conditional frame joins, NRep descent and FinPt-to-FinEnd), LFC1-LFC23, CAF1-CAF11, MFC1-MFC14, UCE1-UCE19 (full original GdEnd consumer), OCH1-OCH30, KIZ original InpZ, KTR1-KTR13, RNC16, PGK1-PGK6

\emph{Remaining producer scope.} ACP proves full FinPtOf at every original IsCD test ring from the supplied Space and literal OffFin, using the sufficient IsLiftA point realization through k+m\_\allowbreak A. It also proves full flag\_\allowbreak loc, o1, NicePrimeEx, LocProp, GlobProp, ChartOfFlat and GlobOfCharts. These abstract propagation and finite-point implications do not yet supply OffFin or the complete rank-two/rank-three propagation producers. Remaining arithmetic suppliers are the actual deformation charts and flatness, sharp component/dimension bounds, preconnected core and seed incidence, rank-two-to-three corner and exceptional ideal, and the actual automorphic seeds, patching/modularity and classicality outputs, and the original finPtDown descent from the extended coefficient frame before returning to the original FinEnd consumer. Keep separate rank-two/rank-three rings and ideals on the same AFE marked frame, with original P retained where used. The actual terminal route uses stronger GdEndAlt/HostEndAlt including WitA/WitAQ; ACP19 proves only their projection to the original endpoints. AFE full Gen/Host entries, common-frame joins, NRep descent and FinPt-to-FinEnd remain available, including Host ell=3; UCE supplies the original GdEnd consumer.

\emph{Source.} \href{https://github.com/CocoSgt/riemann-hypothesis-lean/blob/7cbf7a0c32fca1badf05931879a37c3f8a2b03b4/src/Development/AutomorphicArithmetic/AutomorphicVarieties/AutomorphicCases/AutomorphyUnramifiedCoefficients.lean#L24-L27}{AutomorphyUnramifiedCoefficients}.

\emph{Source.} \href{https://github.com/CocoSgt/riemann-hypothesis-lean/blob/7cbf7a0c32fca1badf05931879a37c3f8a2b03b4/src/Development/AutomorphicArithmetic/AutomorphicVarieties/AutomorphicCases/AutomorphyUnramifiedCoefficients.lean#L12-L13}{FinA}.

\emph{Source.} \href{https://github.com/CocoSgt/riemann-hypothesis-lean/blob/7cbf7a0c32fca1badf05931879a37c3f8a2b03b4/src/Development/AutomorphicArithmetic/AutomorphicVarieties/AutomorphicCases/AutomorphyUnramifiedCoefficients.lean#L15-L17}{AutB}.

\emph{Source.} \href{https://github.com/CocoSgt/riemann-hypothesis-lean/blob/7cbf7a0c32fca1badf05931879a37c3f8a2b03b4/src/Development/AutomorphicArithmetic/AutomorphicVarieties/AutomorphicCases/AutomorphyUnramifiedCoefficients.lean#L19-L22}{LemW}.



\paragraph{6. TASK\_\allowbreak CASE3\_\allowbreak AUT\_\allowbreak JOIN — B — Complete the original Case-III automorphy/image join}

Install the actual Case-III ProModAlt, SeedIVAlt and shared PMExAlt with the inherited CohomologyHyp/Ferrero-Washington/HFit/SeedII facts into the original whole Case-III automorphy/image conclusion.

\emph{Source.} \href{https://github.com/CocoSgt/riemann-hypothesis-lean/blob/7cbf7a0c32fca1badf05931879a37c3f8a2b03b4/src/Development/AutomorphicArithmetic/AutomorphicVarieties/ModularityArgument/ThresholdArguments/DenseThresholdCases/PromodularityCaseThree.lean#L43-L44}{AutomorphyRankFourDense}.

\emph{Source.} \href{https://github.com/CocoSgt/riemann-hypothesis-lean/blob/7cbf7a0c32fca1badf05931879a37c3f8a2b03b4/src/Development/AutomorphicArithmetic/AutomorphicVarieties/ModularityArgument/ThresholdArguments/DenseThresholdCases/CuFinal.lean#L44-L44}{main4\_\allowbreak alt}.



\paragraph{7. TASK\_\allowbreak CASE3\_\allowbreak PROMOD\_\allowbreak JOIN — C — Reconstruct the actual Case-III chain-to-promodularity proof}

Use the original tensor-factor tower, ChainRankFour, ChainDesc4, seeded rank-four chain reduction, actual PMExAlt and SeedIVAlt to construct ProModAlt on every original Ctx3 star point at the same fixed iota, preserving all finite-coefficient, inertia and support maps.

\emph{Source.} \href{https://github.com/CocoSgt/riemann-hypothesis-lean/blob/7cbf7a0c32fca1badf05931879a37c3f8a2b03b4/src/Development/AutomorphicArithmetic/AutomorphicVarieties/ModularityArgument/ThresholdArguments/DenseThresholdCases/PromodularityCaseThreeOfLemmas.lean#L131-L153}{UpperTheory.V.Cu.promodularity\_\allowbreak case\_\allowbreak three\_\allowbreak of\_\allowbreak chain\_\allowbreak reduction}.

\emph{Source.} \href{https://github.com/CocoSgt/riemann-hypothesis-lean/blob/7cbf7a0c32fca1badf05931879a37c3f8a2b03b4/src/Development/AutomorphicArithmetic/AutomorphicVarieties/ModularityArgument/ThresholdArguments/DenseThresholdCases/PromodularityCaseThree.lean#L40-L42}{PromodularityCaseThreeDense}.



\paragraph{8. SEED\_\allowbreak IV\_\allowbreak ALT — C — Actual Case-III seed with prechosen residual datum}

Base Ctx3/Four tower, iota and Amb -> w,s before every later Dies tower -> compatible finite-coefficient seed with exact residual support/root entries.

\emph{Source.} \href{https://github.com/CocoSgt/riemann-hypothesis-lean/blob/7cbf7a0c32fca1badf05931879a37c3f8a2b03b4/src/Development/AutomorphicArithmetic/AutomorphicVarieties/ModularityArgument/ThresholdArguments/DenseThresholdCases/PromodularityCaseThree.lean#L19-L25}{SeedIVAlt}.



\paragraph{9. TENSOR\_\allowbreak FACTOR\_\allowbreak TOWER — B — Actual tensor-factor tower and characters}

Original finite-field, tame, factor and character constraints; source-level construction is sizeable and has not been reconstructed here.

\emph{Source.} \href{https://github.com/CocoSgt/riemann-hypothesis-lean/blob/7cbf7a0c32fca1badf05931879a37c3f8a2b03b4/src/Development/AutomorphicArithmetic/AutomorphicVarieties/ModularityArgument/ThresholdArguments/DenseThresholdCases/PromodularityCaseThreeOfLemmas.lean#L131-L153}{UpperTheory.V.Cu.promodularity\_\allowbreak case\_\allowbreak three\_\allowbreak of\_\allowbreak chain\_\allowbreak reduction}.



\paragraph{10. CHAIN\_\allowbreak RANK\_\allowbreak FOUR — C — Actual SomeIn-to-AllIn propagation}

Same source Space and PM ideal, residual point Good4, original bound 30(C+1)\textasciicircum{}2 <= d; geometry must preserve all identities.

\emph{Source.} \href{https://github.com/CocoSgt/riemann-hypothesis-lean/blob/7cbf7a0c32fca1badf05931879a37c3f8a2b03b4/src/Development/AutomorphicArithmetic/AutomorphicVarieties/ModularityArgument/ThresholdArguments/DenseThresholdCases/PromodularityCaseThree.lean#L26-L32}{ChainRankFour}.



\paragraph{11. CHAIN\_\allowbreak DESC4 — B — Original residual chain descent}

Finite residual descent retaining source Chain4 compatibility.

\emph{Source.} \href{https://github.com/CocoSgt/riemann-hypothesis-lean/blob/7cbf7a0c32fca1badf05931879a37c3f8a2b03b4/src/Development/AutomorphicArithmetic/AutomorphicVarieties/ModularityArgument/ThresholdArguments/DenseThresholdCases/PromodularityCaseThreeOfLemmas.lean#L131-L153}{UpperTheory.V.Cu.promodularity\_\allowbreak case\_\allowbreak three\_\allowbreak of\_\allowbreak chain\_\allowbreak reduction}.



\paragraph{12. RANK4\_\allowbreak CHAIN\_\allowbreak REDUCTION — C — Seeded reduction to the original residual chain}

Select T1 before every later Deep tower; retain one compatible seed and Chain4 for every original residual lift over Q.

\emph{Source.} \href{https://github.com/CocoSgt/riemann-hypothesis-lean/blob/7cbf7a0c32fca1badf05931879a37c3f8a2b03b4/src/Development/AutomorphicArithmetic/AutomorphicVarieties/ModularityArgument/ThresholdArguments/DenseThresholdCases/PromodularityCaseThree.lean#L33-L39}{RankFourChainReductionDense}.

\emph{Source.} \href{https://github.com/CocoSgt/riemann-hypothesis-lean/blob/7cbf7a0c32fca1badf05931879a37c3f8a2b03b4/src/Development/AutomorphicArithmetic/AutomorphicVarieties/ModularityArgument/ThresholdArguments/DenseThresholdCases/CuPChainIV.lean#L16-L82}{UpperTheory.V.Cu.chainIV\_\allowbreak alt}.



\paragraph{13. TASK\_\allowbreak COUNT\_\allowbreak V\_\allowbreak INSTALL — B — Complete the original CountV finite assembly}

From exactly the listed source host/cohomology inputs, complete countV\_\allowbreak of\_\allowbreak finiteness\_\allowbreak duality\_\allowbreak euler\_\allowbreak host, including original host/frame formation, normalized length, the CaseV guard and R1<=2\textasciicircum{}c2.

\emph{Completed inputs (version 0.53).} ACP full finite-point and abstract propagation implications; actual rank-specific and counting inputs retained, AFE full Host entry/auxiliary frame/local witnesses including ell=3, LFC1-LFC23, CAF9-CAF11, MFC1-MFC14, Inherited CountV consumers

\emph{Remaining producer scope.} ACP proves full FinPtOf at every original IsCD test ring from the supplied Space and literal OffFin, using the sufficient IsLiftA point realization through k+m\_\allowbreak A. It also proves full flag\_\allowbreak loc, o1, NicePrimeEx, LocProp, GlobProp, ChartOfFlat and GlobOfCharts. These abstract propagation and finite-point implications do not yet supply OffFin or the complete rank-two/rank-three propagation producers. Remaining arithmetic suppliers are the actual deformation charts and flatness, sharp component/dimension bounds, preconnected core and seed incidence, rank-two-to-three corner and exceptional ideal, and the actual automorphic seeds, patching/modularity and classicality outputs, and the original finPtDown descent from the extended coefficient frame before returning to the original FinEnd consumer. Keep separate rank-two/rank-three rings and ideals on the same AFE marked frame, with original P retained where used. The actual terminal route uses stronger GdEndAlt/HostEndAlt including WitA/WitAQ; ACP19 proves only their projection to the original endpoints. AFE full Gen/Host entries, common-frame joins, NRep descent and FinPt-to-FinEnd remain available, including Host ell=3; UCE supplies the original GdEnd consumer. Complete the original F3 local coefficient bridges and F4 common-overfield readings, with valuation/norm/GoodCoef/iq and the entire coefficient field, and actual complex Hecke standard/no-Jordan plus arithmetic character data on the same form datum. Construct the actual base rows, host/frame lengths, nonzero chains and group ties, MultEq support, eigensystem/carrier/parts, all-embedding logLam, HD/d3b and NEb. Keep the stronger second DivRows3 scope distinct from the completed first DivDef3 and row-pair integrability. The full arithmetic entry and conditional common-frame construction do not by themselves install HOST\_\allowbreak NUMBER\_\allowbreak S, HOST\_\allowbreak FRAME\_\allowbreak LENGTH, TASK\_\allowbreak HOST\_\allowbreak VAN\_\allowbreak E\_\allowbreak INSTALL, TASK\_\allowbreak HOST\_\allowbreak XPP\_\allowbreak INSTALL or TASK\_\allowbreak COUNT\_\allowbreak V\_\allowbreak INSTALL.

\emph{Source.} \href{https://github.com/CocoSgt/riemann-hypothesis-lean/blob/7cbf7a0c32fca1badf05931879a37c3f8a2b03b4/src/Development/AutomorphicArithmetic/AutomorphicVarieties/HeightAndVanishing/Counts.lean#L69-L69}{CountV}.

\emph{Source.} \href{https://github.com/CocoSgt/riemann-hypothesis-lean/blob/7cbf7a0c32fca1badf05931879a37c3f8a2b03b4/src/Development/AutomorphicArithmetic/AutomorphicVarieties/ModularityArgument/ThresholdArguments/DenseThresholdCases/CountVOf.lean#L33-L60}{UpperTheory.V.UnitaryThreeInertia.countV\_\allowbreak of\_\allowbreak finiteness\_\allowbreak duality\_\allowbreak euler\_\allowbreak host}.



\paragraph{14. THM\_\allowbreak S6\_\allowbreak S — C — Original unitary-three ordinary support and automorphy}

Original host/frame and fixed iota -> one marked g1 -> every universal R,u -> weight-zero finiteness and every Type-I BPrime\_\allowbreak s point.

\emph{Source.} \href{https://github.com/CocoSgt/riemann-hypothesis-lean/blob/7cbf7a0c32fca1badf05931879a37c3f8a2b03b4/src/Development/AutomorphicArithmetic/AutomorphicVarieties/ModularityArgument/ThresholdArguments/DenseThresholdCases/ThmS6CoreS.lean#L17-L26}{ThmS6\_\allowbreak s}.

\emph{Source.} \href{https://github.com/CocoSgt/riemann-hypothesis-lean/blob/7cbf7a0c32fca1badf05931879a37c3f8a2b03b4/src/Development/AutomorphicArithmetic/AutomorphicVarieties/ModularityArgument/ThresholdArguments/DenseThresholdCases/ThmS6CoreS.lean#L12-L16}{BPrime\_\allowbreak s}.



\paragraph{15. HOST\_\allowbreak FRAME\_\allowbreak LENGTH — B — Original Case-V host/frame and normalized-length installation}

Actual HostData, frame, local-condition equality, abstract count and normalized original Selmer length must all refer to the same objects.

\emph{Completed inputs (version 0.53).} AFE/LFC full entry and conditional common frame, CAF and MFC, DNE weight/host-degree dictionary at supplied tables

\emph{Remaining producer scope.} Complete the original F3 local coefficient bridges and F4 common-overfield readings, with valuation/norm/GoodCoef/iq and the entire coefficient field, and actual complex Hecke standard/no-Jordan plus arithmetic character data on the same form datum. Construct the actual base rows, host/frame lengths, nonzero chains and group ties, MultEq support, eigensystem/carrier/parts, all-embedding logLam, HD/d3b and NEb. Keep the stronger second DivRows3 scope distinct from the completed first DivDef3 and row-pair integrability. The full arithmetic entry and conditional common-frame construction do not by themselves install HOST\_\allowbreak NUMBER\_\allowbreak S, HOST\_\allowbreak FRAME\_\allowbreak LENGTH, TASK\_\allowbreak HOST\_\allowbreak VAN\_\allowbreak E\_\allowbreak INSTALL, TASK\_\allowbreak HOST\_\allowbreak XPP\_\allowbreak INSTALL or TASK\_\allowbreak COUNT\_\allowbreak V\_\allowbreak INSTALL.

\emph{Source.} \href{https://github.com/CocoSgt/riemann-hypothesis-lean/blob/7cbf7a0c32fca1badf05931879a37c3f8a2b03b4/src/Development/AutomorphicArithmetic/AutomorphicVarieties/ModularityArgument/ThresholdArguments/DenseThresholdCases/CountVOf.lean#L33-L60}{UpperTheory.V.UnitaryThreeInertia.countV\_\allowbreak of\_\allowbreak finiteness\_\allowbreak duality\_\allowbreak euler\_\allowbreak host}.



\paragraph{16. HOST\_\allowbreak NUMBER\_\allowbreak S — C — Uniform original-host rank-three catalogue}

One exponent before all closures, coefficient rings, levels, characters, frames, embeddings and marked points; one separating trace function controls all later finite pairwise nonisomorphic collections.

\emph{Completed inputs (version 0.53).} AFE full Host entry at original scope, LFC full conditional common frame, CAF and MFC, Inherited host-count consumers

\emph{Remaining producer scope.} Complete the original F3 local coefficient bridges and F4 common-overfield readings, with valuation/norm/GoodCoef/iq and the entire coefficient field, and actual complex Hecke standard/no-Jordan plus arithmetic character data on the same form datum. Construct the actual base rows, host/frame lengths, nonzero chains and group ties, MultEq support, eigensystem/carrier/parts, all-embedding logLam, HD/d3b and NEb. Keep the stronger second DivRows3 scope distinct from the completed first DivDef3 and row-pair integrability. The full arithmetic entry and conditional common-frame construction do not by themselves install HOST\_\allowbreak NUMBER\_\allowbreak S, HOST\_\allowbreak FRAME\_\allowbreak LENGTH, TASK\_\allowbreak HOST\_\allowbreak VAN\_\allowbreak E\_\allowbreak INSTALL, TASK\_\allowbreak HOST\_\allowbreak XPP\_\allowbreak INSTALL or TASK\_\allowbreak COUNT\_\allowbreak V\_\allowbreak INSTALL.

\emph{Source.} \href{https://github.com/CocoSgt/riemann-hypothesis-lean/blob/7cbf7a0c32fca1badf05931879a37c3f8a2b03b4/src/Development/AutomorphicArithmetic/AutomorphicVarieties/ModularityArgument/ThresholdArguments/DenseThresholdCases/HostNumberS.lean#L18-L28}{HostNumber\_\allowbreak s}.



\paragraph{17. TASK\_\allowbreak HOST\_\allowbreak VAN\_\allowbreak E\_\allowbreak INSTALL — B — Transport K/duality to literal HostVanE}

Use the already supplied LemmaK/local-global duality and finiteness on the actual CountV host/frame and exact original coefficient modules.

\emph{Source.} \href{https://github.com/CocoSgt/riemann-hypothesis-lean/blob/7cbf7a0c32fca1badf05931879a37c3f8a2b03b4/src/Development/AutomorphicArithmetic/AutomorphicVarieties/UnitaryThreeDeformations/HostNumber.lean#L9-L12}{HostVanE}.

\emph{Source.} \href{https://github.com/CocoSgt/riemann-hypothesis-lean/blob/7cbf7a0c32fca1badf05931879a37c3f8a2b03b4/src/Development/AutomorphicArithmetic/AutomorphicVarieties/LevelCohomology/SoHostVanE.lean#L11-L18}{UpperTheory.V.X.So.hostVanE}.



\paragraph{18. TASK\_\allowbreak HOST\_\allowbreak XPP\_\allowbreak INSTALL — B — Transport CorX to literal HostXpp}

Preserve Psi\textasciicircum{}-2 eta, the prime-to-ell quotient, actual local conditions and the original CountV host while installing BR7.

\emph{Source.} \href{https://github.com/CocoSgt/riemann-hypothesis-lean/blob/7cbf7a0c32fca1badf05931879a37c3f8a2b03b4/src/Development/AutomorphicArithmetic/AutomorphicVarieties/UnitaryThreeDeformations/HostNumber.lean#L14-L17}{HostXpp}.

\emph{Source.} \href{https://github.com/CocoSgt/riemann-hypothesis-lean/blob/7cbf7a0c32fca1badf05931879a37c3f8a2b03b4/src/Development/AutomorphicArithmetic/AutomorphicVarieties/LevelCohomology/Xpp.lean#L11-L84}{UpperTheory.V.X.xpp}.



\paragraph{19. DESIGN\_\allowbreak CHOICES\_\allowbreak COMMON — C — Construct the original common design family and bundle}

For the one target HostExW, construct the original compatible DesignChoices on Cdes=WithTw(WithSm Cprime), where Cprime=WithChN(A intersect target class), A is the one Built construction class, and cT(eps,y)=cTmin(y). Retain every original field, tuple, word, metric and common bundle duty. DSZ10 and RMP11 give size/rank properties of class members, not member existence. The original RecExP/RecAtG and common object/row route remains coupled. Construct ONE original common design family R and its recursive/geometric half with original Inputs/Pins/Holds/ExAt/InGeo and required CoreP. PKT now reconstructs full exists\_\allowbreak pkT, rowPk\_\allowbreak nonempty and exists\_\allowbreak moved with Copy/Coh on those original given inputs and a supplied fitting later stage: one catalogue, both ChT, complete LfT, actual integral lattices and all metric/measure/period constants are retained; one numerical record precedes every eta and V/S/N/G. RDR supplies one Side0 before all readers, ReadPk/pins and full ReadL with exact o.eta=q.eta. The initial actual geometry, Fits/division/central conditions, primitive rational tests and eta agreement must still be produced. Inherited CSR supplies finite columns, one ReadC before later twists, period/evaluation and scalar/error transports, RowSideDetermined and same-R Bundle/BTw/BWit/BunU consumers. RPI supplies exact RowPairInt for actual pairs and finite translates with original mu\_\allowbreak b/fd0; first literal DivDef3 at the same updated R,F,Y,Z,C,sEE follows from its retained original readings, Pins/Holds, stored nonzero IsItW chain, five real-group predicates and group tie. Remaining common fields include host/frame, MultEq nonzero-part support, that original nonzero chain/group tie, eigensystem/carrier/parts and all-embedding logLam, HD/d3b, NEb, actual rational tests and base rows. Retain every simultaneous family hypothesis and eventual threshold. Stronger second DivRows3 remains separate; universal quotient-height CuspRapidDecay and actual common member existence are not inferred from the comparison constructors. Version 0.41 supplies the finite phase-correct rational theta route: exact coefficient extension, split representatives, square-law/PhiIs, Pins from CoversAt, class-sum/point-identity transfer to OneMapOverL, actual-good-projector RatG and both-level LiftTestB. Its original higher covers, all-tuple class-sum uniqueness including phase rescaling, one PointIdentCore before O, base ClassSumsSpan and geometric LatFull remain explicit producers. MLat is available from the original LatFull input. This sharpens the rational-test work without claiming actual common geometry, Fits or eta agreement. Version0.46 proves finite generation of the original whole integral cusp lattice by actual integral-frame evaluation. It also proves the exact identity-cover, analytic unfolding, original-field coefficient descent, global clearing-scalar and both-level lattice-transport implications. Full spanning retains the supplied original IsNab connection and the five explicit principal-chart, level-system, raw-Gaussian, short-rule/dictionary and ramified-jet producers. The higher tuple covers, actual common class sums, PointIdentCore and ClassSumsSpan remain separate. CURRENT VERSION0.47 QUALIFICATION: Construct one original common geometric family and the higher CoversAt/all-tuple class sums and uniqueness with original phase rescalings, one PointIdentCore with common lambda,z, and base ClassSumsSpan. CNC now supplies the full original IsNab; PCL supplies the conditional trace-dual principal-lattice and quantitative normalization. Actual uniformization/HasPrinc and the five principal-chart, finite-level-system, Gaussian/reproduction, short-rule/simultaneous-dictionary and ramified-jet producer groups remain. Keep all original both-level, coefficient-field, Fits and eta agreements. Completed new original producers: ICL1-ICL21, CNC1-CNC16, PCL1-PCL12. Any earlier narrative naming these completed producers as outstanding is historical; this qualification is the current scope authority. CURRENT VERSION0.48 QUALIFICATION: Construct the original finite holomorphic FamilyUnif/N1c, normalized theta families and local lifting, then the common principal charts and finite-level theta systems. PUN now supplies the exact fixed-space models, all-family center lattices, A=1 exponent matching and individual primitive points with the full point-equality criterion; CNC supplies the full canonical connection and PCL the quantitative trace-dual normalization. Complete higher CoversAt and all-tuple class sums/uniqueness with the original phase rescalings, one PointIdentCore with common lambda,z, and base ClassSumsSpan. The Gaussian/reproduction, short-rule/simultaneous-dictionary and ramified-jet producers, and all original both-level, coefficient-field, Fits and eta agreements remain explicit. Completed current producers or conditional consumers: ICL1-ICL21, CNC1-CNC16, PCL1-PCL12, PUN1-PUN24, GCF1-GCF18 at their stated scopes. Earlier frontier descriptions are historical; this qualification is the current scope authority. Original whole-contract premises and quantifiers are retained. CURRENT VERSION0.49 QUALIFICATION: Construct the adjoining simultaneous principal framed points and test-ring dictionary. Full finite holomorphic FamilyUnif/N1c, normalized theta families and local lifts are now supplied by HUF. The next bounded target is allTest\_\allowbreak of\_\allowbreak secBaseChange: represent FrPt by the Lie-eigenframe bundles and theta-section frame line, use the full n=1 SecBaseChangeSt on every nonflat coefficient base change, and prove finite-generated subalgebra descent with the exact TRing/FromTest transport. Then supply the full FramesOK/CoverTri/FibTri and IsUnf/hasPrinc package, including adapted mod-eight symplectic basis, normalized theta multiplier/weight and arithmetic invariance. Continue higher CoversAt, all-tuple class sums and uniqueness with original phase rescalings, one PointIdentCore with common lambda,z, base ClassSumsSpan, Gaussian/reproduction, short-rule dictionaries and ramified jets, retaining both-level, coefficient-field, Fits and eta agreements. Common-family nonemptiness remains distinct from compatibility of an already supplied member. Completed current producers or conditional consumers: ICL1-ICL21, CNC1-CNC16, PCL1-PCL12, PUN1-PUN24, HUF1-HUF21 (full original FamilyUnif/N1c), GCF1-GCF18 and WRS1-WRS10 at their stated scopes. Earlier frontier descriptions are historical; this qualification is the current scope authority. Original whole-contract premises and quantifiers are retained.

\emph{Completed inputs (version 0.53).} FCD1-FCD18 (full homogeneous-chart descent, graded presentation and Noetherian cube model), GSL1-GSL26 (full conditional section laws, universal basis and projective-model algebra), NFL1-NFL22, TDR1-TDR13, NFR1-NFR13, HUF1-HUF21, PUN1-PUN24, ICL1-ICL21, CNC1-CNC16, PCL1-PCL12, GCF/WRS at original supplied-member scope

\emph{Remaining producer scope.} FCD supplies the full original all-chart homogeneous-matrix descent, finite graded presentation, finite graded Cech cohomology and exists\_\allowbreak noetherian\_\allowbreak model\_\allowbreak grCube. One Noetherian R0, one cube C0 and the actual comparisons are chosen before every base change; the nonempty vertices are flat. GSL8 consumes this model. The remaining section-model suppliers are the actual twisted-section extension and power-killing/gluing results, actual dehomogenization and section-cube comparison, affine line-bundle finite projectivity and chart smoothness, on the fixed original scheme and bundles. Do not reopen the algebraic Noetherian-model theorem or retain it as an opaque assumption. The independent PolBaseChangeSt, Serre/multiplication/invariant descent, geometric integrality, principal rank one, torsion section decomposition/growth and theorem-of-square generators remain. Supply those original producers on the same polarized abelian scheme before GSL26 and the completed TDR/NFR/NFL uniformization. Moduli and HUF/PUN foundations remain explicit. Later finite-level systems, Gaussian/reproduction, simultaneous dictionaries, ramified jets, all-tuple class sums, phase-sensitive uniqueness, common PointIdentCore and base ClassSumsSpan remain, retaining both levels, whole fields, Fits and eta agreement.



\paragraph{20. LOCAL\_\allowbreak NBUILT\_\allowbreak COMMON — C — Actual local specialization on the common selected family}

Construct, with the required closed-class property on the same Des(C,3), one eventual family of actual NBuilt certificates for the original Fm.D(x,i), retaining the certificate before forgetting to TSel. HLP now reconstructs full HeckeLift0, all-index commutation, HeckeStd0 and HeckeIntTw from the original MtInput and complex Hecke facts, in the full ambient GenAt range k<=nV. Every maximal ideal is an actual supported residual character class; literal MtSub is the O-span of projections of the original MtLat generators, is finite free, and localizes onto the exact NCov.part. Literal Psi has its exact range and off-space projection error, both original span assertions hold, and ActEx/ActHK are constructed on the same lattice under original semisimplicity. Original MjLat relation quotients, universal coefficient-map relation fields, selected tame generators, coefficient embeddings and all identity transports give the complete nonempty\_\allowbreak spec constructor on its full printed input list. A single compositum combines the four original coefficient requirements with big fixed beforehand. Whole ndSpec\_\allowbreak exists still needs actual common geometric/choice MtInput, complex HeckeStdPlain and the no-Jordan producer, original F3 bridges with valuation/GoodCoef/iq, original F4 common-overfield reading of the entire coefficient field, and exact arithmetic character/reciprocity/admissibility suppliers. Original AutSpec, Galois specialization and the remaining NBuilt fields on that same installation also remain. Do not replace the actual projected lattice or maps by abstract objects with the same dimensions.

\emph{Completed inputs (version 0.53).} FCD full algebraic model, ACP abstract arithmetic suppliers, GMR same-setting measure normalization, GSL/NFL full conditional section laws and uniformization, AFE/LFC full entry and conditional common frame, TDR/NFR and CAF/MFC, WRS selected built-class extraction at membership

\emph{Remaining producer scope.} Continue the distinct retained-defect, period-height and normalized spectral producers, and the original compatible second-layer selection on one common arithmetic/geometric family. GLOBAL\_\allowbreak LOCAL\_\allowbreak GAIN remains fully closed at its original v0.49 scope. Complete local comparison theorems do not themselves construct the whole common family. Common design existence, NBUILT compatibility, original true-column records, shared fields/embeddings, Fits and eta agreement, simultaneous rows and eventual-family thresholds remain explicit. A stronger historical WPt/TSel construction may serve other producers but is not needed to reopen the proved selected-class local-gain estimate. FCD supplies the full original all-chart homogeneous-matrix descent, finite graded presentation, finite graded Cech cohomology and exists\_\allowbreak noetherian\_\allowbreak model\_\allowbreak grCube. One Noetherian R0, one cube C0 and the actual comparisons are chosen before every base change; the nonempty vertices are flat. GSL8 consumes this model. The remaining section-model suppliers are the actual twisted-section extension and power-killing/gluing results, actual dehomogenization and section-cube comparison, affine line-bundle finite projectivity and chart smoothness, on the fixed original scheme and bundles. Do not reopen the algebraic Noetherian-model theorem or retain it as an opaque assumption. The independent PolBaseChangeSt, Serre/multiplication/invariant descent, geometric integrality, principal rank one, torsion section decomposition/growth and theorem-of-square generators remain. Supply those original producers on the same polarized abelian scheme before GSL26 and the completed TDR/NFR/NFL uniformization. Moduli and HUF/PUN foundations remain explicit. Later finite-level systems, Gaussian/reproduction, simultaneous dictionaries, ramified jets, all-tuple class sums, phase-sensitive uniqueness, common PointIdentCore and base ClassSumsSpan remain, retaining both levels, whole fields, Fits and eta agreement.



\paragraph{21. SECOND\_\allowbreak LAYER\_\allowbreak COMPATIBLE\_\allowbreak SELECTION — C — Original second-layer selector and class duties}

For the same A, supply Duty on A and a Closed common original TSel family whose selected records are the original true-column rec7 records (or have the entire dependent tuple transported), with their real witness and original coefficient-field root data. Existence/compatibility remains arithmetic; merely abstractly isomorphic fields or independently selected columns do not suffice. CSR21–22 now construct the same-R Bundle/BTw/BWit/BunU outputs from the exact selected realization/row-side packet and NEb, with the original EvP/ParEv quantifier order. Actual common selection, the five stated realization packets and NEb remain upstream; none is replaced by an independently chosen compatible record. Version 0.38 closes the selected analytic extension and reconstructs the first updated DivDef3 consumer; it leaves the original common selection, geometric readings, rational tests, base rows, host/parts/NEb and exact retained DI premises upstream, with unchanged EvP/ParEv thresholds. Version 0.39 supplies the original given-data reader constructors and exact RowPairInt/first DivDef3 on original premises. It leaves one actual common geometric/comparison/test member, eta agreement, compatible selection, rational tests, base rows, host/parts/NEb and original EvP/ParEv thresholds upstream. Version 0.40 also supplies original comparison/RowPk/fitting-stage Moved algorithms and exact common data transport. The actual initial geometric member, fitting tests, eta agreement and compatible second-layer selector with original eventual thresholds remain upstream. Version 0.41 supplies the finite phase-correct rational theta route: exact coefficient extension, split representatives, square-law/PhiIs, Pins from CoversAt, class-sum/point-identity transfer to OneMapOverL, actual-good-projector RatG and both-level LiftTestB. Its original higher covers, all-tuple class-sum uniqueness including phase rescaling, one PointIdentCore before O, base ClassSumsSpan and geometric LatFull remain explicit producers. MLat is available from the original LatFull input. This sharpens the rational-test work without claiming actual common geometry, Fits or eta agreement. CURRENT VERSION0.49 QUALIFICATION: Continue the distinct global denominator, period/height, normalized spectral and second-layer compatible-selection producers, together with existence of the original common arithmetic/geometric family. GLOBAL\_\allowbreak LOCAL\_\allowbreak GAIN and its actual GroupCompatibilityFamily supplier are now complete on the unchanged selected built class; no additional WPt/TSel realization is needed for that estimate. Do not reopen local first-growth, mean, auxiliary admission, selected D3 tops/Adm2F, the single GIs2 stage or the exact comparison/kernel transport. A full RecT or TSel may remain useful for other tasks, but is not claimed by this local-gain closure. Completed current producers or conditional consumers: GCF1-GCF18 at their stated scopes, WRS1-WRS10 original selected-class compatibility/local gain. Earlier frontier descriptions are historical; this qualification is the current scope authority. Original whole-contract premises and quantifiers are retained.

\emph{Completed inputs (version 0.53).} WLD full generic lattice constructors, GMR actual-setting normalization, FCD/ACP exact nested suppliers, WRS1-WRS10 original selected-class compatibility/local gain, DNE full denominator consumer, DFD/SSC canonical-data and signature comparisons, GSL/NFL and AFE/LFC conditional constructors, Inherited exact same-record realization/transport consumers

\emph{Remaining producer scope.} Continue the distinct retained-defect, period-height and normalized spectral producers, and the original compatible second-layer selection on one common arithmetic/geometric family. GLOBAL\_\allowbreak LOCAL\_\allowbreak GAIN remains fully closed at its original v0.49 scope. Complete local comparison theorems do not themselves construct the whole common family. Common design existence, NBUILT compatibility, original true-column records, shared fields/embeddings, Fits and eta agreement, simultaneous rows and eventual-family thresholds remain explicit. A stronger historical WPt/TSel construction may serve other producers but is not needed to reopen the proved selected-class local-gain estimate.



\paragraph{22. BASE\_\allowbreak THETA\_\allowbreak READINGS — C — Actual base reading on the common bundle}

For every family in the original class, eventually on the SAME selected Bundle and for each i<m, supply the actual R,F,Read,ReadL, weak models/Weil data and phase-one pins. Identify their LiftTie.theta, mtRead, mjRead and ThDef3 basis with that Bundle.P.A3, P.Zs, P.S0 and its stored maps. Exact field, embeddings, metrics, measures, row family and weak choices are retained. This remains one part of the coupled BTH12 producer. PKT now supplies the complete original comparison packet, RowPk and fitting-stage Moved with Copy/Coh on original given geometry, Inputs/Pins/Holds/ExAt/InGeo and Fits. Its ChT/LfT/Reps, actual lattice carriers, original Haar and norm constants, period maps, representatives and all-place RrT are constructed together; its one numerical record precedes every eta/V/S/N/G. RDR gives the original level-zero side, moved-Read and ReadPk/pin constructors on those objects, with one Side0 before every reader and all twelve Read ties. Full ReadL retains the exact o.eta(j,v)=q.eta(j,v) guard for every 1<=j<=m and place v; arbitrary Over does not imply it. Inherited CSR fixes one ReadC before all TwArgs, transports both metrics and all thirteen reading maps, and gives actual period/evaluation with kappa=(c*nP)\textasciicircum{}(-1/2). RPI supplies exact row-pair integrability, so first updated DivDef3 follows from the retained original Pins/Holds, nonzero chain, five real-group predicates and group tie. Initial theta extension remains complete. Remaining work is actual initial common geometry and its original predicates, Fits/primitive rational tests, eta agreement and outer family/eventual/same-Bundle duties, reFace compatibility, MultEq support, original chain/group data, actual base rows and any stronger DivRows3 theta square. No new compatible reading, family threshold or Haar normalization is substituted. Version 0.41 supplies the finite phase-correct rational theta route: exact coefficient extension, split representatives, square-law/PhiIs, Pins from CoversAt, class-sum/point-identity transfer to OneMapOverL, actual-good-projector RatG and both-level LiftTestB. Its original higher covers, all-tuple class-sum uniqueness including phase rescaling, one PointIdentCore before O, base ClassSumsSpan and geometric LatFull remain explicit producers. MLat is available from the original LatFull input. This sharpens the rational-test work without claiming actual common geometry, Fits or eta agreement.

\emph{Completed inputs (version 0.53).} FCD1-FCD18 (full homogeneous-chart descent, graded presentation and Noetherian cube model), GSL1-GSL26, NFL1-NFL22, TDR1-TDR13, NFR1-NFR13, HUF1-HUF21, PUN1-PUN24, CNC1-CNC16, Inherited actual-reading consumers at original inputs

\emph{Remaining producer scope.} FCD supplies the full original all-chart homogeneous-matrix descent, finite graded presentation, finite graded Cech cohomology and exists\_\allowbreak noetherian\_\allowbreak model\_\allowbreak grCube. One Noetherian R0, one cube C0 and the actual comparisons are chosen before every base change; the nonempty vertices are flat. GSL8 consumes this model. The remaining section-model suppliers are the actual twisted-section extension and power-killing/gluing results, actual dehomogenization and section-cube comparison, affine line-bundle finite projectivity and chart smoothness, on the fixed original scheme and bundles. Do not reopen the algebraic Noetherian-model theorem or retain it as an opaque assumption. The independent PolBaseChangeSt, Serre/multiplication/invariant descent, geometric integrality, principal rank one, torsion section decomposition/growth and theorem-of-square generators remain. Supply those original producers on the same polarized abelian scheme before GSL26 and the completed TDR/NFR/NFL uniformization. Moduli and HUF/PUN foundations remain explicit. Later finite-level systems, Gaussian/reproduction, simultaneous dictionaries, ramified jets, all-tuple class sums, phase-sensitive uniqueness, common PointIdentCore and base ClassSumsSpan remain, retaining both levels, whole fields, Fits and eta agreement. Complete the original F3 local coefficient bridges and F4 common-overfield readings, with valuation/norm/GoodCoef/iq and the entire coefficient field, and actual complex Hecke standard/no-Jordan plus arithmetic character data on the same form datum. Construct the actual base rows, host/frame lengths, nonzero chains and group ties, MultEq support, eigensystem/carrier/parts, all-embedding logLam, HD/d3b and NEb. Keep the stronger second DivRows3 scope distinct from the completed first DivDef3 and row-pair integrability. The full arithmetic entry and conditional common-frame construction do not by themselves install HOST\_\allowbreak NUMBER\_\allowbreak S, HOST\_\allowbreak FRAME\_\allowbreak LENGTH, TASK\_\allowbreak HOST\_\allowbreak VAN\_\allowbreak E\_\allowbreak INSTALL, TASK\_\allowbreak HOST\_\allowbreak XPP\_\allowbreak INSTALL or TASK\_\allowbreak COUNT\_\allowbreak V\_\allowbreak INSTALL.



\paragraph{23. BASE\_\allowbreak THETA\_\allowbreak RATIONAL\_\allowbreak TESTS — C — Simultaneous original base rational tests}

On the SAME original Bundle/reading/theta map, for every i<m supply one finite separating complex-linear test family whose values lie in the fixed embedded coefficient field on theta of EVERY original rational source vector and on EVERY stored target-coordinate basis vector. BTH8/BTH11 use this same LiftTestB for rational adjoint image and matrix coordinates. RTG now reconstructs prescribed-field coefficient extension, original split class representatives/RepsOK, the self-dual split-lattice square law and finite residue PhiIs; every Pins field follows from the actual CoversAt. With the base coefficient record literally fixed, exact class-sum existence/uniqueness including phase-rescaled O, PointIdentCore with common lambda,z before every O, and base ClassSumsSpan imply all-tuple finite lift span and OneMapOverL. The phase is exactly cj(twAt embLL A3.A.v0)(R.phs i). RatG uses the actual good projector eOp(y), and full LiftTestB follows at BOTH original levels. These finite implication proofs are complete. Remaining original producers are the prescribed higher covers, those actual class sums/uniqueness, PointIdentCore, base ClassSumsSpan and the base/common geometric datum; exact MLat is supplied only from original LatFull, whose geometric production remains separate. Finite dimension or an unrelated rational theta model cannot replace these same-object identities. Original outer eventual-family/Bundle quantifiers remain. Version0.46 proves finite generation of the original whole integral cusp lattice by actual integral-frame evaluation. It also proves the exact identity-cover, analytic unfolding, original-field coefficient descent, global clearing-scalar and both-level lattice-transport implications. Full spanning retains the supplied original IsNab connection and the five explicit principal-chart, level-system, raw-Gaussian, short-rule/dictionary and ramified-jet producers. The higher tuple covers, actual common class sums, PointIdentCore and ClassSumsSpan remain separate. CURRENT VERSION0.47 QUALIFICATION: Construct one original common geometric family and the higher CoversAt/all-tuple class sums and uniqueness with original phase rescalings, one PointIdentCore with common lambda,z, and base ClassSumsSpan. CNC now supplies the full original IsNab; PCL supplies the conditional trace-dual principal-lattice and quantitative normalization. Actual uniformization/HasPrinc and the five principal-chart, finite-level-system, Gaussian/reproduction, short-rule/simultaneous-dictionary and ramified-jet producer groups remain. Keep all original both-level, coefficient-field, Fits and eta agreements. Completed new original producers: ICL1-ICL21, CNC1-CNC16, PCL1-PCL12. Any earlier narrative naming these completed producers as outstanding is historical; this qualification is the current scope authority. CURRENT VERSION0.48 QUALIFICATION: Construct the original finite holomorphic FamilyUnif/N1c, normalized theta families and local lifting, then the common principal charts and finite-level theta systems. PUN now supplies the exact fixed-space models, all-family center lattices, A=1 exponent matching and individual primitive points with the full point-equality criterion; CNC supplies the full canonical connection and PCL the quantitative trace-dual normalization. Complete higher CoversAt and all-tuple class sums/uniqueness with the original phase rescalings, one PointIdentCore with common lambda,z, and base ClassSumsSpan. The Gaussian/reproduction, short-rule/simultaneous-dictionary and ramified-jet producers, and all original both-level, coefficient-field, Fits and eta agreements remain explicit. Completed current producers or conditional consumers: ICL1-ICL21, CNC1-CNC16, PCL1-PCL12, PUN1-PUN24. Earlier frontier descriptions are historical; this qualification is the current scope authority. Original whole-contract premises and quantifiers are retained. CURRENT VERSION0.49 QUALIFICATION: Construct the adjoining simultaneous principal framed points and test-ring dictionary. Full finite holomorphic FamilyUnif/N1c, normalized theta families and local lifts are now supplied by HUF. The next bounded target is allTest\_\allowbreak of\_\allowbreak secBaseChange: represent FrPt by the Lie-eigenframe bundles and theta-section frame line, use the full n=1 SecBaseChangeSt on every nonflat coefficient base change, and prove finite-generated subalgebra descent with the exact TRing/FromTest transport. Then supply the full FramesOK/CoverTri/FibTri and IsUnf/hasPrinc package, including adapted mod-eight symplectic basis, normalized theta multiplier/weight and arithmetic invariance. Continue higher CoversAt, all-tuple class sums and uniqueness with original phase rescalings, one PointIdentCore with common lambda,z, base ClassSumsSpan, Gaussian/reproduction, short-rule dictionaries and ramified jets, retaining both-level, coefficient-field, Fits and eta agreements. Completed current producers or conditional consumers: ICL1-ICL21, CNC1-CNC16, PCL1-PCL12, PUN1-PUN24, HUF1-HUF21 (full original FamilyUnif/N1c and theta family). Earlier frontier descriptions are historical; this qualification is the current scope authority. Original whole-contract premises and quantifiers are retained.

\emph{Completed inputs (version 0.53).} FCD1-FCD18 (full homogeneous-chart descent, graded presentation and Noetherian cube model), GSL1-GSL26, NFL1-NFL22, TDR1-TDR13, NFR1-NFR13, HUF1-HUF21, PUN1-PUN24, ICL1-ICL21, CNC1-CNC16, PCL1-PCL12, Inherited both-level RatG/LiftTestB consumers

\emph{Remaining producer scope.} FCD supplies the full original all-chart homogeneous-matrix descent, finite graded presentation, finite graded Cech cohomology and exists\_\allowbreak noetherian\_\allowbreak model\_\allowbreak grCube. One Noetherian R0, one cube C0 and the actual comparisons are chosen before every base change; the nonempty vertices are flat. GSL8 consumes this model. The remaining section-model suppliers are the actual twisted-section extension and power-killing/gluing results, actual dehomogenization and section-cube comparison, affine line-bundle finite projectivity and chart smoothness, on the fixed original scheme and bundles. Do not reopen the algebraic Noetherian-model theorem or retain it as an opaque assumption. The independent PolBaseChangeSt, Serre/multiplication/invariant descent, geometric integrality, principal rank one, torsion section decomposition/growth and theorem-of-square generators remain. Supply those original producers on the same polarized abelian scheme before GSL26 and the completed TDR/NFR/NFL uniformization. Moduli and HUF/PUN foundations remain explicit. Later finite-level systems, Gaussian/reproduction, simultaneous dictionaries, ramified jets, all-tuple class sums, phase-sensitive uniqueness, common PointIdentCore and base ClassSumsSpan remain, retaining both levels, whole fields, Fits and eta agreement.



\paragraph{24. SLR\_\allowbreak BASE\_\allowbreak ROWS\_\allowbreak PRODUCER — C — Actual base row identities on that same bundle}

For the SAME selected Bundle, common original fields and readings used by the other three BTH12 inputs, supply actual RowsHold0\_\allowbreak 3 uniformly eventually for every family. Its original primitive row realization, spectrum/type and basis identifications must be retained. BTH supplies a theta consumer only; it does not construct the base row identity. No separate witness may be chosen merely to satisfy this row requirement.



\clearpage
\section{One coefficient ring for all flat charts}
\label{sec:flat-chart-descent-v053}

\status{FCD1--FCD18 prove the original homogeneous-matrix flat-chart
descent theorem, the finite graded presentation theorem, and the
Noetherian graded-cube model theorem.  The original coefficient ring is
arbitrary.  One coefficient subring, one presentation and one comparison
map serve every chart.  The subsequent construction of this graded
module from actual geometric sections remains separate; no global
\texttt{Perf} or section law is asserted here.  The accepted formal
source is fixed at
\texttt{7cbf7a0c32fca1badf05931879a37c3f8a2b03b4}.}

The section-model argument needs a finite complex that continues to
describe the same sections after every change of coefficients.  The
linear algebra producing that complex was proved in GSL5--GSL10.  Its
remaining input is a suitable model over a Noetherian coefficient ring.
The present section constructs the algebraic part of that model.
The important choice is quite concrete: retain the coefficients of the
presentation \emph{and of certificates for all the chart properties}.
Keeping only the presentation coefficients would give no reason for
flatness to survive over the smaller ring.

All rings are commutative with identity, including the zero ring.
Modules need not be faithful, and no coefficient extension below is
assumed flat.  If $\iota$ is a finite ordered set, write
\[
 S_R=R[X_i\mid i\in\iota],\qquad
 u_i=X_i-1,\qquad X_J=\prod_{i\in J}X_i.
\]
For a matrix $\Phi\in\operatorname{Mat}_{a\times b}(S_R)$,
$\operatorname{coker}\Phi$ means
$S_R^a/\Phi S_R^b$.  Thus the source's
\texttt{Quot\allowbreak SMul\allowbreak Top} $u_i(\operatorname{coker}\Phi)$ is exactly
\[
 Q_i(\Phi)=S_R^a/(\Phi S_R^b+u_iS_R^a).                 \tag{FCD.1}\label{FCD.1}
\]
A $\mathbb Z$-graded module has a direct-sum decomposition; in
particular, every one of its elements has only finitely many nonzero
homogeneous components.  The variables have degree one.  A free
module with basis degrees $d_1,\ldots,d_a$ has degree-$n$ part
$\bigoplus_l(S_R)_{n-d_l}e_l$.

We use the usual elementary properties of tensor products, localization
and projective modules, and Hilbert's basis theorem.  The needed
Artin--Rees, completion and smoothness arguments are given in FCD5.
These conventional algebraic inputs are distinct from
the geometric producers retained after GSL26.  The source map at the
end identifies both the original declarations and the pinned Mathlib
completion results used for this background.

\subsection{Finite certificates attached to one presentation}

% BEGIN FCD1
\paragraph{FCD1: finite relations for a prescribed set of generators.}
\label{claim:FCD1}
Let $B$ be a finitely presented $R$-algebra, and let
$\pi:R[X_1,\ldots,X_m]\twoheadrightarrow B$ be a surjective
$R$-algebra map.  Its kernel is a finitely generated ideal, even when
$R$ is not Noetherian.  Similarly, if $Q$ is a finitely presented
$S$-module and $F\twoheadrightarrow Q$ is a surjection from a
finitely generated $S$-module, its kernel is finitely generated.

\begin{proof}
Choose a finite presentation
$B=R[Y_1,\ldots,Y_t]/(h_1,\ldots,h_r)$.
Lift each image of $Y_l$ to a polynomial $f_l(X)$, and express each
$\pi(X_j)$ by a polynomial $v_j(Y)$ in this presentation.  The
polynomials
\[
 h_k(f(X)),\qquad X_j-v_j(f(X))                         \tag{FCD.2}\label{FCD.2}
\]
belong to $\ker\pi$.  They generate it.  Indeed, if $p(X)$ maps to
zero, then $p(v(Y))$ belongs to $(h_1,\ldots,h_r)$.
After substituting $Y=f(X)$, the polynomial $p(v(f(X)))$ is in
the ideal generated by the first family in \eqref{FCD.2}.
For any polynomial $p$ the difference $p(X)-p(Z)$ lies in
$(X_1-Z_1,\ldots,X_m-Z_m)$: check a monomial by replacing one
factor at a time, and then add.  With $Z=v(f(X))$ this proves the
assertion about $\ker\pi$.

For the module assertion, first let $E=S^a\twoheadrightarrow Q$
be finite free.  Take a known finite presentation
$S^b\to S^c\xrightarrow{q}Q\to0$, and lift the generators in
both directions to maps $\alpha:S^c\to E$ and
$\beta:E\to S^c$.  If $f:E\to Q$ is the given surjection,
these maps satisfy $f\alpha=q$ and $q\beta=f$.
The kernel of $f$ is generated by the columns of $1-\alpha\beta$
and the images under $\alpha$ of the finitely many given relations.
For $z\in\ker f$, write
$z=(1-\alpha\beta)z+\alpha\beta z$; the second term has the
required form because $q\beta z=0$.  Every displayed generator is
in the kernel.  Finally choose a finite free surjection $E\to F$.
The finitely generated kernel of the composite $E\to Q$ maps
surjectively onto $\ker(F\to Q)$.  This proves the general case.
\end{proof}
% END FCD1

% BEGIN FCD2
\paragraph{FCD2: a generalized inverse from the actual projective quotient.}
\label{claim:FCD2}
Let $S\twoheadrightarrow B$ be a surjective algebra map $\pi$,
with kernel $I=(g_1,\ldots,g_r)$.  Suppose $u\in I$ and
$\epsilon:S^a\twoheadrightarrow P$ is $\pi$-semilinear, where
$P$ is projective over $B$.  Suppose also that
\[
 \ker\epsilon=\Phi S^b+uS^a.                          \tag{FCD.3}\label{FCD.3}
\]
There are matrices $\Psi\in\operatorname{Mat}_{b\times a}(S)$
and $E_k\in\operatorname{Mat}_{a\times b}(S)$ such that
\[
 \Phi\Psi\Phi=\Phi+\sum_{k=1}^r g_kE_k.              \tag{FCD.4}\label{FCD.4}
\]

\begin{proof}
Put $p_l=\epsilon(e_l)$.  Semilinearity gives
$\epsilon(y)=\sum_l\pi(y_l)p_l$.  Consequently the map
$F:B^a\to P$, $F(\beta)=\sum_l\beta_lp_l$, is surjective.
Projectivity supplies a $B$-linear section $s:P\to B^a$, so
$Fs=1_P$.  For each $l$, lift $e_l-s(p_l)\in B^a$ to
$y_l\in S^a$.  Its image under $\epsilon$ is zero.  By
\eqref{FCD.3}, write $y_l=\Phi w_l+u w'_l$.  Since $\pi(u)=0$,
the matrix $\Psi$ with columns $w_l$ satisfies
\[
 \overline\Phi\,\overline\Psi=1_{B^a}-sF.           \tag{FCD.5}\label{FCD.5}
\]
Here bars mean application of $\pi$ to the entries.  The columns
of $\Phi$ lie in $\ker\epsilon$, hence $F\overline\Phi=0$.
Multiplying \eqref{FCD.5} on the right by $\overline\Phi$ gives
$\overline\Phi\,\overline\Psi\,\overline\Phi
=\overline\Phi$.  Every entry of
$\Phi\Psi\Phi-\Phi$ therefore lies in $I$.
Express its finitely many entries as sums $\sum_k g_k e_{k,l,j}$
and assemble $E_k=(e_{k,l,j})$.  This is \eqref{FCD.4}.

The section $s$ and its lifted columns were chosen for the given
$P,\epsilon,\Phi$.  No separately chosen quotient module or
identification is introduced.
\end{proof}
% END FCD2

% BEGIN FCD3
\paragraph{FCD3: all certificates on a single chart.}
\label{claim:FCD3}
Fix $i\in\iota$, and suppose the chart data consist of a smooth
$R$-algebra $B$, a surjection $\pi:S_R\to B$ with
$\pi(X_i)=1$, a projective $B$-module $P$, and a surjective
$\pi$-semilinear map $\epsilon:S_R^a\to P$ satisfying
\eqref{FCD.3} with $u=u_i$.
There is a finite family $g_1,\ldots,g_r$ generating $\ker\pi$,
and finite polynomial data $c_k,C_k,C'_k,\Psi,E_k$ such that
\begin{align}
 u_i&=\sum_k c_kg_k,\nonumber\\
 g_k1_a&=\Phi C_k+u_iC'_k,\nonumber\\
 \Phi\Psi\Phi&=\Phi+\sum_k g_kE_k.                   \tag{FCD.6}\label{FCD.6}
\end{align}
Here $C_k$ has size $b\times a$, $C'_k$ has size $a\times a$,
and the last two sizes are as in FCD2.  The quotient $S_R/(g_k)$
is formally smooth over $R$.

\begin{proof}
Smoothness includes finite presentation and formal smoothness.
FCD1 gives the finite generating family for $\ker\pi$.
Since $\pi(u_i)=0$, express $u_i$ in that ideal to obtain $c_k$.
For every $k,l$, semilinearity implies
$\epsilon(g_ke_l)=\pi(g_k)p_l=0$.
Use \eqref{FCD.3} to write
$g_ke_l=\Phi w_{k,l}+u_iw'_{k,l}$; the vectors $w_{k,l}$ and
$w'_{k,l}$ are the columns of $C_k,C'_k$.  This proves the
middle identity of \eqref{FCD.6}.  FCD2 proves the last identity.
Finally the actual map $\pi$ identifies $S_R/(g_k)$ with $B$
as an $R$-algebra, and transports its formal smoothness.  Only
finitely many matrices and entries have been chosen.
\end{proof}
% END FCD3

% BEGIN FCD4
\paragraph{FCD4: formal smoothness expressed by finite polynomial identities.}
\label{claim:FCD4}
Let $S=R[X_i\mid i\in\iota]$ and $I=(g_1,\ldots,g_r)$.
The quotient $A=S/I$ is formally smooth over $R$ if and only if
there are polynomials $a_{ik}$ and $b_{jkl}$ such that, with
\[
 y_i=X_i+\sum_k a_{ik}g_k,
\]
one has
\[
 g_j(y)=\sum_{k,l}b_{jkl}g_kg_l\quad(1\le j\le r).   \tag{FCD.7}\label{FCD.7}
\]
This criterion is valid over an arbitrary coefficient ring.

\begin{proof}
If $A$ is formally smooth, lift $1_A$ across the square-zero
surjection $S/I^2\to S/I$ to an $R$-algebra section
$\sigma:A\to S/I^2$.  Lift each $\sigma(\overline X_i)$ to a
polynomial $y_i\in S$.  Because the reduction of $\sigma$ is the
identity, $y_i-X_i\in I$, giving $a_{ik}$.  Evaluation at $y$
followed by reduction modulo $I^2$ equals the composite
$S\to A\xrightarrow{\sigma}S/I^2$: these maps agree on $R$
and on every variable.  Thus $g_j(y)\in I^2$.
The products $g_kg_l$ generate $I^2$, giving the coefficients
$b_{jkl}$ in \eqref{FCD.7}.

Conversely \eqref{FCD.7} defines an $R$-algebra map
$\sigma:A\to S/I^2$ by $\overline X_i\mapsto\overline y_i$.
Its reduction is the identity.  To verify formal smoothness directly,
let $D\to D/J$ be an $R$-algebra quotient with $J^2=0$, and
let $h:A\to D/J$.  Choose lifts in $D$ of the finitely many
$h(\overline X_i)$.  Evaluation gives $t:S\to D$, with
$t(I)\subseteq J$, hence $t(I^2)=0$.  The induced map
$\overline t:S/I^2\to D$ composed with $\sigma$ lifts $h$.
For a nilpotent ideal $J$, lift successively through
$D/J^{n+1}\to D/J^n$: its kernel $J^n/J^{n+1}$ is square-zero
for $n\ge1$.  This proves the full infinitesimal lifting property.

The coefficients of \eqref{FCD.7} are the finite evidence of formal
smoothness that can be retained in a smaller ring.  A bare assertion
that the original chart was smooth would not by itself provide such
evidence after descent.
\end{proof}
% END FCD4

% BEGIN FCD5
\paragraph{FCD5: the flatness consequence over a Noetherian coefficient ring.}
\label{claim:FCD5}
If $R_0$ is Noetherian, $S_0=R_0[X_i\mid i\in\iota]$, and
$A_0=S_0/I_0$ is formally smooth over $R_0$, then $A_0$ is
flat over $R_0$.

\begin{proof}
Hilbert's basis theorem makes $S_0$ Noetherian.  Let
$\widehat S_0=\varprojlim_{n\ge1}S_0/I_0^n$.
Starting with $1_{A_0}:A_0\to S_0/I_0$, formal smoothness gives
compatible lifts $\sigma_n:A_0\to S_0/I_0^n$.
At the next step the kernel $I_0^n/I_0^{n+1}$ has square zero,
so the lift can be chosen to reduce to the \emph{already chosen}
$\sigma_n$.  The resulting $\sigma:A_0\to\widehat S_0$
is a section of the first projection $p:\widehat S_0\to A_0$.

For completeness, the completion fact used here is that the
$I$-adic completion of a Noetherian ring $S$ is flat over $S$.
Here is its usual module argument.  For an inclusion $N\subseteq M$
of finite $S$-modules, Artin--Rees says that the filtrations
$I^nN$ and $N\cap I^nM$ define the same completion: for a fixed
$c$ and $n\ge c$,
\[
 I^nN\subseteq N\cap I^nM\subseteq I^{n-c}N.
\]
The required inclusion has a short proof.  The Rees ring
$\mathcal R=\bigoplus_{n\ge0}I^nt^n=S[It]$ is a finitely
generated $S$-algebra, hence Noetherian.  The module
$\bigoplus_{n\ge0}I^nMt^n$ is finite over $\mathcal R$,
generated by any finite generating family of $M$ in degree zero.
Its graded submodule
$\bigoplus_{n\ge0}(N\cap I^nM)t^n$ is therefore finite.
Choose finitely many homogeneous generators, all in degrees at
most $c$.  In degree $n\ge c$ they express every element of
$N\cap I^nM$ as a sum of elements of
$I^{n-d}(N\cap I^dM)$ with $d\le c$, each contained in
$I^{n-c}N$.  The other inclusion is immediate.  This proves
the Artin--Rees assertion used here.

The finite-level sequences with left term
$N/(N\cap I^nM)$, middle term $M/I^nM$ and right term
$(M/N)/I^n(M/N)$ are exact.  The transition maps on the left
are surjective.  Consequently taking inverse limits preserves this
short exact sequence.  Explicitly, a compatible right-hand sequence
can be lifted recursively: after any lift at the next level, its
discrepancy from the previously chosen lift is in the left term and
can be corrected using that surjectivity.  Artin--Rees then identifies
the left limit with $\widehat N$.  Thus completion is exact on
finite modules.

A finite presentation $S^b\to S^a\to M\to0$ now identifies
$\widehat M$ with $\widehat S\otimes_SM$: both are the
cokernel of the same matrix $\widehat S^b\to\widehat S^a$.
In particular, for every ideal $J\subseteq S$, which is finite
because $S$ is Noetherian, the map
$\widehat S\otimes_SJ\to\widehat S$ is injective.
The ideal criterion for flatness gives flatness of $\widehat S$.
This argument uses neither a separatedness assumption on $S$
nor an assumption that $I$ lies in its Jacobson radical.

Apply this to $S_0,I_0$.  Since $S_0$ is free over $R_0$,
$\widehat S_0$ is flat over $R_0$ by composition of tensor
functors.  The maps $\sigma,p$ exhibit $A_0$ as an $R_0$-module
retract of $\widehat S_0$.  A retract of a flat module is flat:
after tensoring an injection, its induced map is a direct summand
of an injection.  Hence $A_0$ is $R_0$-flat.
\end{proof}
% END FCD5

% BEGIN FCD6
\paragraph{FCD6: the explicit retract for a dehomogenized cokernel.}
\label{claim:FCD6}
Let $S$ be an $R$-algebra, $I=(g_1,\ldots,g_r)$, and suppose
$S/I$ is flat over $R$.  Suppose $u,\Phi$ and matrices satisfy
the three identities \eqref{FCD.6}, with $u$ in place of $u_i$.
Then
$Q=S^a/(\Phi S^b+uS^a)$ is $R$-flat.  More precisely, $Q$ is
an explicit $R$-module retract of $(S/I)^a$.

\begin{proof}
Write $\operatorname{pr}:S^a\to Q$ for the quotient map and
$q:S^a\to(S/I)^a$ for reduction of the coordinates.
Applying $g_k1_a=\Phi C_k+uC'_k$ to a vector shows that
$g_k$ annihilates $Q$.  Thus $I$ annihilates $Q$, and there is
a well-defined map
\[
 p:(S/I)^a\longrightarrow Q,\qquad p(q(v))=\operatorname{pr}(v).
\]
Define the other map by the concrete formula
\[
 j:Q\longrightarrow(S/I)^a,\qquad
 j(\operatorname{pr}(v))=q(v-\Phi\Psi v).             \tag{FCD.8}\label{FCD.8}
\]
To check that this is well-defined, first apply $1-\Phi\Psi$
to $\Phi w$.  By the last identity of \eqref{FCD.6} the result
is $-\sum_k g_kE_kw$, whose reduction is zero.  Secondly,
$(1-\Phi\Psi)uv$ reduces to zero because
$u=\sum_kc_kg_k\in I$.  These are exactly the two families
of relations in $Q$.

The composite satisfies
\[
 p j(\operatorname{pr}(v))
 =\operatorname{pr}(v-\Phi\Psi v)=\operatorname{pr}(v).
\]
Hence $pj=1_Q$.  A finite direct sum of the $R$-flat module
$S/I$ is flat, and FCD5's retract argument applies.  All maps
were given before any further coefficient change.  Tensoring the
identity $pj=1_Q$ therefore preserves this retract over every
new coefficient ring as well.
\end{proof}
% END FCD6

% BEGIN FCD7
\paragraph{FCD7: descend all finitely many identities to one coefficient ring.}
\label{claim:FCD7}
Let a finite family of polynomials over an arbitrary ring $R$ be given.
There is a Noetherian subring $R_0\subseteq R$ containing their
coefficients.  Their unique coefficientwise lifts to $S_{R_0}$
satisfy every polynomial identity between the original family,
including identities involving substitution and matrix multiplication.
If a specified polynomial is homogeneous of an integer degree $d$,
its lift has that same degree.

\begin{proof}
Let $T\subseteq R$ be the union of the coefficient sets of the
finitely many polynomials, and let $R_0$ be the subring generated
by $T$.  This is the image of a polynomial ring
$\mathbb Z[t_1,\ldots,t_N]$, so it is Noetherian by Hilbert's
basis theorem and passage to a quotient.  The argument also covers
positive characteristic and the zero ring.

Each polynomial lifts by viewing its coefficients as elements of
$R_0$.  The map $S_{R_0}\to S_R$ is injective, since equality
of polynomials means equality of every coefficient.  It commutes
with sums, products and substitution in the variables.  Thus the
difference of the two lifted sides of an identity maps to zero
and is already zero in $S_{R_0}$.  Matrix identities reduce
entrywise to these operations.

For the homogeneity assertion, the coefficient of every monomial
of degree different from $d$ is zero after inclusion in $R$.
Injectivity makes it zero in $R_0$.  This works also for negative
$d$, when the only homogeneous polynomial is zero.  No division
or cancellation in $R$ has occurred.
\end{proof}
% END FCD7

% BEGIN FCD8
\paragraph{FCD8: the original simultaneous flat-chart descent theorem.}
\label{claim:FCD8}
Let $R$ be arbitrary, $\iota$ finite and ordered, and let
$d:\{1,\ldots,a\}\to\mathbb Z$ and
$e:\{1,\ldots,b\}\to\mathbb Z$.  Suppose
$\Phi\in\operatorname{Mat}_{a\times b}(S_R)$ has
$\Phi_{lj}$ homogeneous of degree $e_j-d_l$.
For \emph{every} $i\in\iota$, suppose there exist a smooth
$R$-algebra $B_i$, a surjective $R$-algebra map $\pi_i:S_R\to B_i$
with $\pi_i(X_i)=1$, a projective $B_i$-module $P_i$, and a
surjective $\pi_i$-semilinear map $\epsilon_i:S_R^a\to P_i$
such that
\[
 \epsilon_i(y)=0
 \quad\Longleftrightarrow\quad
 y=\Phi w+(X_i-1)w'\quad\hbox{for some }w,w'.         \tag{FCD.9}\label{FCD.9}
\]
Then there exist one Noetherian subring $R_0\subseteq R$ and
one matrix $\Phi_0$ over $S_{R_0}$ such that its entrywise
image is $\Phi$, its entries have the same degrees $e_j-d_l$,
and
\[
 Q_i(\Phi_0)\ \text{is flat over }R_0
 \quad\text{for every }i\in\iota.                   \tag{FCD.10}\label{FCD.10}
\]
This is the full conclusion of the source theorem
\begin{center}
\texttt{exists\_\allowbreak{}noetherian\_\allowbreak{}subring\_\allowbreak{}}\allowbreak\texttt{homogeneous\_\allowbreak matrix\_\allowbreak flat\_\allowbreak charts}.
\end{center}

\begin{proof}
Choose one set of the supplied data for each of the finitely many
$i$.  Apply FCD3, and then FCD4, to obtain its families
$g_{i,k},c_{i,k},C_{i,k},C'_{i,k},\Psi_i,E_{i,k}$ and the
formal-smoothness polynomials $a_{i,v,k},b_{i,j,k,l}$.
Collect the coefficients of \emph{all} these polynomials and
of all entries of $\Phi$.  The index set, every relation family
and every matrix are finite.  FCD7 therefore gives one $R_0$
and coefficientwise lifts of every family, including $\Phi_0$.
It also preserves the homogeneous degrees of $\Phi$.

Fix any chart $i$.  Its lifted equations \eqref{FCD.7} prove
formal smoothness of
\[
 B_{i,0}=S_{R_0}/(g_{i,1,0},\ldots,g_{i,r_i,0})
\]
by FCD4.  The ring $R_0$ is Noetherian, so FCD5 proves that
$B_{i,0}$ is $R_0$-flat.  All three equations
\eqref{FCD.6} hold over $S_{R_0}$ by FCD7.  Apply FCD6 with
$u=X_i-1$ and this same $\Phi_0$ to obtain \eqref{FCD.10}.
This works for every $i$ without enlarging or changing $R_0$.

Notice the order: the finite chart witnesses are chosen first;
the common coefficient ring is then formed from their union;
all flatness assertions are proved on that one ring.  The input
$R$ itself has never been assumed Noetherian or reduced.
\end{proof}
% END FCD8

\subsection{A finite graded presentation that is correct on every chart}

% BEGIN FCD9
\paragraph{FCD9: three elementary dehomogenization identities.}
\label{claim:FCD9}
Let $M=\bigoplus_{n\in\mathbb Z}M_n$ be graded, and put
$x=X_i$ for a fixed variable.  Then:
\begin{enumerate}
\item multiplication by $x-1$ on $M$ is injective;
\item if a homogeneous $z\in M_n$ satisfies $z=(x-1)w$,
then $x^kz=0$ for some $k\ge0$;
\item for each $y\in M$, and every sufficiently large integer $N$,
there are $y_N\in M_N$ and $v\in M$ with $y=y_N+(x-1)v$.
\end{enumerate}
These assertions do not require multiplication by $x$ to be injective.

\begin{proof}
Write $w=\sum_jw_j$.  The degree-$j$ component of $(x-1)w$
is $xw_{j-1}-w_j$.  If this is zero for every $j$, start below
the finite support of $w$ and use $w_j=xw_{j-1}$ successively;
every component is zero.  This proves the first assertion.

For the second assertion the same recurrence holds except in degree
$n$.  It gives $w_{n-1}=0$, then $z=-w_n$, and
$w_{n+k}=x^kw_n$ for $k\ge0$.  Above the finite support of
$w$ the left side is zero, so some $x^kz$ is zero.

For the third assertion choose $N$ at least every degree in the
support of $y=\sum_jy_j$, and set
\[
 y_N=\sum_j x^{N-j}y_j\in M_N.
\]
The geometric-sum identity
$x^k-1=(x-1)(1+x+\cdots+x^{k-1})$ shows term by term that
$y-y_N$ is a multiple of $x-1$.  The formula also covers $k=0$
by taking an empty sum.  This supplies the required $v$ explicitly.
\end{proof}
% END FCD9

% BEGIN FCD10
\paragraph{FCD10: finite generators up to powers of each chart variable.}
\label{claim:FCD10}
If $M/(X_i-1)M$ is finitely generated over $S_R$, there is a
finite homogeneous family in $M$ such that every $y\in M$ has
some $X_i^ky$ in its $S_R$-span.  If this hypothesis holds for
every $i\in\iota$, one finite homogeneous family works for
all $i$, with the exponent allowed to depend on $i$ and $y$.
More generally, if each chart already has a finite family with
this power-span property, their homogeneous components give one
homogeneous family with the same property on all charts.

\begin{proof}
For one chart, take finitely many generators of the quotient.
By FCD9, lift each to a homogeneous element of $M$.  Let $N$
be their span.  It is a graded submodule: multiplying a homogeneous
generator by the homogeneous components of a polynomial decomposes
every element of $N$ into components that still belong to $N$.
Every $y\in M$ has $y-s\in(x-1)M$ for some $s\in N$.

First suppose $y$ has degree $n$.  Choose a common degree $D$
at least $n$ and all degrees occurring in $s$.  Homogenizing $s$
as in FCD9 gives $s_D\in N\cap M_D$.  The two elements
$x^{D-n}y$ and $s_D$ have the same image in $M/(x-1)M$.
Their difference is homogeneous, so FCD9 gives
$x^k(x^{D-n}y-s_D)=0$ for some $k$.  Hence
$x^{k+D-n}y=x^ks_D\in N$.  For a general element, apply
this argument to each of its finitely many homogeneous components
and increase to a common exponent.  Their sum then has the same
power-span property.

For finitely many charts take the union of the finite homogeneous
families just constructed.  In the more general formulation, first
take the union of the given finite families, and then replace each
element by all of its finitely many homogeneous components.  The
new span contains the old span on every chart.  The empty chart
index set is harmless: the empty family satisfies the vacuous
chart conditions.
\end{proof}
% END FCD10

% BEGIN FCD11
\paragraph{FCD11: the relations also have finite dehomogenizations.}
\label{claim:FCD11}
Let $\varphi:F\to G$ be a graded $S_R$-linear map with $F$
finitely generated.  Fix $i$, put $x=X_i$ and $u=x-1$, and
assume that every $z\in G$ has some $x^kz\in\operatorname{im}\varphi$.
If $G/uG$ is finitely presented, then
$(\ker\varphi)/u(\ker\varphi)$ is finitely generated.

\begin{proof}
Write $K=\ker\varphi$.  The composite
$F\to G\to G/uG$ is surjective: $x^kz$ and $z$ have the
same class modulo $uG$.  By FCD1 its kernel $J$ is finitely
generated.  We claim
\[
 J=K+uF,\qquad K\cap uF=uK.                          \tag{FCD.11}\label{FCD.11}
\]
The easy inclusion in the first equality follows by applying
$\varphi$.  For its converse, suppose $\varphi(y)=uw$.
By FCD9, write $y=y'+uv_0$ with $y'$ homogeneous.
Then $\varphi(y')=u(w-\varphi(v_0))$ is homogeneous.
FCD9 gives $x^k\varphi(y')=0$.  Therefore $k'=x^ky'$
belongs to $K$.  Write $x^ky'=y'+uv_1$ by the geometric-sum
identity; then $y=k'+u(v_0-v_1)$, as required.

For the second equality, if $uv\in K$, then
$u\varphi(v)=0$.  Multiplication by $u$ is injective on the
graded module $G$, by FCD9, so $v\in K$.  The reverse inclusion
is immediate.  The isomorphism theorem applied to \eqref{FCD.11}
now gives
\[
 K/uK\simeq (K+uF)/uF=J/uF.
\]
This is a quotient of the finitely generated module $J$ and is
therefore finitely generated.  No Noetherian hypothesis on $S_R$
or finite-presentation hypothesis on $F$ was used.
\end{proof}
% END FCD11

% BEGIN FCD12
\paragraph{FCD12: the full finite graded presentation up to chart torsion.}
\label{claim:FCD12}
Let $G$ be a graded $S_R$-module such that each
$G/(X_i-1)G$ is finitely presented over $S_R$.
There exist finite ranks $a,b$, integer degree lists $d,e$, a
matrix $\Phi$ with $\Phi_{lj}$ homogeneous of degree $e_j-d_l$,
and a graded map
\[
 \psi:N=\operatorname{coker}\Phi\longrightarrow G
\]
with the following two properties for every $i$:
\begin{align}
 \psi(y)=0&\ \Longrightarrow\ X_i^ky=0
       \quad\hbox{for some }k\ge0,\nonumber\\
 z\in G&\ \Longrightarrow\ X_i^kz\in\operatorname{im}\psi
       \quad\hbox{for some }k\ge0.                  \tag{FCD.12}\label{FCD.12}
\end{align}
These are the original quantifiers of
\texttt{exists\_\allowbreak{}graded\_\allowbreak{}presentation\_\allowbreak{}saturating}.

\begin{proof}
FCD10 gives finitely many homogeneous $g_l\in G_{d_l}$ whose
span has the second property in \eqref{FCD.12}.  Let
$F=\bigoplus_lS_Re_l$, with $\deg e_l=d_l$, and define the
graded map $\varphi(e_l)=g_l$.  Its kernel $K$ is graded:
if a sum of homogeneous components maps to zero, the direct-sum
decomposition in $G$ makes each component map to zero.

For every chart, FCD11 proves that $K/(X_i-1)K$ is finitely
generated.  Apply FCD10 to $K$ to choose a finite homogeneous
family $\kappa_j\in K_{e_j}$ whose span is correct up to powers
of every $X_i$.  Let $\Phi$ have these columns.  In the chosen
graded free basis its $l,j$ entry has degree $e_j-d_l$.
Since all columns lie in $K$, the map $\varphi$ descends to the
displayed graded $\psi:N\to G$.

Its image is exactly the image of $\varphi$, giving the second
assertion of \eqref{FCD.12}.  If $\psi(y)=0$, represent $y$
by $v\in F$.  Then $v\in K$.  For the chosen chart $i$, some
$X_i^kv$ is a combination of the $\kappa_j$, hence is in
$\Phi S_R^b$.  Its class $X_i^ky$ is zero.  This proves the
first assertion, for arbitrary $y$, not just homogeneous $y$.

This constructs the finite presentation needed on the charts without
asserting that $G$ itself is a finite $S_R$-module.
\end{proof}
% END FCD12

% BEGIN FCD13
\paragraph{FCD13: actual dehomogenization gives the required finite quotients.}
\label{claim:FCD13}
Let $G$ be graded, $\pi:S_R\twoheadrightarrow B$ an
$R$-algebra map with $\pi(X_i)=1$, and
$\epsilon:G\twoheadrightarrow P$ a $\pi$-semilinear map.
Suppose every homogeneous $y$ with $\epsilon(y)=0$ is killed
by some power of $X_i$.  Then
\[
 \ker\epsilon=(X_i-1)G,
 \qquad G/(X_i-1)G\simeq P                          \tag{FCD.13}\label{FCD.13}
\]
as $S_R$-modules, using $\pi$ for the action on $P$.
In particular the quotient is finite if $P$ is finite over $B$.
It is finitely presented over $S_R$ if $B$ is finitely presented
over $R$ and $P$ is finitely presented over $B$.

\begin{proof}
The inclusion $(X_i-1)G\subseteq\ker\epsilon$ follows from
semilinearity.  If $\epsilon(y)=0$, write
$y=y_N+(X_i-1)v$ with $y_N$ homogeneous, using FCD9.
Then $\epsilon(y_N)=0$, so $X_i^ky_N=0$ for some $k$.
The geometric-sum identity gives
\[
 y_N=-(X_i-1)(1+X_i+\cdots+X_i^{k-1})y_N,
\]
and $y$ also belongs to $(X_i-1)G$.  Surjectivity of
$\epsilon$ proves the isomorphism in \eqref{FCD.13}.

Because $\pi$ is surjective, a finite set of $B$-module generators
of $P$ also generates it over $S_R$.  For finite presentation,
FCD1 gives finitely many generators $g_k$ of $\ker\pi$.
Take a finite presentation $B^t\to B^s\to P\to0$ and lift its
relation columns to $S_R^s$.  These lifted columns together with
$g_ke_l$ generate the kernel of $S_R^s\to P$: reduce any
kernel vector to $B^s$, subtract the lifted relation combination,
and then express the remaining coordinates in $(g_k)$.
This gives an actual finite presentation over $S_R$ and proves
the assertion.
\end{proof}
% END FCD13

% BEGIN FCD14
\paragraph{FCD14: the same dehomogenization after the finite presentation.}
\label{claim:FCD14}
Let $\psi:N\to G$ be graded, and suppose it satisfies both
conditions \eqref{FCD.12} for a fixed $i$.  Let
$\epsilon:G\twoheadrightarrow P$ be semilinear for a ring map
$\pi:S_R\to B$ with $\pi(X_i)=1$, and suppose its homogeneous
kernel is killed by powers of $X_i$.  Then $\epsilon\psi$ is
surjective and
\[
 \epsilon\psi(y)=0\quad\Longleftrightarrow\quad
 y=(X_i-1)v\text{ for some }v\in N.                 \tag{FCD.14}\label{FCD.14}
\]
Only $N$ needs a direct-sum grading for this assertion; for $G$
it suffices that $\psi$ takes homogeneous elements to the stated
homogeneous pieces.

\begin{proof}
Given $p\in P$, choose $z\in G$ with $\epsilon(z)=p$.
By the second part of \eqref{FCD.12}, there are $k,y$ with
$\psi(y)=X_i^kz$.  Semilinearity and $\pi(X_i)=1$ give
$\epsilon\psi(y)=p$.

If $y$ is homogeneous and $\epsilon\psi(y)=0$, the chart
kernel hypothesis gives $X_i^k\psi(y)=0$.
Thus $\psi(X_i^ky)=0$, and the first part of
\eqref{FCD.12} gives $X_i^{j+k}y=0$ for some $j$.
The proof of FCD13's kernel identity, applied to the graded
module $N$ and to $\epsilon\psi$, now proves
\eqref{FCD.14}.  That proof only uses homogeneous pieces in
the target to invoke the stated kernel hypothesis, as claimed.
\end{proof}
% END FCD14

\subsection{The Noetherian graded-cube model}

For a graded $S_R$-module $M$ and a degree $n$, define its
localization cube by
\[
 C(M,n)_J=(M[X_J^{-1}])_n,\qquad J\subseteq\iota,    \tag{FCD.15}\label{FCD.15}
\]
with the actual localization maps for inclusions of subsets.  At
$J=\varnothing$ the entry is $M_n$.  This is the source's
\texttt{gr\allowbreak Cube}; the degree pieces are its \texttt{hom\allowbreak Loc}.
The ordinary \v Cech complex has
\[
 \check C^p(M,n)=\bigoplus_{|J|=p+1}C(M,n)_J,
\]
and alternating localization differentials.  We also use the
augmented cube complex $A^p(M,n)=\bigoplus_{|J|=p}C(M,n)_J$,
whose degree-zero term is $M_n$.  These conventions will keep the
top and bottom cohomology in the next argument explicit.

% BEGIN FCD15
\paragraph{FCD15: finite graded \v Cech cohomology, with its monomial proof.}
\label{claim:FCD15}
Suppose $R_0$ is Noetherian, $\iota$ is finite, and $M$ is a
finite graded $S_{R_0}$-module.  For every integer $n$ and
every $p\ge0$, the cohomology $H^p(\check C(M,n))$ is a finite
$R_0$-module.

\begin{proof}
Put $m=|\iota|$.  First compute the augmented complex for one
free generator of degree $d$.  The localized degree-$n$ terms
have Laurent monomial bases $X^\alpha$ with
\[
 \alpha\in\mathbb Z^{\iota},\qquad
 \sum_i\alpha_i=n-d,
 \qquad\{i:\alpha_i<0\}\subseteq J.
\]
Indeed localizing a polynomial ring at $X_J$ allows precisely
the negative exponents in $J$.  Every element has a unique finite
such expansion, over arbitrary $R_0$, so the complex is the
direct sum of its complexes for individual exponent vectors.

Fix $\alpha$ and put $N(\alpha)=\{i:\alpha_i<0\}$.
Its summands are one copy of $R_0$ for every subset $J$ containing
$N(\alpha)$, in degree $|J|$.  With basis $e_J$, the differential
adds one index $j\notin J$ with sign
$(-1)^{\#\{i\in J:i<j\}}$.
If $N(\alpha)\ne\iota$, choose $j_0\notin N(\alpha)$.
Define a degree-minus-one map by
\[
 h(e_J)=
 \begin{cases}
 (-1)^{\#\{i\in J:i<j_0\}}e_{J\setminus\{j_0\}},&j_0\in J,\\
 0,&j_0\notin J.
 \end{cases}
\]
The indicated subset still contains $N(\alpha)$.
In $dh+hd$, the term removing and adding $j_0$ is $e_J$;
the terms involving any other index cancel in pairs, since
exchanging the order of insertion and deletion changes its sign.
Thus $dh+hd=1$ and this monomial complex is contractible.

If $N(\alpha)=\iota$, only the top term $J=\iota$ occurs,
giving one possible top cohomology class.  There are finitely many
such $\alpha$ for fixed $n-d$.  When $m>0$, writing
$\alpha_i=-1-\beta_i$ reduces the count to nonnegative integers
with fixed sum $\sum_i\beta_i=d-n-m$.  Each is bounded by that
sum when the sum is nonnegative, and there are none otherwise.
When $m=0$, there is one empty exponent vector if $n=d$ and
none otherwise.  Consequently every cohomology group of this
augmented free complex is finite.  Finite direct sums prove the
same assertion for finite graded free modules of arbitrary shifts.

Now let $M$ be any finite graded module.  It has a finite graded
free surjection $F\to M$: replace any finite generating family
by its homogeneous components.  Its kernel $K$ is graded and
finite because $S_{R_0}$ is Noetherian.  Localization is exact,
and taking a homogeneous component of an exact sequence of
graded modules is exact.  Hence
\[
 0\longrightarrow A^\bullet(K,n)
 \longrightarrow A^\bullet(F,n)
 \longrightarrow A^\bullet(M,n)\longrightarrow0
\]
is an exact sequence of complexes.  Prove finiteness of augmented
cohomology by descending induction on $p$, uniformly for all
finite graded modules.  Above $m$ the complexes are zero.  At
degree $p$, the long exact sequence gives
\[
 H^p(A(F,n))\longrightarrow H^p(A(M,n))
 \longrightarrow H^{p+1}(A(K,n)).
\]
The first module is finite by the free calculation; the last is
finite by the induction hypothesis.  The middle module is an
extension of a quotient of the first by a submodule of the last.
Both are finite because $R_0$ is Noetherian.  This completes
the induction without any finite-projective-dimension assumption.

Finally $M_n$ is finite over $R_0$: in the chosen free surjection,
$F_n$ is a finite sum of polynomial pieces of fixed nonnegative
degree, each with finitely many monomials, and it surjects onto
$M_n$.  For $p\ge1$, ordinary \v Cech cohomology is the
augmented cohomology in degree $p+1$.  In degree zero there is
the exact sequence
\[
 0\longrightarrow\operatorname{im}(M_n\to A^1(M,n))
 \longrightarrow H^0(\check C(M,n))
 \longrightarrow H^1(A(M,n))\longrightarrow0.
\]
Both outer modules are finite.  This proves the claim, including
the case $\iota=\varnothing$, when the ordinary complex is zero.
\end{proof}
% END FCD15

% BEGIN FCD16
\paragraph{FCD16: flatness on a singleton implies flatness on every containing vertex.}
\label{claim:FCD16}
For any graded $S_R$-module $M$ there is an $R$-linear
isomorphism
\[
 (M[X_i^{-1}])_0\simeq M/(X_i-1)M.                  \tag{FCD.16}\label{FCD.16}
\]
If this module is $R$-flat, then for every $J$ containing $i$
and every integer $n$, $(M[X_J^{-1}])_n$ is $R$-flat.

\begin{proof}
Represent a degree-zero element on the left of \eqref{FCD.16}
by $y/X_i^k$, with $k\ge0$ and $y\in M_k$, and send it to
the class of $y$.  Any degree-zero fraction has such a
representative: choose a common denominator and take its
homogeneous numerator of the required degree.  Equality of two
fractions means that an additional power of $X_i$ kills their
cross-multiplied difference.  Modulo $X_i-1$, this is exactly
equality of the two numerator classes, so the map is well-defined.

If the class of $y$ is zero, then $y=(X_i-1)w$ and is
homogeneous.  FCD9 gives $X_i^ly=0$, so the original fraction
was zero.  For surjectivity, take any class represented by $z$;
FCD9 gives a homogeneous representative $z_N\in M_N$ for
some $N\ge0$.  The fraction $z_N/X_i^N$ maps to it.

Multiplication by the invertible homogeneous element $X_i$
identifies consecutive degrees of $M[X_i^{-1}]$.
Therefore flatness of its degree-zero piece gives flatness of
every piece, and their direct sum $M[X_i^{-1}]$ is $R$-flat.
Further localize by the product of the remaining $X_j$ in $J$.
Localization of an $R$-flat module by an endomorphism commuting
with $R$ is $R$-flat: it is the filtered direct limit of copies
of that module with transition given by that endomorphism, and
filtered direct limits preserve injections after tensoring.
This proves flatness of $M[X_J^{-1}]$.
Its grading is a direct sum, so every homogeneous piece is a
direct summand and is flat as well.  The grading and its directness
can be checked by multiplying a finite sum of homogeneous fractions
by a common homogeneous denominator and using directness in $M$.
\end{proof}
% END FCD16

% BEGIN FCD17
\paragraph{FCD17: canonical scalar extension and localization comparisons.}
\label{claim:FCD17}
Let $R_0\to R$ be any ring map and let a homogeneous matrix
$\Phi_0$ over $S_{R_0}$ map to $\Phi$ over $S_R$, with the
same free-module shifts.  Write $N_0=\operatorname{coker}\Phi_0$
and $N=\operatorname{coker}\Phi$.  There is a canonical
isomorphism of cubes, for every degree $n$,
\[
 R\otimes_{R_0} C(N_0,n)\simeq C(N,n),               \tag{FCD.17}\label{FCD.17}
\]
induced by coefficientwise application of the ring map to the
representatives.  If $\psi:N\to G$ is graded and satisfies
\eqref{FCD.12} for every $i$, its induced map
$C(N,n)_J\to C(G,n)_J$ is bijective for every nonempty $J$.
All these maps commute with the restriction maps between vertices.

\begin{proof}
The map $R\otimes_{R_0}S_{R_0}\to S_R$ sending
$r\otimes p$ to $r\,p_R$ is an isomorphism: the monomials
are a basis on both sides.  Apply right exactness of tensor
product to the presentation
$S_{R_0}^b\xrightarrow{\Phi_0}S_{R_0}^a\to N_0\to0$.
It gives the canonical isomorphism
$R\otimes_{R_0}N_0\simeq N$, sending
$1\otimes[v]$ to $[v_R]$.  This uses a cokernel, not a
kernel, so flatness of $R$ over $R_0$ is unnecessary.

Localizing commutes with this scalar extension.  Explicitly the
map sends $r\otimes(y/X_J^k)$ to $r\,y_R/X_J^k$;
the inverse follows from the universal property of inverting
$X_J$, or equivalently from the same filtered-limit description
used in FCD16.  These maps preserve degrees.  Taking a degree
piece also commutes with tensor product here: a graded module
is the direct sum of its pieces, tensor product commutes with
that direct sum, and its coordinate projections split the
individual pieces.  This proves \eqref{FCD.17}, including the
empty vertex.

For the second assertion, choose $i\in J$.  The kernel and
cokernel of $\psi$ have every element killed by a power of
$X_i$, by \eqref{FCD.12}.  Localizing at $X_J$ kills both,
because $X_i$ becomes invertible.  Exactness of localization
then gives an isomorphism $N[X_J^{-1}]\simeq G[X_J^{-1}]$,
and its degree-$n$ component is bijective.
Each map sends a numerator to its specified image and leaves
the localization denominators in place.  Enlarging $J$ has
the same effect before or after applying these formulas.
This proves the asserted commutation and hence the cube,
rather than merely vertexwise, statement.
\end{proof}
% END FCD17

% BEGIN FCD18
\paragraph{FCD18: the original Noetherian graded-cube model.}
\label{claim:FCD18}
Let $R$ be arbitrary and $G$ a graded $S_R$-module.  For each
$i\in\iota$, suppose there exist a smooth $R$-algebra $B_i$,
a surjective $R$-algebra map $\pi_i:S_R\to B_i$ with
$\pi_i(X_i)=1$, a finitely presented projective $B_i$-module
$P_i$, and a surjective $\pi_i$-semilinear map
$\epsilon_i:G\to P_i$.  Assume that whenever $y$ is
homogeneous and $\epsilon_i(y)=0$, some $X_i^ky$ is zero.
Then there are a single Noetherian subring $R_0\subseteq R$,
a cube $C_0$ of $R_0$-modules, and an actual cube map
\[
 \theta:R\otimes_{R_0}C_0\longrightarrow C(G,0)       \tag{FCD.18}\label{FCD.18}
\]
such that:
\begin{enumerate}
\item $(C_0)_J$ is $R_0$-flat for every nonempty $J$;
\item every $H^p(\check C(C_0))$ is finite over $R_0$;
\item $\theta_J$ is bijective for every nonempty $J$.
\end{enumerate}
This is the full original
\texttt{exists\_\allowbreak{}noetherian\_\allowbreak{}model\_\allowbreak{}gr\allowbreak{}Cube}, with its original
chart hypotheses and quantifier order.

\begin{proof}
FCD13 makes every $G/(X_i-1)G$ finitely presented over $S_R$.
FCD12 therefore supplies one homogeneous matrix $\Phi$ and
one graded $\psi:N=\operatorname{coker}\Phi\to G$ satisfying
\eqref{FCD.12} on all charts.
Let $q:S_R^a\to N$ be its quotient map.
For each $i$, FCD14 says that $\epsilon_i\psi$ is surjective
with kernel $(X_i-1)N$.  Its composite with $q$ consequently
has the exact kernel
\[
 \ker(\epsilon_i\psi q)=\Phi S_R^b+(X_i-1)S_R^a.
\]
Indeed $q(y)=(X_i-1)q(v)$ is equivalent to
$y-(X_i-1)v\in\Phi S_R^b$.  Thus the ORIGINAL $B_i,P_i$
and these composite maps satisfy every hypothesis of FCD8.

Apply FCD8 to obtain one $R_0$ and $\Phi_0$.  Give
$N_0=\operatorname{coker}\Phi_0$ the same quotient grading,
and define, rather than choose independently,
\[
 C_0=C(N_0,0).
\]
FCD17 supplies the canonical first map in the composite
\[
 R\otimes_{R_0}C(N_0,0)
 \xrightarrow{\ \simeq\ } C(N,0)
 \xrightarrow{C(\psi,0)} C(G,0).
\]
Define $\theta$ to be exactly this composite.

For a nonempty $J$, choose $i\in J$.
FCD8 makes $N_0/(X_i-1)N_0$ flat over $R_0$; FCD16
therefore makes $(N_0[X_J^{-1}])_0=(C_0)_J$ flat.
The module $N_0$ is finite over the Noetherian polynomial ring
$S_{R_0}$, so FCD15 gives finite \v Cech cohomology in every
degree.  Finally FCD17 makes both factors of $\theta_J$
bijective for every nonempty $J$.  The maps commute with all
restrictions by that same claim, so the conclusion concerns
the actual cube map, as required.

The choices $\Phi,\psi,R_0,\Phi_0,C_0,\theta$ have all
been made for the fixed original $G$ and its chart data.
No later coefficient algebra has been introduced.  In particular,
the original map $R_0\to R$ need not be flat.
\end{proof}
% END FCD18

\subsection{Exact use in the existing section proof}

FCD18 now supplies the Noetherian-cube input previously retained in
GSL8 whenever its original graded chart hypotheses are available.
Applying GSL8 to $C(G,0)$ gives one bounded complex of finite
projective $R$-modules and one map to its \v Cech complex that
remains a quasi-isomorphism after every $R\to R'$.  This is a
direct reuse of the already proved complex algebra.

For the geometric application, the intended $G$ is the graded sum of
the actual global sections of $M\otimes H^t$, and the $X_i$ act by
the fixed sections $s_i$ of the very ample bundle $H$.
The remaining producers must identify the chart algebra with
$\Gamma(D(s_i),\mathcal O)$ and its module with the actual
sections of $M\otimes H^0$ there, prove the required smoothness,
finite presentation and projectivity, and prove the dehomogenization
and actual-section comparison.  They must produce the maps on this
same graded module and the same cover.  After these are supplied,
the source's final map is the composite of $\theta$, the geometric
localization comparison, and the natural identification
$M\otimes H^0=M$; GSL8--GSL10 then give every literal
\texttt{Perf} field.  The current section does not replace these
geometric arguments by an abstract identification of dimensions.

The other retained GSL producers---polarized base change, geometric
$H^1$ vanishing, geometric integrality and principal rank one,
torsion-section decomposition and growth, and the multiplication
and square generators---are unchanged.  The original \texttt{Moduli}
record and the HUF/PUN foundations also remain.  Neither a global
section law nor a selected whole RH/GRH obligation is closed by the
present algebraic model theorem.

\subsection{Source correspondence}

All following paths are relative to
\texttt{src/\allowbreak Development/\allowbreak Abelian\allowbreak Variety\allowbreak Theory/\allowbreak Proofs/\allowbreak } at the
accepted source pin.  In this list \texttt{C/\allowbreak } abbreviates
\texttt{Coherent\allowbreak Cohomology\allowbreak Proofs/\allowbreak }.  The local snapshots retain the
complete files; the provenance records bind the stated ranges and
each numbered proof block.

\begin{itemize}
\item FCD1--FCD3 and FCD7--FCD8:
\texttt{C/\allowbreak{}Exists\allowbreak{}Noetherian\allowbreak{}Subring\allowbreak{}Homogeneous\allowbreak{}Matrix\allowbreak{}Flat\allowbreak{}Char.lean},
lines 11--142 and 186--287.  FCD1 expands the finite-presentation
facts used at the start of the chart-witness proof.
\item FCD4:
\texttt{C/\allowbreak{}Exists\allowbreak{}Witnesses\allowbreak{}Of\allowbreak{}Formally\allowbreak{}Smooth\allowbreak{}Quotient.lean},
lines 9--72, and the converse certificate in
\texttt{C/\allowbreak{}Exists\allowbreak{}Noetherian\allowbreak{}Subring\allowbreak{}Homogeneous\allowbreak{}Matrix\allowbreak{}Flat\allowbreak{}Char.lean},
lines 145--182.
\item FCD5 expands the Noetherian flatness consequence used in that
same file, lines 253--268.  The dependency pin recorded by the
repository's \texttt{lake-manifest.\allowbreak json} is Mathlib
\texttt{0df444a360eaa60ab8c11dca51a86af692955474}; the supporting
files are \texttt{Mathlib/\allowbreak{}Ring\allowbreak{}Theory/\allowbreak{}Smooth/\allowbreak{}Flat.lean},
\texttt{Smooth/\allowbreak{}Adic\allowbreak{}Completion.lean}, and
\texttt{Adic\allowbreak{}Completion/\allowbreak{}As\allowbreak{}Tensor\allowbreak{}Product.lean} under
\texttt{Mathlib/\allowbreak Ring\allowbreak Theory/\allowbreak }.
\item FCD6:
\texttt{C/\allowbreak{}Flat\allowbreak{}Quot\allowbreak{}SMul\allowbreak{}Top\allowbreak{}Coker\allowbreak{}Of\allowbreak{}Matrix\allowbreak{}Witnesses.lean}, lines 10--117.
\item FCD9--FCD12:
\texttt{C/\allowbreak{}Exists\allowbreak{}XPow\allowbreak{}Smul\allowbreak{}Eq\allowbreak{}Zero\allowbreak{}Of\allowbreak{}Eq\allowbreak{}XSub\allowbreak{}One\allowbreak{}Smul.lean},
\texttt{C/\allowbreak{}Exists\allowbreak{}XPow\allowbreak{}Smul\allowbreak{}Mem\allowbreak{}Span\allowbreak{}Of\allowbreak{}Finite\allowbreak{}Quot.lean},
\texttt{C/\allowbreak{}Exists\allowbreak{}Homogeneous\allowbreak{}Family\allowbreak{}Saturating.lean},
\texttt{C/\allowbreak{}Exists\allowbreak{}Pow\allowbreak{}Smul\allowbreak{}Eq\allowbreak{}Add\allowbreak{}Sub\allowbreak{}One\allowbreak{}Smul.lean}, and
\texttt{C/\allowbreak{}Exists\allowbreak{}Graded\allowbreak{}Presentation\allowbreak{}Saturating.lean}.
These are respectively the degree recurrence, finite chart generation,
homogeneous union, finite relation module and full presentation.
\item FCD13--FCD14:
\texttt{C/\allowbreak{}Dehomogenising\allowbreak{}Eq\allowbreak{}Zero\allowbreak{}Iff.lean},
\texttt{C/\allowbreak{}Finite\allowbreak{}Quot\allowbreak{}SMul\allowbreak{}Top\allowbreak{}XSub\allowbreak{}One\allowbreak{}Of\allowbreak{}Chart.lean}, and
\texttt{C/\allowbreak{}Dehomogenising\allowbreak{}Comp\allowbreak{}Of\allowbreak{}Saturating.lean}, lines 22--42.
\item FCD15:
\texttt{C/\allowbreak{}Cube\allowbreak{}Complex\allowbreak{}Poly\allowbreak{}Grading\allowbreak{}Finite\allowbreak{}Exact.lean},
\texttt{C/\allowbreak{}Cube\allowbreak{}Complex\allowbreak{}Gr\allowbreak{}Free\allowbreak{}Finite\allowbreak{}And\allowbreak{}Exact\allowbreak{}At.lean},
\texttt{C/\allowbreak{}Finite\allowbreak{}Homology\allowbreak{}Cube\allowbreak{}Complex\allowbreak{}Graded\allowbreak{}Module.lean},
\texttt{C/\allowbreak{}Finite\allowbreak{}Gr\allowbreak{}Cech\allowbreak{}Homology.lean}, and the explicit monomial
and \v Cech comparison declarations in the manifest.
\item FCD16--FCD17:
\texttt{C/\allowbreak{}Nonempty\allowbreak{}Hom\allowbreak{}Loc\allowbreak{}Singleton\allowbreak{}Equiv\allowbreak{}Quot\allowbreak{}SMul\allowbreak{}Top.lean},
\texttt{C/\allowbreak{}Flat\allowbreak{}Hom\allowbreak{}Loc\allowbreak{}Of\allowbreak{}Flat\allowbreak{}Hom\allowbreak{}Loc\allowbreak{}Singleton.lean},
\texttt{C/\allowbreak{}Is\allowbreak{}Base\allowbreak{}Change\allowbreak{}Coker\allowbreak{}Matrix\allowbreak{}Mv\allowbreak{}Polynomial.lean}, and
\texttt{C/\allowbreak{}Nonempty\allowbreak{}Base\allowbreak{}Change\allowbreak{}Gr\allowbreak{}Cube\allowbreak{}Iso.lean}.
\item FCD18:
\texttt{C/\allowbreak{}Dehomogenising\allowbreak{}Comp\allowbreak{}Of\allowbreak{}Saturating.lean}, lines 44--114.
The next geometric consumer is
\texttt{C/\allowbreak{}Exists\allowbreak{}Finite\allowbreak{}Projective\allowbreak{}Cech\allowbreak{}Complex\allowbreak{}Line\allowbreak{}Bundle.lean},
lines 19--43; the already reconstructed algebra it invokes is
GSL8--GSL10.
\end{itemize}


\section{Finite characteristic points and propagation across components}
\label{sec:arithmetic-component-propagation}

\paragraph{What this chapter proves.}
The arithmetic construction uses two modularity realizations on one
frame, one in rank two and one in rank three.  Their modularity ideals
belong to different rings.  Component propagation makes the rank-two
ideal vanish on every relevant component, uses that information in
rank three, and produces an ideal controlling every nontriangular
rank-three point.  The last step is the original theorem
\(\texttt{Fin\allowbreak Pt\allowbreak Of}\): that ideal forces the characteristic polynomial
of every permitted characteristic-\(\ell\) point to be the prescribed
residual polynomial.

We prove this last theorem in full, including its original
\(\texttt{Is\allowbreak CD}\) test rings.  We also prove the original local and
global component-propagation principles, their transfer through flat
local charts, and the one-dimensional prime-avoidance theorem used
to find the local propagation points.  These are substantive inputs
to the two rank-specific propagation theorems.  The separate
deformation-space dimension, connectedness and rank-two-to-rank-three
corner constructions are stated explicitly at the end; their
conclusions are not assumed to be already reconstructed here.

There is one useful simplification in the finite-point argument.
The source realizes a point whose reduction may be triangular.
The actual \(\texttt{Fin\allowbreak Pt}\) consumer assumes the stronger,
literal condition \(\texttt{Is\allowbreak Lift\allowbreak A}\), which includes diagonal
reduction.  For these points the elementary subring
\(k+\mathfrak m_A\) gives the required realization directly.
No arbitrary coefficient extension or new frame is chosen.

\paragraph{Notation and the supplied universal space.}
Let \(O\) be the original complete discrete valuation ring with
finite residue field \(k\) of characteristic \(\ell\), and let
\(F\) be an original frame satisfying \(F.\mathrm{Ok}(C_0)\).
Write \(G_0=F.G_0\), \(G_v=F.G_v\), \(I_v=F.I_v\), and
\(\Psi=F.\Psi\).  The local maps are
\(\operatorname{loc}_{0,v}:G_v\to G_0\); their two branches are
\(\operatorname{loc}_{2,(v,0)}=\operatorname{loc}_{0,v}\) and
\(\operatorname{loc}_{2,(v,1)}=\operatorname{conj}\circ
\operatorname{loc}_{0,v}\).  Put
\[
 w(2)=(1,-1),\qquad w(3)=(1,0,-1).
                                                        \tag{ACP.1}
\]
The frame involution has
\(g^*=\operatorname{conj}(g^{-1})\) and
\(\Psi(g^*)=\Psi(g)^{-1}\).  Let \(J_n\) be the matrix with
ones on the antidiagonal and zeros elsewhere.

An \(\mathrm{NRep}(F,n,C)\) consists of a representation
\(r:G_0\to\mathrm{GL}_n(C)\) and characters
\(\psi_{v,i}:G_v\to C^\times\), with
\[
 r(g^*)=J_nr(g)^tJ_n,
 \qquad r(F.g_0)-r((F.g_0)^*)\ \text{diagonal},
                                                        \tag{ACP.2}
\]
and the ordered identity
\[
 \prod_{i=0}^{n-1}
 \bigl(r(\operatorname{loc}_{0,v}s_i)
                         -\psi_{v,i}(s_i)I_n\bigr)=0
 \quad\text{for every }(s_0,\ldots,s_{n-1})\in G_v^n.
                                                        \tag{ACP.3}
\]
Products in (ACP.3) have the displayed order.
Adic continuity means that every matrix entry and every character
value is locally constant modulo each power of the maximal ideal.
The condition \(\mathrm{IsLiftA}(O,r,\psi)\) means this continuity,
the specified residual characters \(\Psi^{w(n)_i}\), and zero
residual off-diagonal entries.  In particular it says more than
equality of residual semisimplifications.

A supplied \(S:\mathrm{Space}(O,F,n)\) has a complete Noetherian
local ring \(R\) with residue field \(k\), and a universal
\((r^{\mathrm{univ}},\psi^{\mathrm{univ}})\) for these continuous
diagonal-residual lifts over complete Noetherian local
\(O\)-algebras with residue field \(k\).  Write
\[
 \begin{split}
 \mathcal B&=\overline{O[\text{weight-zero entry monomials},
                         \ \psi^{\mathrm{univ}}_{v,i}(s)]},\\
 \mathcal A&=\mathcal B[\,r^{\mathrm{univ}}(g)_{ij}\,]
                                          \ \subset R.
 \end{split}                                             \tag{ACP.4}
\]
The closure is the \(\mathfrak m_R\)-adic closure.  The weight of
an entry \((i,j)\) is \(w(n)_i-w(n)_j\).  Thus a monomial belongs
to the displayed weight-zero set precisely when its entry weights
sum to zero.  Let \(u\) denote the universal representation over
\(\mathcal A\), let \(\iota:\mathcal B\hookrightarrow\mathcal A\)
be the inclusion, and set
\[
 \begin{split}
 \mathfrak m_{\mathcal B}&=\mathcal B\cap\mathfrak m_R,
 \quad\pi(\mathfrak p)=\iota^{-1}(\mathfrak p),\\
 \lambda&=\mathfrak m_O\mathcal B+
  \bigl(\psi_{v,i}^{\mathrm{univ}}(s)-1:
       s\in I_v,\ \Psi(\operatorname{loc}_v s)=1\bigr).
 \end{split}                                             \tag{ACP.5}
\]
The exact predicate \(S.\mathrm{FinLam}(I)\) is
\(\mathfrak m_{\mathcal B}\subset\sqrt{I+\lambda}\).

Let \(\mathfrak t_+\), respectively \(\mathfrak t_-\), be the
ideals of \(\mathcal A\) generated by all entries below,
respectively above, the diagonal.  A prime \(\mathfrak p\) is
\emph{off triangular} if neither ideal is contained in
\(\mathfrak p\).  The original condition \(\mathrm{OffFin}(S)\)
asserts the existence of one ideal \(I\subset\mathcal B\) with
\[
 \mathfrak m_{\mathcal B}\subset\sqrt{I+\lambda},
 \qquad I\subset\pi(\mathfrak p)
 \quad\text{for every off-triangular prime }\mathfrak p.
                                                        \tag{ACP.6}
\]
This same ideal must work before the later test ring and point are
chosen.

We use elementary local algebra, the principal ideal theorem,
finite prime avoidance, the Krull intersection theorem for
Noetherian local rings, and flat going down.  The arguments that
apply these facts, and the relevant constructions, are given below.
The frame's local generator and uniform class-exponent fields are
the literal fields reconstructed in LFC3 and LFC6; they are not
replaced by new finiteness assumptions depending on a test point.

\subsection{Realizing the exact test points without a larger construction}

% BEGIN ACP1
\paragraph{ACP1. The complete ring with the original residue constants.}
\label{claim:ACP1}
Let \(A\) satisfy the original \(\mathrm{IsCD}(O,A)\), and
suppose the specified maps form a scalar tower
\(O\to k\overset{j}{\longrightarrow}A\).
Thus \(A\) is a complete discrete valuation domain, its residue
field \(k_A\) is finite, and the given \(O\)-map is local.
Then
\[
 C=j(k)+\mathfrak m_A
   =\{a\in A:\overline a\in\overline j(k)\subset k_A\}
                                                        \tag{ACP.7}
\]
is a complete Noetherian local \(O\)-algebra with residue field
\(k\).  Its maximal ideal is \(\mathfrak m_A\), viewed as a
subset of \(C\), and
\[
 \mathfrak m_C^N=C\cap\mathfrak m_A^N\quad(N\ge0),
 \qquad \mathfrak m_C^N=\mathfrak m_A^N\quad(N\ge1)
                                                        \tag{ACP.8}
\]
under the inclusion in \(A\).  The \(O\)-structure is the
specified one.  In particular \(C\) satisfies the original
\(\mathrm{IsCO}(O,C)\) contract.

\begin{proof}
A unital map from the field \(k\) into the domain \(A\) is
injective.  Its composite with \(A\to k_A\) is also injective:
a nonzero element of \(k\) has its inverse already in \(A\),
so cannot lie in \(\mathfrak m_A\).  We may identify \(k\)
with these two images.  Since \(k_A\) is a finite set,
\([k_A:k]\) is finite.  This is supplied by \(\mathrm{IsCD}\);
no finite-degree hypothesis has been added.

The inverse image description in (ACP.7) makes \(C\) a subring.
If \(c\in C\setminus\mathfrak m_A\), its inverse in \(A\)
has residue \(\overline c^{-1}\in k\), hence also lies in \(C\).
Therefore \(C\) is local and \(\mathfrak m_C=\mathfrak m_A\).
Its residue map has image exactly \(k\), and the map from \(O\)
induces its prescribed residue isomorphism.

Choose a uniformizer \(t\) of \(A\).  For \(N\ge1\), every
element \(t^Na\in\mathfrak m_A^N\) is the product of \(N-1\)
copies of \(t\in\mathfrak m_C\) and the element
\(ta\in\mathfrak m_C\).  This proves the nontrivial inclusion
in (ACP.8); the reverse inclusion is immediate.  The assertion
for \(N=0\) is its intersection version.  The subring \(C\) is
closed in \(A\), since it is the inverse image of a subset of the
discrete residue field.  Equation (ACP.8) identifies its adic
Cauchy sequences and their limits with those in this closed
subring.  Hence \(C\) is complete and separated.

Here is a direct Noetherian argument.  Evaluation at \(t\) embeds
\(k[[T]]\) in \(A\).  Indeed, the first nonzero coefficient of
a series gives its exact \(t\)-adic order, so evaluation is
injective.  Let \(b_1,\ldots,b_d\in A\) lift a \(k\)-basis of
\(k_A\).  Every \(a\in A\) has a convergent expansion
\[
 a=\sum_{n\ge0}t^n\sum_{i=1}^d c_{n,i}b_i,
 \qquad c_{n,i}\in k.
                                                        \tag{ACP.9}
\]
To construct it, match the residue of \(a\) by a linear
combination of the \(b_i\), subtract, divide by \(t\), and
repeat.  Completeness makes the remainders tend to zero.
Regrouping the finite number of basis vectors shows that \(A\)
is a finite module over \(k[[t]]\).  The ring \(C\) contains
\(k[[t]]\) and is a submodule of that finite module.  Since
\(k[[T]]\) is Noetherian, \(C\) is a finite Noetherian module
over it, and therefore is a Noetherian ring.  If the zero-maximal-
ideal case is allowed in a DVR convention, then \(A=k_A\) and
\(C=k\), and all these assertions are immediate.

Finally \(O\to C\) factors as its original residue map followed
by the specified embedding of \(k\).  It is local, and every
element of \(C\) differs from an element of \(O\) by an element
of \(\mathfrak m_C\).  Together with completeness and
Noetherianity these are exactly all the \(\mathrm{IsCO}\) fields.
\end{proof}
% END ACP1

% BEGIN ACP2
\paragraph{ACP2. Every original finite-point test is realized on the supplied space.}
\label{claim:ACP2}
Let \(n=2\) or \(3\), let \(S:\mathrm{Space}(O,F,n)\) be
supplied, and let \(A\) and its coefficient structures be as in
ACP1.  For every \((r,\psi):\mathrm{NRep}(F,n,A)\) satisfying
\(\mathrm{IsLiftA}(O,r,\psi)\), there is an \(O\)-algebra map
\[
 \xi:\mathcal A\longrightarrow A
 \quad\text{with}\quad \xi_*u=(r,\psi).
                                                        \tag{ACP.10}
\]
Moreover \(\mathfrak m_{\mathcal B}\mathcal A\) maps into
\(\mathfrak m_A\).  This is the exact realization needed for
the original \(\mathrm{FinPt}\) consumer, on its original space.

\begin{proof}
Use the explicitly constructed \(C=k+\mathfrak m_A\).
Every off-diagonal matrix entry lies in \(\mathfrak m_A\),
by the residual-diagonal field of \(\mathrm{IsLiftA}\).
Each diagonal entry has residue
\(\Psi(g)^{w(n)_i}\in k\), so lies in \(C\).
Each full ordinary character value has residue
\(\Psi(\operatorname{loc}_v s)^{w(n)_i}\in k\), so also lies
in \(C\).  Its inverse is the character value at \(s^{-1}\),
and lies in \(C\) as well.  Thus these are characters to
\(C^\times\), not merely functions with values in the subring.

The representation and all its normalization, polarization and
ordered-product identities restrict to \(C\), because
\(C\hookrightarrow A\) is injective.  Equation (ACP.8) shows
that local constancy modulo powers of \(\mathfrak m_A\)
is exactly the required local constancy modulo powers of
\(\mathfrak m_C\) for these functions.  Reduction to the
residue field of \(C\) is the original diagonal representation
over \(k\), and the ordinary characters have their original
residual values.  This is an \(\mathrm{IsLift}(O,-)\) over
the \(\mathrm{IsCO}\)-algebra \(C\).

The universal property included in the supplied \(S\) therefore
gives \(f:R\to C\) carrying its universal representation to
this restricted one.  Restrict \(f\) to \(\mathcal A\) and
include \(C\) in \(A\); this is \(\xi\).
No assertion of existence of \(S\) has been used here: it was
part of the input before the test ring was chosen.

For the last assertion, an \(O\)-algebra map between two
\(\mathrm{IsCO}\)-algebras preserves their residue maps.
To see this directly, let \(x\in\mathfrak m_R\).  If \(f(x)\)
had nonzero residue, choose \(a\in O^\times\) with that residue.
Then \(x-a\) would be a unit in \(R\), while
\(f(x)-a\in\mathfrak m_C\) would be a nonunit, a contradiction.
It follows that \(f(\mathfrak m_R)\subset\mathfrak m_C\).
The stated image condition follows from
\(\mathfrak m_{\mathcal B}\subset\mathfrak m_R\).
\end{proof}
% END ACP2

\subsection{The local polynomial identity and rigidity of diagonal characters}

% BEGIN ACP3
\paragraph{ACP3. An ordered product with a separator gives the ordered flag.}
\label{claim:ACP3}
Let \(C\) be any commutative local ring, \(H\) any group,
\(r:H\to\mathrm{GL}_n(C)\), and
\(\psi_i:H\to C^\times\) for \(0\le i<n\).
Suppose an element \(s_0\in H\) satisfies
\[
 \psi_i(s_0)-\psi_j(s_0)\in C^\times\quad(i\ne j),
 \qquad
 \det(X-\overline{r(s_0)})
       =\prod_i(X-\overline{\psi_i(s_0)}),
                                                        \tag{ACP.11}
\]
and suppose (ACP.3), with \(H\) in place of \(G_v\), holds
for every tuple.  There is one basis of \(C^n\) in which
every \(r(s)\) is upper triangular with ordered diagonal
\((\psi_0(s),\ldots,\psi_{n-1}(s))\).
This proves the full original \(\texttt{flag\_\allowbreak loc}\) statement;
no domain, completeness or continuity hypothesis is necessary.

\begin{proof}
Put \(T=r(s_0)\) and \(a_i=\psi_i(s_0)\).
The constant tuple in the ordered-product identity gives
\(\prod_i(T-a_i)=0\).  Define
\[
 E_i=\prod_{j\ne i}(T-a_j)(a_i-a_j)^{-1}.
                                                        \tag{ACP.12}
\]
Lagrange interpolation modulo the monic polynomial
\(\prod_i(X-a_i)\) gives
\(\sum E_i=1\), \(E_iE_j=0\) for \(i\ne j\),
\(E_i^2=E_i\), and \(TE_i=a_iE_i\).
These polynomial identities hold over \(C\), since the displayed
differences are units.  Modulo the maximal ideal, (ACP.11) has
distinct roots and each \(\overline E_i\) has a one-dimensional
image.  Choose lifts \(v_i\) of a basis of those images.
The vectors \(E_iv_i\) reduce to a basis, so the matrix having
them as columns has unit determinant.  They form a basis of
\(C^n\), in which the \(E_i\)'s are the coordinate projectors.

Fix \(s\in H\) and indices \(b\le a\).  In (ACP.3) put
\(s\) in position \(a\) and \(s_0\) in every other position.
Multiply on the left by \(E_a\) and on the right by \(E_b\).
Each factor before position \(a\) becomes the unit scalar
\(a_a-a_j\), and each factor after it becomes the unit scalar
\(a_b-a_j\).  Cancelling their product yields
\[
          E_a\bigl(r(s)-\psi_a(s)\bigr)E_b=0.
                                                        \tag{ACP.13}
\]
For \(b<a\), orthogonality gives \(E_ar(s)E_b=0\);
for \(b=a\), it gives \(E_ar(s)E_a=\psi_a(s)E_a\).
These are exactly the lower-zero and ordered-diagonal assertions
in the chosen basis, simultaneously for every \(s\).
For \(n=0\), the unique basis of the zero module gives the same
conclusion directly.
\end{proof}
% END ACP3

% BEGIN ACP4
\paragraph{ACP4. The literal local characteristic polynomial on the universal space.}
\label{claim:ACP4}
For every supplied \(S:\mathrm{Space}(O,F,n)\), with
\(n=2\) or \(3\), one has
\[
 \det\bigl(X-u(\operatorname{loc}_{0,v}s)\bigr)
       =\prod_i\bigl(X-\psi^{\mathrm{univ}}_{v,i}(s)\bigr)
                  \quad\text{in }\mathcal A[X]
                                                        \tag{ACP.14}
\]
for every \(v,s\).  This is the original \(\texttt{o1}\)
assertion, and it remains true after every algebra map out of
\(\mathcal A\).

\begin{proof}
The \(f3\) field of \(F.\mathrm{Ok}(C_0)\) gives
\(s_0\in G_v\) with
\(\Psi(\operatorname{loc}_v s_0)^2\ne1\).
For the two lists of weights (ACP.1), the corresponding powers
of this value are pairwise distinct.  Thus the differences of
\(\psi^{\mathrm{univ}}_{v,i}(s_0)\) are units of the local
ring \(R\).  The residual-diagonal field of the universal lift
gives (ACP.11), with \(C=R\).  Apply ACP3 to the restricted
universal representation and its exact local characters.
The characteristic polynomial of an upper-triangular matrix is
the product of its diagonal factors, proving (ACP.14) over \(R\).
All its coefficients belong to \(\mathcal A\), and
\(\mathcal A\hookrightarrow R\) is injective, so the identity
holds there.  Polynomial evaluation under any later algebra map
preserves the identity; no flatness of that map is needed.
\end{proof}
% END ACP4

% BEGIN ACP5
\paragraph{ACP5. Principal units have no torsion in a characteristic-prime local domain.}
\label{claim:ACP5}
Let \(C\) be a local domain of characteristic \(p>0\), with
\(p\) prime.  If \(x-1\in\mathfrak m_C\) and \(x^h=1\)
for an integer \(h>0\), then \(x=1\).
This is the full original \(\texttt{one\allowbreak Unit\allowbreak TF}\) assertion.

\begin{proof}
Write \(h=p^am\) with \(p\nmid m\), and put \(y=x^{p^a}\).
Then \(y^m=1\) and \(y-1=(x-1)^{p^a}\in\mathfrak m_C\).
The geometric sum \(1+y+\cdots+y^{m-1}\) has residue \(m\),
which is a nonzero element of the prime field and hence a unit.
The identity
\[
          (y-1)(1+y+\cdots+y^{m-1})=y^m-1=0
\]
gives \(y=1\).  Thus \((x-1)^{p^a}=0\); the domain
hypothesis gives \(x=1\).
\end{proof}
% END ACP5

% BEGIN ACP6
\paragraph{ACP6. A bounded finite-character quotient determines the whole character.}
\label{claim:ACP6}
Let \(G\) be a topological group, \(T\subset G\), and \(h>0\).
Assume every homomorphism from \(G\) to a finite abelian
\(p\)-group, with open kernel and trivial on \(T\), has exponent
dividing \(h\).  Let \(C\) be a Noetherian local domain of
characteristic \(p\) with finite residue field.  Suppose
\(a,b:G\to C^\times\) satisfy:
\(a\) is adically continuous, \(\ker b\) is open,
\(a\equiv b\pmod{\mathfrak m_C}\), and \(a=b\) on \(T\).
Then \(a=b\) on \(G\).  Completeness of \(C\) is not required.

\begin{proof}
Put \(\chi=ab^{-1}\).  Its values lie in
\(1+\mathfrak m_C\).  Each ring
\(C/\mathfrak m_C^N\) is finite: its successive ideal
quotients are finite-dimensional vector spaces over the finite
residue field.  The image \(B_N\) of \(\chi\) in its unit
group is an abelian \(p\)-group.  Indeed, for
\(z\in\mathfrak m_C\), the characteristic-\(p\) identity
\((1+z)^{p^N}=1+z^{p^N}\) and \(p^N\ge N\) show that its
reduction is annihilated by \(p^N\).

The character \(\chi_N:G\to B_N\) has open kernel.
To verify this without any unstated topology on unit groups,
intersect the open kernel of \(b\) with a neighborhood of
\(1\) on which \(a\) is \(1\) modulo \(\mathfrak m_C^N\).
This intersection lies in \(\ker\chi_N\); a subgroup containing
an open identity neighborhood is open.  Moreover \(\chi_N\)
is trivial on \(T\).  By the uniform assumption,
\[
                  \chi(g)^h-1\in\mathfrak m_C^N
                  \quad\text{for every }N,g.
\]
Krull intersection in the Noetherian local domain gives
\(\chi(g)^h=1\).  ACP5, applied to \(\chi(g)\), gives
\(\chi(g)=1\) for every \(g\), as required.
\end{proof}
% END ACP6

% BEGIN ACP7
\paragraph{ACP7. The frame determines every multiplicative diagonal entry.}
\label{claim:ACP7}
Let \(C\) be a Noetherian local domain with finite residue field,
let \(j:k\to C\) be a ring map, and let
\((r,\psi):\mathrm{NRep}(F,3,C)\) be adically continuous.
Suppose, for every \(g,h\in G_0\) and \(i\),
\[
 \begin{split}
 r(gh)_{ii}&=r(g)_{ii}r(h)_{ii},\\
 r(g)_{ii}&\equiv j(\Psi(g)^{w_i})\pmod{\mathfrak m_C},\\
 s\in I_v,\ \Psi(\operatorname{loc}_v s)=1
       &\Longrightarrow r(\operatorname{loc}_{0,v}s)_{ii}=1.
 \end{split}                                             \tag{ACP.15}
\]
Then \(r(g)_{ii}=j(\Psi(g)^{w_i})\) for every \(g,i\).
This is the full diagonal-rigidity step used in
\(\texttt{Fin\allowbreak Pt\allowbreak Of}\).

\begin{proof}
The embedding \(j\) makes \(C\) a ring of characteristic
\(\ell\).  Fix \(i\).  Multiplicativity and the equation at
\(g^{-1}\) make \(a_i(g)=r(g)_{ii}\) a character to
\(C^\times\), with inverse \(r(g^{-1})_{ii}\).
The reference character is \(b_i(g)=j(\Psi(g)^{w_i})\).
Its kernel is open, since it contains the open kernel of the
frame character.  The first two requirements of ACP6 are now
the given continuity and congruence.

Choose the local detectors \(u_{v,a}\) from the exact
\(\mathrm{gen}\) field, and the exponent \(h>0\) from
\(\mathrm{cl}\).  These are fixed by the frame before \(C\)
or \(a_i\) is considered.  Put
\[
 T=\{\operatorname{loc}_{2,(v,b)}(u_{v,a}):
                        v\in V,\ b\in\{0,1\},\ a<d_v\}.
                                                        \tag{ACP.16}
\]
A finite abelian \(\ell\)-group character with open kernel
which kills \(T\) kills each full local inertia image:
apply \(\mathrm{gen}\) to its pullback along each continuous
\(\operatorname{loc}_{2,(v,b)}\).  The \(\mathrm{cl}\) field
then gives the same exponent \(h\) for that character.  Thus
the finite-character hypothesis of ACP6 holds for (ACP.16).

For branch \(b=0\), the third equation of (ACP.15) applies
to every detector because \(u_{v,a}\in I_v\) and its residual
frame character is \(1\).  For branch \(b=1\), use
\(\operatorname{conj}(g)=(g^{-1})^*\) in (ACP.2):
\[
 r(\operatorname{conj}g)_{ii}=r(g^{-1})_{2-i,2-i}.
\]
When all diagonal values at \(g\) are one, their inverse
values on the right are also one.  The reference character is
one on both branches, by the corresponding frame identity
for \(\Psi\).  Hence \(a_i=b_i\) on \(T\).
ACP6 gives equality everywhere, separately for each \(i\),
with no change in the exponent or detector choices.
\end{proof}
% END ACP7

\subsection{The complete finite-point theorem}

% BEGIN ACP8
\paragraph{ACP8. The five cycle products determine a three-by-three characteristic polynomial.}
\label{claim:ACP8}
For a matrix \(M=(m_{ij})\) over any commutative ring, suppose
\[
 m_{01}m_{10}=m_{02}m_{20}=m_{12}m_{21}=0,
 \qquad m_{01}m_{12}m_{20}=m_{02}m_{21}m_{10}=0.
                                                        \tag{ACP.17}
\]
Then \(\det(X-M)=\prod_{i=0}^2(X-m_{ii})\).

\begin{proof}
Expand the determinant over its six permutations.  The identity
permutation contributes \(\prod_i(X-m_{ii})\).  The three
transpositions contribute respectively
\[
 -(X-m_{22})m_{01}m_{10},\quad
 -(X-m_{11})m_{02}m_{20},\quad
 -(X-m_{00})m_{12}m_{21},
\]
and the two three-cycles contribute
\(-m_{01}m_{12}m_{20}\) and \(-m_{02}m_{21}m_{10}\).
They all vanish by (ACP.17).
\end{proof}
% END ACP8

% BEGIN ACP9
\paragraph{ACP9. Off-triangular test points collapse the weight-zero algebra.}
\label{claim:ACP9}
Let \(S:\mathrm{Space}(O,F,3)\) satisfy
\(\mathrm{OffFin}(S)\).  Let \(A\) have the original
\(\mathrm{IsCD}\) and scalar-tower structures, and let
\((r,\psi)\) be an \(\mathrm{IsLiftA}\) with
\[
 s\in I_v,\ \Psi(\operatorname{loc}_v s)=1
                         \Longrightarrow\psi_{v,i}(s)=1.
                                                        \tag{ACP.18}
\]
Choose \(\xi\) by ACP2.  If
\(\ker\xi\subset\mathcal A\) is off triangular, then
\[
 \xi(b)=j(\operatorname{res}_R b)\quad(b\in\mathcal B).
                                                        \tag{ACP.19}
\]
Consequently every weight-zero product of entries of \(r\)
which contains an off-diagonal entry is zero, and
\(r(g)_{ii}=j(\Psi(g)^{w_i})\) for every \(g,i\).

\begin{proof}
Write \(\varphi=\xi\circ\iota:\mathcal B\to A\).
Both \(\ker\xi\) and \(\ker\varphi\) are prime, because
\(A\) is a domain.  The ideal \(I\) from (ACP.6) lies in
\(\ker\varphi\).  The ideal \(\lambda\) also lies there:
its coefficient generators vanish because \(O\to A\) factors
through \(k\), and its ordinary-character generators vanish
by (ACP.18).  Hence
\[
 \mathfrak m_{\mathcal B}\subset\sqrt{I+\lambda}
                                      \subset\ker\varphi.
\]
Given \(b\in\mathcal B\), choose \(a\in O\) with residue
\(\operatorname{res}_R b\).  The element \(b-a\) lies in
\(\mathfrak m_{\mathcal B}\), so
\(\varphi(b)=\varphi(a)=j(\operatorname{res}_R b)\).
This proves (ACP.19), including the original coefficient map.

A weight-zero entry monomial belongs to \(\mathcal B\) by
definition (ACP.4).  If it contains an off-diagonal entry,
its universal residue is zero, since the universal lift is
diagonal on reduction.  Equation (ACP.19) makes its image zero.
A single diagonal entry is itself a weight-zero monomial, and
its residue is \(\Psi(g)^{w_i}\), proving the final assertion.
In particular the five products in (ACP.17) vanish: the weights
around each indicated cycle telescope to zero.
\end{proof}
% END ACP9

% BEGIN ACP10
\paragraph{ACP10. Triangular test points have the same exact diagonal characters.}
\label{claim:ACP10}
Keep the data and condition (ACP.18) of ACP9.  If \(\ker\xi\)
is not off triangular, then all the matrices \(r(g)\) are
upper triangular in the given basis, or all are lower triangular.
In either case
\[
                 r(g)_{ii}=j(\Psi(g)^{w_i})
                 \quad\text{for every }g,i.
                                                        \tag{ACP.20}
\]

\begin{proof}
Failure of the off-triangular condition says that
\(\mathfrak t_+\subset\ker\xi\) or
\(\mathfrak t_-\subset\ker\xi\).  These are assertions
about all generators of one fixed ideal, so they give one
triangular direction simultaneously for every \(g\).
Diagonal entries of products of matrices with this direction
multiply.  Thus the first line of (ACP.15) holds.
The second is the residual field of \(\mathrm{IsLiftA}\).

Take \(s\in I_v\) with \(\Psi(\operatorname{loc}_v s)=1\).
ACP4 and (ACP.18) give
\[
 \prod_{i=0}^2\bigl(X-r(\operatorname{loc}_{0,v}s)_{ii}\bigr)
                                                    =(X-1)^3.
\]
Evaluate at the \(i\)-th diagonal entry.  Its difference from
one has cube zero, hence is zero in the domain \(A\).
This proves the third line of (ACP.15).  Since a DVR is
Noetherian and \(\mathrm{IsCD}\) supplies a finite residue
field, ACP7 applies and proves (ACP.20).

Every product in (ACP.17) also vanishes in either triangular
direction.  Thus this case has exactly the same polynomial
conclusion as the off-triangular case, without requiring
\(I\subset\ker\varphi\) here.
\end{proof}
% END ACP10

% BEGIN ACP11
\paragraph{ACP11. Full original finite-point control and its fixed-frame application.}
\label{claim:ACP11}
For the original coefficient DVR \(O\), every
\(F.\mathrm{Ok}(C_0)\), and every supplied
\(S:\mathrm{Space}(O,F,3)\),
\[
               \mathrm{OffFin}(S)\ \Longrightarrow
                         \mathrm{FinPt}(O,F,3).
                                                        \tag{ACP.21}
\]
Explicitly, for every original \(\mathrm{IsCD}(O,A)\), every
specified scalar tower \(O\to k\to A\), and every
\(\mathrm{IsLiftA}\) satisfying (ACP.18),
\[
 \det(X-r(g))=
   \prod_{i=0}^2\bigl(X-j(\Psi(g)^{w_i})\bigr)
                            \quad(g\in G_0).
                                                        \tag{ACP.22}
\]
There is no added characteristic-zero, coefficient-flatness,
finite-module, or finite-extension-of-fraction-fields assumption.
Neither \(\mathrm{Big}\) nor an automorphy premise is needed
for this implication.  This proves the full original
\(\texttt{Fin\allowbreak Pt\allowbreak Of}\) contract.

\begin{proof}
Fix an arbitrary test ring and point under the displayed
quantifiers.  ACP1--2 realize it on the already supplied space.
Apply ACP9 or ACP10 according as its prime kernel is off
triangular or not.  In both cases the diagonal entries are
the exact specified characters and the five products (ACP.17)
vanish.  ACP8 therefore gives (ACP.22) for every \(g\).
The choice of the test point was arbitrary, so the original
universal quantifiers have all been retained.

For the later common-frame construction, keep the frame,
link and marked \(g_1\) supplied by AFE15 or AFE16.
Once a space \(S\) on that frame has been constructed with
\(\mathrm{OffFin}(S)\), (ACP.21) supplies the precise premise
of AFE20.  That already proved transport chooses
\(U=G_L\) before every later ordinary point and returns
\(\mathrm{FinEnd}\) on the same \(g_1\).  This application
does not assert that \(S\) or \(\mathrm{OffFin}(S)\) has been
constructed merely by stating them.  It identifies the exact
consumer of the finite-point theorem just proved.
\end{proof}
% END ACP11

\subsection{One-dimensional primes that avoid every forbidden locus}

% BEGIN ACP12
\paragraph{ACP12. A one-dimensional specialization avoiding a prescribed element.}
\label{claim:ACP12}
Let \((R,\mathfrak m)\) be a Noetherian local ring,
\(Q\subset R\) a prime ideal, and \(f\in\mathfrak m\setminus Q\).
There is a prime \(P\supset Q\) such that
\[
                    f\notin P,\qquad \dim R/P=1.
                                                        \tag{ACP.23}
\]

\begin{proof}
By the ascending-chain condition choose a prime \(P\) maximal
among primes containing \(Q\) and avoiding \(f\).
The local domain \(B=R/P\) has a nonzero nonunit
\(\bar f\), and every nonzero prime of \(B\) contains
\(\bar f\): otherwise its inverse image would contradict
the maximality of \(P\).

We prove directly that this forces \(\dim B=1\).
Let \(P_1,\ldots,P_s\) be the finitely many minimal primes
over \((\bar f)\).  They have height at most one by the
principal ideal theorem and are nonzero because \(B\) is a
domain and \(\bar f\ne0\).
Given \(0\ne x\in\mathfrak m_B\), choose a prime \(Q_x\)
minimal over \((x)\) and contained in \(\mathfrak m_B\).
It is nonzero and has height at most one.  By the preceding
property it contains \(\bar f\), and hence contains some
\(P_i\).  Strict containment would give the chain
\((0)\subsetneq P_i\subsetneq Q_x\), contradicting that
height bound.  Therefore \(Q_x=P_i\), so \(x\in P_i\).
Zero belongs to the same finite union.  We have proved
\(\mathfrak m_B\subset\bigcup_iP_i\).
Finite prime avoidance gives \(\mathfrak m_B\subset P_i\)
for one index, and thus \(\mathfrak m_B=P_i\).
Its height is at most one and at least one, since it contains
the nonzero element \(\bar f\).  A local ring has dimension
equal to the height of its maximal ideal, so \(\dim B=1\).
\end{proof}
% END ACP12

% BEGIN ACP13
\paragraph{ACP13. The full finite-family nice-prime existence theorem.}
\label{claim:ACP13}
Let \((R,\mathfrak m)\) be any Noetherian local ring,
\(\varpi\in\mathfrak m\), and \(N\ge1\).
Let \((\mathfrak a_i)_{i\in E}\) be a finite family of ideals
with \(\dim R/\mathfrak a_i<N\), and let \(I\) be an ideal
with \(\dim R/I\ge N+1\).  Then there is a prime \(P\)
such that
\[
 I\subset P,\qquad \varpi\in P,\qquad \dim R/P=1,
 \qquad \mathfrak a_i\not\subset P\quad(i\in E).
                                                        \tag{ACP.24}
\]
This proves the original \(\mathrm{NicePrimeEx}\), including
arbitrary \(I\), arbitrary \(\varpi\in\mathfrak m\), and an
arbitrary finite indexing type.  The family may be empty.

\begin{proof}
First
\(\dim R/(I,\varpi)\ge N\).
Here the familiar one-equation dimension bound requires no
catenarity assumption.  In a Noetherian local ring \(T\),
lift a system of parameters of \(T/(x)\), of length
\(d=\dim T/(x)\).  Together with \(x\), these generate an
ideal with maximal-ideal radical.  The generalized principal
ideal theorem gives \(\dim T\le d+1\).
Apply this to \(T=R/I\); the quotients are nonzero since
\(I\) and \(\varpi\) lie in the maximal ideal.

Choose a minimal prime \(Q\) over \((I,\varpi)\) with
\(\dim R/Q\ge N\).  Such a prime exists because there are
finitely many minimal primes and the dimension of a Noetherian
quotient is the maximum of the dimensions of its irreducible
components.  No \(\mathfrak a_i\) is contained in \(Q\):
that containment would imply
\(\dim R/\mathfrak a_i\ge\dim R/Q\ge N\).
Choose \(a_i\in\mathfrak a_i\setminus Q\).
As \(N\ge1\), the prime \(Q\) is not maximal; choose
\(g\in\mathfrak m\setminus Q\), and put
\[
                           f=g\prod_{i\in E}a_i.
\]
It belongs to \(\mathfrak m\) and avoids the prime \(Q\).
ACP12 gives a prime \(P\supset Q\), of quotient dimension
one, avoiding \(f\).  It contains \(I\) and \(\varpi\).
If \(\mathfrak a_i\subset P\), then \(a_i\in P\) and
\(f\in P\), a contradiction.  For an empty family the
product is one and exactly the same argument applies.
\end{proof}
% END ACP13

\subsection{Propagation inside a local deformation chart}

\paragraph{The exact local propagation data.}
Let \(R\) be a Noetherian ring and use the order of ideal
inclusion on \(\operatorname{Spec}R\).  For an ideal \(I\),
write \(\operatorname{qdim}(I)=\dim(R/I)\); the zero quotient
has dimension \(-\infty\), as in the source.  Fix an integer
\(D\ge0\), subsets
\(\mathcal T,\mathcal E,\mathcal G,\mathcal N,\mathcal H\)
of \(\operatorname{Spec}R\), and an ideal \(\mathfrak s\).
They denote the type-I, exceptional-type, good, nice, and thick
primes.  A \emph{component prime} is a minimal prime over
\(\mathfrak s\).  Two such primes \(P,Q\) are adjacent if
\[
                         \dim R/(P+Q)\ge D.
                                                        \tag{ACP.25}
\]
Two nice exceptional primes \(p,q\in\mathcal N\cap\mathcal E\)
are \emph{linked} when there is a thick component prime
\(e\in\mathcal H\) contained in both.  A chain always means
a finite chain, and a chain of length zero is permitted.

The original \(\mathrm{LocData.Ok}\) consists of the following
ten requirements.  We spell them out to distinguish the
propagation theorem from its later geometric inputs.
\begin{enumerate}
\item Good primes are upward closed.
\item Type-I primes are downward closed.
\item Every nice prime is type I or exceptional.
\item Every good prime \(I\) with \(\operatorname{qdim}(I)\ge D\)
      is contained in a nice prime.
\item If \(e\) is a thick component prime and
      \(e\subset I\) with \(\operatorname{qdim}(I)\ge D\),
      then \(I\) is contained in a nice exceptional prime.
\item Every type-I minimal prime of \(R\) is a component prime.
\item Any two component primes are connected by adjacency steps
      (ACP.25).
\item Every component prime is type I or thick.
\item If a nice type-I prime is good, all primes below it are good.
\item If a nice exceptional prime \(p\) has a good type-I prime
      below it, then along every chain of links from \(p\) to
      \(q\), every type-I prime below \(q\) is good.
\end{enumerate}
The last two are the literal \(P1\) and
\(\mathrm{flipTwoThree}\) properties.  They are inputs to
this local theorem; the theorem proves how they propagate across
all component primes.

% BEGIN ACP14
\paragraph{ACP14. The full local component-propagation theorem.}
\label{claim:ACP14}
Under exactly the ten conditions above, if \(I\) is good,
type I, and \(\operatorname{qdim}(I)\ge D\), then every
type-I minimal prime of \(R\) is good.  This is the full
original \(\mathrm{LocProp}\), with no local-ring hypothesis
on \(R\) and no extra disjointness condition on the five sets.

\begin{proof}
We first record a consequence that applies at any nice prime
\(p\).  If there is a good type-I prime below \(p\), then
every type-I prime below \(p\) is good.  If \(p\) is type I,
upward closure first makes \(p\) good and condition 9 applies.
If it is exceptional, apply condition 10 to the chain of
length zero.  We will call this the propagation property at
\(p\).

We propagate the following precise information along component
primes.  A component prime \(Q\) is in one of two states:
\begin{itemize}
\item \(Q\) is type I and good;
\item \(Q\) is thick, and there is a nice exceptional prime
      \(p_0\) with a good type-I prime below it, together with
      a link chain from \(p_0\) to a nice exceptional prime
      \(p\) containing \(Q\).
\end{itemize}
In either state, if \(Q\) is also type I, it is good.
For the second state this follows from condition 10.

Here is a useful extension of either state.  Suppose
\(Q\subset M\) and \(\operatorname{qdim}(M)\ge D\).
Then there is a nice prime \(p'\supset M\) such that every
type-I prime below \(p'\) is good.  In the first state,
\(M\) is good by upward closure; condition 4 gives \(p'\),
and the propagation property at \(p'\) applies to \(Q\).
In the second state, condition 5 applied to the thick
component \(Q\) gives a nice exceptional \(p'\supset M\).
The primes \(p,p'\) are linked by \(Q\).  Append this link
to the existing chain and apply condition 10.

Now suppose \(Q,Q'\) are adjacent and the stated information
is available at \(Q\).  There is a prime \(M\supset Q+Q'\)
with \(\operatorname{qdim}(M)\ge D\): take the first prime
in a chain witnessing the dimension bound for \(R/(Q+Q')\).
The preceding extension gives a nice \(p'\supset M\) at
which every type-I prime below it is good.
If \(Q'\) is type I, it is therefore good, giving the first
state.  Otherwise condition 8 makes \(Q'\) thick.
Condition 5, applied to \(Q'\subset M\), gives a nice
exceptional prime \(q\supset M\).

If \(Q\) was in the first state, the good type-I prime
\(Q\subset q\) makes \(q\) itself a starting prime with
the required witness.  Use the length-zero chain at \(q\)
to put \(Q'\) in the second state.
If \(Q\) was in the second state, its previous ending prime
\(p\) and the new \(q\) are linked by the thick component
\(Q\), because both contain \(Q\).  Appending this link
puts \(Q'\) in the second state.  This proves the induction
step along every adjacency chain.  Notice that no assertion
that a thick prime is not type I was used.

To begin the induction, choose a minimal prime \(P_0\subset I\).
Condition 2 makes it type I, and condition 6 makes it a
component prime.  Condition 4 gives a nice \(p\supset I\).
The propagation property at \(p\), using the given good
type-I prime \(I\), makes \(P_0\) good.  Thus \(P_0\) is in
the first state.  Any type-I minimal prime \(P\) of \(R\)
is also a component prime, and condition 7 gives a finite
adjacency chain from \(P_0\) to \(P\).  The induction makes
\(P\) good, proving the assertion.
\end{proof}
% END ACP14

\subsection{From local component propagation to the whole spectrum}

\paragraph{The exact global propagation data.}
Let \(X\) be a Noetherian topological space, with subsets
\(T,G,C,Z,E\).  They are the type-I locus, good locus, core,
chosen closed-point locus, and exceptional locus, respectively.
Call an irreducible component \(Y\) relevant if \(Y\cap T\ne\varnothing\).
The original \(\mathrm{GlobData.Ok}\) asks that:
\begin{enumerate}
\item \(Z\subset C\);
\item \(C\) is preconnected;
\item every closed set meeting \(C\) meets \(Z\);
\item \(C\subset\overline{Z\setminus E}\);
\item every relevant component meets \(C\);
\item whenever two relevant components meet at a point of \(Z\),
      goodness of the whole first component implies goodness of
      the whole second;
\item every point of \(Z\setminus E\) belongs to a relevant component;
\item at least one relevant component is contained in \(G\).
\end{enumerate}
Neither openness of \(T\) nor closedness of \(G\) is assumed
in this abstract global theorem.  Those additional properties
will enter the chart-transfer theorem separately.

% BEGIN ACP15
\paragraph{ACP15. The full global component-propagation theorem.}
\label{claim:ACP15}
Under exactly the eight global conditions above, every relevant
component is contained in \(G\); in particular \(T\subset G\).
This is the full original \(\mathrm{GlobProp}\).

\begin{proof}
A Noetherian topological space has finitely many irreducible
components.  Let \(G_0\) be the union of its relevant
components that are contained in \(G\), and let \(B_0\) be
the union of the remaining relevant components.  Both are
closed, being finite unions of closed irreducible components.
Condition 7 places \(Z\setminus E\) in \(G_0\cup B_0\),
so condition 4 places the whole core \(C\) in this closed
union.  Condition 8, together with condition 5, shows that
\(C\cap G_0\ne\varnothing\).

Suppose a bad relevant component existed.  By condition 5
it would also give \(C\cap B_0\ne\varnothing\).
Preconnectedness of \(C\), applied to the two closed sets
covering it, gives a point of \(C\cap G_0\cap B_0\).
Condition 3, applied to the closed set \(G_0\cap B_0\),
therefore gives a point \(z\in Z\cap G_0\cap B_0\).
Choose a good relevant component \(Y\) and a bad relevant
component \(Y'\) containing \(z\).  Condition 6 makes
\(Y'\subset G\), contradicting its definition.
Thus no bad relevant component exists.

Every point of \(T\) belongs to an irreducible component,
which is automatically relevant.  The proved component
assertion therefore gives \(T\subset G\).
\end{proof}
% END ACP15

% BEGIN ACP16
\paragraph{ACP16. Flat local maps give the full required component chart.}
\label{claim:ACP16}
Let \(j:\operatorname{Spec}C\to X\) be an open embedding
into an arbitrary topological space.  Let \(R\) be a
commutative local \(C\)-algebra, flat over \(C\), and put
\[
 f=j\circ\operatorname{Spec}(C\to R),
 \qquad x=f(\mathfrak m_R).
\]
Then \(f\) is continuous, sends the closed point to \(x\),
and has the following two properties:
\begin{enumerate}
\item for every minimal prime \(P\) of \(R\),
      \(\overline{\{f(P)\}}\) is an irreducible component of \(X\);
\item every irreducible component \(Y\) of \(X\) through \(x\)
      is \(\overline{\{f(P)\}}\) for some minimal prime \(P\) of \(R\).
\end{enumerate}
This is the full original \(\mathrm{ChartOfFlat}\).  Neither
\(C\), \(R\), nor \(X\) is assumed Noetherian in this claim.

\begin{proof}
We recall explicitly the going-down consequence of flatness.
If \(P\in\operatorname{Spec}R\), with contraction \(p\),
and \(q\subset p\), there is \(P'\subset P\) contracting to
\(q\).  Indeed, \(C_p\to R_P\) is flat and local, hence
faithfully flat.  To see faithfulness, for any proper ideal
\(J\subset C_p\) its extension is contained in the maximal
ideal of \(R_P\).  A nonzero module contains a nonzero cyclic
submodule; flatness preserves that injection, and the cyclic
submodule remains nonzero after tensoring.  The fiber over
\(qC_p\) is consequently nonzero.  A prime of that fiber
provides the required \(P'\).

We also use the following elementary fact about the open chart.
If an irreducible set \(Y\subset X\) contains \(j(p)\), then
there is \(q\subset p\) in \(\operatorname{Spec}C\) such that
\(Y\subset\overline{\{j(q)\}}\).
The nonempty intersection \(Y\cap j(\operatorname{Spec}C)\)
is dense in \(Y\).  Its inverse image is irreducible; its
closure in \(\operatorname{Spec}C\) is \(V(q)\) for a prime
\(q\).  It contains \(p\), so \(q\subset p\).
Continuity of \(j\) and the stated density give the assertion.

Now let \(P\) be a minimal prime of \(R\), and suppose an
irreducible set \(Y\) contains \(\overline{\{f(P)\}}\).
Apply the preceding fact to obtain \(q\subset P\cap C\)
with \(Y\subset\overline{\{j(q)\}}\).  Going down gives
\(P'\subset P\) contracting to \(q\).  Minimality forces
\(P'=P\), so \(Y\subset\overline{\{f(P)\}}\).
The latter closed irreducible set is therefore maximal among
irreducible sets, proving the first property.

Conversely let \(Y\) be an irreducible component through \(x\).
The open-chart fact gives a prime
\(q\subset\mathfrak m_R\cap C\) with
\(Y\subset\overline{\{j(q)\}}\).  Going down from
\(\mathfrak m_R\) gives a prime \(P'\) over \(q\).
Choose a minimal prime \(P\subset P'\).
Specialization then gives
\(Y\subset\overline{\{f(P)\}}\).
Both maximality of \(Y\) and irreducibility of the latter
closure force equality.  Continuity and the closed-point
identity follow directly from the definition of \(f\).
\end{proof}
% END ACP16

% BEGIN ACP17
\paragraph{ACP17. The full transfer from a local chart to global components.}
\label{claim:ACP17}
Let \(T\subset X\) be open and \(G\subset X\) closed.
Let \(R\) be Noetherian and local and let
\(f:\operatorname{Spec}R\to X\) be any chart with all four
properties in ACP16, centered at \(x\).  Give \(R\) local
propagation data with type-I set \(f^{-1}(T)\), good set
\(f^{-1}(G)\), and threshold \(D\), and suppose the conclusion
of ACP14 holds for those data.  Then all three assertions hold:
\begin{enumerate}
\item If \(\operatorname{qdim}(P)\ge D\) for every minimal
      prime \(P\) of \(R\), goodness propagates between any
      two relevant components of \(X\) through \(x\).
\item If every minimal prime of \(R\) is type I, there is
      a relevant component of \(X\) through \(x\).
\item If there is a good type-I prime \(I\) of \(R\) with
      \(\operatorname{qdim}(I)\ge D\), there is a relevant
      component of \(X\) wholly contained in \(G\).
\end{enumerate}
Here relevant means meeting \(T\).  This is the full original
\(\mathrm{GlobOfCharts}\), for an arbitrary supplied chart;
it is not restricted to the flat charts just constructed.

\begin{proof}
An open set meeting \(\overline{\{y\}}\) contains \(y\).
Also, if \(y\in G\), closedness gives
\(\overline{\{y\}}\subset G\).
These two elementary observations account for both type and
goodness when passing between generic points and components.

For assertion 1, choose minimal primes \(P,P'\) with
\(Y=\overline{\{f(P)\}}\) and
\(Y'=\overline{\{f(P')\}}\), using the chart property for
components through \(x\).  Relevance and openness of \(T\)
make both primes type I.  If \(Y\subset G\), then \(P\)
is good, and its assumed dimension bound allows the local
propagation conclusion to be applied with seed \(P\).
It makes \(P'\) good.  Closedness of \(G\) then gives
\(Y'\subset G\).

For assertion 2, take an irreducible component through \(x\)
and use the chart property to write it as
\(\overline{\{f(P)\}}\) for a minimal \(P\).
By hypothesis \(f(P)\in T\), so this component is relevant.

For assertion 3, choose a minimal prime \(P\subset I\).
The inverse image of \(T\) is open, so it is stable under
generization; hence \(P\) is type I.  The local conclusion,
applied to the given seed \(I\), makes \(P\) good.
The chart's minimal-prime property makes
\(\overline{\{f(P)\}}\) an irreducible component.
It meets \(T\) at \(f(P)\) and is contained in \(G\),
proving the final assertion.
\end{proof}
% END ACP17

% BEGIN ACP18
\paragraph{ACP18. Combining the two rank-three loci gives one finite ideal.}
\label{claim:ACP18}
Let \(\mathcal B\) be any commutative ring with ideals
\(\mathfrak m,\lambda,M,I\).  If
\[
 \mathfrak m\subset\sqrt{M+\lambda},\qquad
 \mathfrak m\subset\sqrt{I+\lambda},
\]
then
\[
                  \mathfrak m\subset\sqrt{(M\cap I)+\lambda}.
                                                        \tag{ACP.26}
\]
In an original rank-three space let
\(\mathfrak e\subset\mathcal A\) be generated by the entries
whose weights are odd.  A type-I prime is off triangular and
does not contain \(\mathfrak e\); an exceptional-type prime
is off triangular and contains \(\mathfrak e\).
If \(M\) controls all type-I primes, \(I\) controls all
exceptional-type primes, and both satisfy \(\mathrm{FinLam}\),
then the single ideal \(M\cap I\) satisfies the full
\(\mathrm{OffFin}\) assertion (ACP.6).

\begin{proof}
Let \(P\) be any prime containing \((M\cap I)+\lambda\).
Since \(MI\subset M\cap I\), primality implies
\(M\subset P\) or \(I\subset P\).  Together with
\(\lambda\subset P\), the corresponding hypothesis forces
\(\mathfrak m\subset P\).  Intersect over all such primes
to obtain (ACP.26).

Every off-triangular prime of \(\mathcal A\) lies in exactly
one of the two indicated cases, according as it contains
\(\mathfrak e\).  In the first case its contraction contains
\(M\), in the second it contains \(I\); either contraction
therefore contains \(M\cap I\).  Equation (ACP.26) supplies
the remaining \(\mathrm{FinLam}\) condition on this same
ideal.  This is the actual final ideal operation in the
rank-three propagation theorem.  It does not by itself
construct the exceptional-type ideal \(I\).
\end{proof}
% END ACP18

\subsection{The exact arithmetic interface and the remaining geometric suppliers}

% BEGIN ACP19
\paragraph{ACP19. Forgetting the extra Alt properties preserves the common marking.}
\label{claim:ACP19}
At the original quantifiers,
\[
 \mathrm{GdEndAlt}\Longrightarrow\mathrm{GdEnd},
 \qquad \mathrm{HostEndAlt}\Longrightarrow\mathrm{HostEnd}.
                                                        \tag{ACP.27}
\]
These implications retain the identical \(g_1\), Marks and
FinEnd.  They do not supply either Alt endpoint.

\begin{proof}
The original rank-three automorphy predicate is
\(\mathrm{AutCsd}_0(\iota,r)\), the existence of the specified
regular algebraic conjugate-self-dual automorphic representation
attached to \(r\).  Its Alt version is literally
\[
 \mathrm{AutCsd}_0(\iota,r)\ \land
             \bigl(\mathrm{Continuous}(r)\Rightarrow
                                      \mathrm{WitA}(\iota,r)\bigr).
                                                        \tag{ACP.28}
\]
The premises of the original and Alt rank-three end conditions
are identical.  Projecting the first conjunct at each allowed
ordinary point gives the original rank-three end condition.
The Alt rank-two end condition likewise has the identical
premises and concludes the original attached-newform assertion
together with \(\mathrm{WitAQ}(\iota,W)\).  Project its first
conjunct at each allowed point and corner representation.

Now fix any initial datum under the original GdEnd quantifiers.
Use its Alt witness \(g_1\), retain Marks and FinEnd, and use
the two projections just described.  This proves GdEnd without
choosing another marker.  The Host statement is the same
argument with its rank-three conclusion only.  Nothing in
this projection constructs the extra \(\mathrm{WitA}\) or
\(\mathrm{WitAQ}\), and neither converse is asserted.
The actual terminal route continues to retain those properties.
\end{proof}
% END ACP19

\paragraph{Where the proved propagation principles are used.}
For the arithmetic spaces, the global topological space is
\(X=\operatorname{Spec}\mathcal A\), the good locus is
\(V(M\mathcal A)\), and the core is
\[
 \{\mathfrak p:\mathfrak m_{\mathcal B}
        \subset\pi(\mathfrak p),\quad
                      \mathfrak t_+\not\subset\mathfrak p\}.
                                                        \tag{ACP.29}
\]
The chosen points are its maximal ideals.  In rank two the
exceptional locus is empty and the local threshold is
\(2d+1\); in rank three the exceptional locus is
\(V(\mathfrak e)\) and the threshold is \(3d+1\).
The original \(P1'\) and \(P23'\) conditions supply the
two local propagation rules in ACP14.  ACP13 supplies the
prime-avoidance mechanism once the required bounds on the
forbidden loci have been proved.  ACP16--17 transfer that
local result to the component-incidence property used by
ACP15.  ACP18 is the final combination of the type-I and
exceptional-type ideals.  ACP11 then supplies the exact
finite-point conclusion used by AFE20.

The still needed geometric inputs are specific.  They are:
the representable local charts and their flatness; the sharp
component-chain and dimension estimates; the bounds for the
triangular and coincident-character loci; connectedness and
the indicated closed-point density of the core; the passage
from the seed prime to a sufficiently large local prime; and
the actual rank-two-to-rank-three corner construction producing
the exceptional-type ideal.  They are the inputs proved in
the source files \texttt{Loc\allowbreak Flat}, \texttt{Chained\allowbreak Bot},
\texttt{Chained\allowbreak Of\allowbreak Comap},
\texttt{Krull\allowbreak Dim\allowbreak Lt\allowbreak Of\allowbreak Not\allowbreak Le\allowbreak Triangular},
\texttt{Preconnected\allowbreak Rank\allowbreak Two}, \texttt{Seedpt}, and
\texttt{Quotient\allowbreak Hom\allowbreak Nonvanishing\allowbreak Bound}, respectively.
This chapter proves neither the full rank-two nor the full
rank-three propagation theorem without those constructions.

The actual later call keeps a single extended frame for both
ranks, while allowing different spaces and modularity ideals.
The rank-two propagation result feeds the rank-three corner
and propagation arguments.  Their \(\mathrm{OffFin}\) result
then enters ACP11.  The automorphic seeds, the modularity
realizations and arithmetic-point classicality remain separate
suppliers.  The completed Dat and Host frame entries, including
Host \(\ell=3\), are reused; no new restriction on their
marked element, specified place, or order of choices has
been introduced.



\section{One character and compatible measures on the actual setting}
\label{sec:global-measure-rescaling-v53}

The purpose of this chapter is the original \texttt{Rescale\allowbreak Step3}.
We start with the supplied Hida setting, its supplied adelic bridge, and
its already normalized measures at infinity.  We construct the fiber
measures which determine the global stabilizer measures and the local
flag-stabilizer measures.  We then use AMD13--AMD17 to fill in the finite
factors without changing those measures at infinity.  The character is
chosen once, before any word, finite place, or auxiliary tuple.

The construction of a fiber measure is different from the assertion
that the volume of an arithmetic quotient equals a Siegel--Weil
constant.  Here the normalization is fixed by local Fourier inversion
and local disintegration.  At almost every finite place we compute its
mass exactly.  The resulting positive, convergent product constructs the
global fiber measure and supplies its nonvanishing.  Consequently the
measure construction below does not require a theta-integral identity.
We state the original source endpoint with its exact
\texttt{Siegel\allowbreak Weil\allowbreak Unitary\allowbreak Convergent} premise; that larger premise, and the
later arithmetic volume comparisons which use it, are not proved in this
chapter.

\subsection{Notation and the precise measure contracts}

Let $E/F$ be a CM extension, with involution $c$, and let
$h(x,y)=x^*Jy$ be a nondegenerate Hermitian form on $E^n$.
For a place $v$ put $k_v=F_v$, $A_v=E\otimes_F k_v$ and
$G_v=U(J)(k_v)$.  The algebra $A_v$ is either a quadratic field or
$k_v\times k_v$; at infinity it is $\mathbb C$ over $\mathbb R$.
All assertions below include the split algebra.  We use column notation
for an ordered tuple, so
\[
 \begin{gathered}
 V_{v,a}=\operatorname{Mat}_{n\times a}(A_v),\qquad
 \gamma_v(X)=X^*JX,\\
 U_{v,a}=\{X:\text{the columns of }X\text{ are }A_v\text{-independent}\}.
 \end{gathered}
 \tag{GMR.1}
\]
For the product of the infinite places the notation has the same
meaning, component by component.  On a split algebra, independence
means independence in both factors.  Write $\mathcal H_{v,a}$ for the
$k_v$-space of Hermitian $a\times a$ matrices.  Its dimension is $a^2$;
$V_{v,a}$ has dimension $2na$ over $k_v$.

An adelic character $\psi$ satisfying \texttt{Is\allowbreak Char\allowbreak AF} is continuous,
unitary, trivial on the diagonal copy of $F$, and nontrivial.  The
normalization \texttt{Is\allowbreak Std\allowbreak Char} is exactly
\[
 \psi_\infty(x)=\exp\left(2\pi i\sum_{\tau:F\hookrightarrow\mathbb R}x_\tau\right).
 \tag{GMR.2}
\]
The pairings defining the source measures \texttt{Sd\allowbreak W} and
\texttt{Sd\allowbreak Herm} are the trace pairings \texttt{BR} and
\texttt{pair\allowbreak Herm}, with this character.  We always use those pairings;
an arbitrary normalization of coordinate Lebesgue or additive Haar
measure is not substituted for a self-dual measure.

Here is the full fiber-family condition used below.  Given self-dual
measures $dX$ on $V_{v,a}$ and $dT$ on $\mathcal H_{v,a}$, a family
$\omega_{v,T}$ satisfies \texttt{Is\allowbreak Fibre\allowbreak Family} when
\[
\begin{split}
 &\omega_{v,T}\bigl(V_{v,a}\setminus(\gamma_v^{-1}(T)\cap U_{v,a})\bigr)=0,\\
 &\int_{U_{v,a}} f(X)\,dX
       =\int_{\mathcal H_{v,a}}\int f(X)\,d\omega_{v,T}(X)\,dT
       \quad(f\geq0\text{ measurable}),
\end{split}
 \tag{GMR.3}
\]
and, for every real continuous compactly supported $f$ whose support is
contained in $U_{v,a}$, the fiber integral is finite and continuous in
\emph{every} $T\in\mathcal H_{v,a}$.  There is no exclusion of singular
Hermitian matrices $T$ in this condition.  Measures are understood on
the ambient tuple space, by extension by zero from the independent
locus.

For a rational nondegenerate Hermitian matrix $T$, the original
\texttt{Is\allowbreak Om\allowbreak T} requires a measure $\nu_T$ supported on the adelic level
set $\gamma=T$ and the following assertion for \emph{all} self-dual
measures and \emph{all} fiber families satisfying (GMR.3).
If
\[
 \Phi(X)=\Phi_\infty(X_\infty)\Phi_f(X_f),\qquad
 \Phi_f(X_f)=\prod_v\Phi_v(X_v),
 \tag{GMR.4}
\]
where $\Phi_\infty$ is continuous with compact support, every $\Phi_v$
is locally constant with compact support, and $\Phi_v$ is the indicator
of the coordinate lattice attached to a single rational $F$-basis of
$E$ at all but finitely many places, then $\Phi$ is integrable,
the product below is multipliable, and
\[
 \int\Phi\,d\nu_T
   =\left(\int\Phi_\infty\,d\omega_{\infty,T}\right)
      \prod_v\int\Phi_v\,d\omega_{v,T}.
 \tag{GMR.5}
\]
This includes the existence quantifiers in the source's
\texttt{All\allowbreak Some}: the self-dual measures and the fiber families must
exist, and the conclusion must hold for every permitted choice.

For an actual tuple $X_0$ and its actual stabilizer $H$, the condition
\texttt{Is\allowbreak Compat\allowbreak T}, or \texttt{Is\allowbreak Compat\allowbreak Loc} at one place, says that
$\nu_T$ is supported on the actual orbit and
\[
 \int_G f(g)\,d\mu_G(g)
  =\int_{G X_0}\int_H f(s(y)h)\,d\mu_H(h)\,d\nu_T(y)
 \tag{GMR.6}
\]
for every measurable $f:G\to[0,\infty]$ and \emph{every} section $s$
with $s(y)X_0=y$ on that orbit.  The original statement does not assume
that this section is measurable.  We prove precisely this version by
showing that the inner integral is independent of the section and has
a measurable representative defined by local sections.

The principal source interfaces are
\texttt{Vol\allowbreak From\allowbreak SW.\allowbreak Is\allowbreak Om\allowbreak T}, \texttt{Vol\allowbreak From\allowbreak SW.\allowbreak Is\allowbreak Compat\allowbreak T},
\texttt{Loc\allowbreak Vol.\allowbreak Is\allowbreak Om}, \texttt{Loc\allowbreak Vol.\allowbreak Loc\allowbreak Tie\allowbreak W}, and
\texttt{Setting\allowbreak Exists\allowbreak Adapted.\allowbreak Rescale\allowbreak Step3}.  Their literal definitions
and the producer bodies are included in the source manifest.  AMD1,
AMD4--AMD8 and AMD13--AMD17 are used at their reviewed original scopes.

% BEGIN GMR1
\paragraph{GMR1. A single standard character and its local conductors.}
\label{claim:GMR1}
For every totally real number field $F$ there is a unique character
$\psi^0$ satisfying \texttt{Is\allowbreak Char\allowbreak AF} and (GMR.2).
Every character satisfying \texttt{Is\allowbreak Char\allowbreak AF} has the form
$\psi(x)=\psi^0(bx)$ for a unique $b\in F^\times$.
In particular its restriction to each completion is nontrivial, and
its finite-place conductor is $\mathcal O_v$ at all but finitely many
places.  These conclusions are made before any tuple or word is chosen.

\begin{proof}
On $\mathbb A_{\mathbb Q}$ take the character whose real factor is
$e^{2\pi ix}$ and whose $p$-adic factor is
$e^{-2\pi i\{x\}_p}$, where $\{x\}_p$ is the finite negative-power
part of the $p$-adic expansion, viewed in $\mathbb Q/\mathbb Z$.
The sum of the negative-power parts of a rational number differs from
that number by an integer.  Thus the product of these factors is one
on $\mathbb Q$.  Compose with the adelic trace from $F$ to $\mathbb Q$.
The local traces over the places above $p$ add to the rational trace;
the infinite trace is the sum over the real embeddings.  We obtain
(GMR.2), triviality on $F$, continuity, and absolute value one.
The infinite factor is nontrivial, so the whole character is nontrivial.
The annihilator of $\mathcal O_v$ for its local trace pairing is the
inverse different.  Outside the finitely many ramified rational primes
this is $\mathcal O_v$ itself.

We recall the character argument which also verifies that no local
factor can disappear for another permitted $\psi$.  A continuous
character of a real vector space is the exponential of a real linear
functional.  A continuous character of a nonarchimedean local field is
of the form $x\mapsto\psi^0_v(b_vx)$: on a compact open additive subgroup
it factors through a finite quotient, whose character group is its
trace dual, and the compatible characters on the increasing groups
$\varpi^{-m}\mathcal O_v$ give the unique parameter $b_v$.
Continuity on the finite adeles implies triviality on
$\prod_{v\notin S}\mathcal O_v$ for some finite $S$.  One way to see
this is to choose a small arc about one containing no nontrivial
subgroup of the unit circle and use continuity on that compact product.
Consequently the parameters $b_v$ form an adele $b$.

The annihilator of diagonal $F$ for $\psi^0$ is diagonal $F$ itself.
For completeness, choose $r\in F$ with $b_f-r\in\prod_v\mathcal O_v$;
this is the Chinese remainder theorem applied to the finitely many
denominators of $b_f$.  Triviality on $\mathcal O_F$ now says that
$b_\infty-r$ lies in the trace-dual lattice of $\mathcal O_F$ in
$F_\infty$.  That lattice is the image of the inverse different, so
$b_\infty-r$ is the infinite image of an element $d\in F$.
Subtract $r+d$.  The remaining parameter has zero infinite component,
and its finite character is trivial on $F$.  The Chinese remainder
theorem also proves that $F$ is dense in $\mathbb A_{F,f}$, so this finite
character is identically one.  Local nondegeneracy of the trace pairing
then makes every remaining finite parameter zero.  This proves the
annihilator assertion and hence $b\in F$.  Nontriviality of $\psi$ gives
$b\ne0$.  At almost every place $b$ is a unit, so the conductor assertion
follows from the assertion for $\psi^0$.

Finally, if the infinite factor is (GMR.2), the infinite images of $b$
are all one, so $b=1$.  This proves uniqueness of the standard
character, as well as the original \texttt{standard\_\allowbreak character\_\allowbreak exists}
assertion.
\end{proof}
% END GMR1

% BEGIN GMR2
\paragraph{GMR2. Self-dual measures exist and are uniquely normalized.}
\label{claim:GMR2}
For every nontrivial local factor supplied by GMR1, the source pairings
on $V_{v,a}$, on $\mathcal H_{v,a}$, and on $A_v^I$ with pairing
$\psi_v(\sum_{i\in I}\operatorname{Tr}_{A_v/k_v}(x_i y_i))$
have unique self-dual additive Haar measures.  The same holds for the
product of the infinite places.  These measures satisfy the literal
source Fourier-inversion predicate, not just invariance under addition.
Whenever an integral lattice is its own annihilator, its self-dual mass
is one.

\begin{proof}
Each pairing is nondegenerate over the base field.  For tuple space
this follows from invertibility of $J$ and nondegeneracy of the trace
pairing of the quadratic etale algebra.  For Hermitian matrices,
diagonal entries give the scalar trace pairing and each pair of
off-diagonal entries gives the trace pairing of $A_v$; the coefficients
which occur are nonzero in characteristic zero.  This also proves the
assertion for a split algebra.  The direct sum pairing on $A_v^I$ is
nondegenerate, including $I=\varnothing$.

Here is the normalization argument.  In real coordinates, Fourier
inversion for $e^{2\pi i x\cdot y}$ uses ordinary Lebesgue measure.
Replacing this pairing by $e^{2\pi i x^tBy}$ multiplies the self-dual
measure by $|\det B|^{1/2}$.  Indeed a change of variable in the two
successive transforms contributes $|\det B|^{-1}$, exactly cancelling
the square of that factor.  The same calculation applies over a local
field.  It can also be obtained directly from a compact open lattice
$L$: the transform of $1_L$ is
\[
 \widehat{1_L}(y)=\mu(L)1_{L^\perp}(y),\qquad
 \widehat{\widehat{1_L}}(z)=\mu(L)\mu(L^\perp)1_{-L}(z).
 \tag{GMR.7}
\]
Starting with any additive Haar measure, there is exactly one positive
rescaling for which $\mu(L)\mu(L^\perp)=1$.  Fourier inversion follows
first on locally constant compactly supported functions by finite
Fourier inversion on lattice quotients, then on continuous compactly
supported functions whenever the transform is integrable.  In the real
case the corresponding extension follows by convolution with Gaussian
approximate identities and dominated convergence for the integrable
transform.  This is precisely the conditional inversion assertion in
the source definition of \texttt{Self\allowbreak Dual}.

To check uniqueness even for that conditional assertion, two additive
Haar measures differ by a positive factor $c$.  Apply the asserted
inversion to a nonzero compactly supported locally constant function
at a finite place, or to a smooth compactly supported function at
infinity.  Its Fourier transform is integrable; two successive
transforms for $c\mu$ multiply the answer for $\mu$ by $c^2$.
Thus $c^2=1$, and positivity gives $c=1$.  Formula (GMR.7) now gives
$\mu(L)=1$ when $L=L^\perp$.  Products of the individual infinite-place
measures give the asserted measure at infinity by Fubini and the same
uniqueness argument.
\end{proof}
% END GMR2

% BEGIN GMR3
\paragraph{GMR3. The Gram map is a submersion on independent tuples.}
\label{claim:GMR3}
Over every $k_v$ and at the product of the infinite places, the map
$\gamma:U_{v,a}\to\mathcal H_{v,a}$ is an analytic submersion.
This remains true when $\gamma(X)$ is singular.  Near every independent
$X$, the fiber has dimension $2na-a^2$ over the base field.

\begin{proof}
Since the columns of $X$ are independent and $J$ is invertible, the
map $X^*J:A_v^n\to A_v^a$ is onto.  Choose a right inverse $W$, so
$X^*JW=1$.  For a Hermitian matrix $S$, put $Y=WS/2$.
The differential of the Gram map is
\[
 D\gamma_X(Y)=X^*JY+Y^*JX=S/2+S^*/2=S.
 \tag{GMR.8}
\]
All divisions are in a characteristic-zero field; the argument also
works at places above two.  Over a split algebra the right inverse
is chosen separately in each factor.  Over a product of infinite
completions it is chosen componentwise.  Thus the differential is
surjective as a map of base-field vector spaces.

Choose $a^2$ of the real or local-field coordinates on which this
differential is invertible and retain the remaining coordinates $z$.
The inverse function theorem applied to $X\mapsto(z(X),\gamma(X))$
gives local analytic coordinates $(z,T)$ in which $\gamma$ is projection
onto $T$.  There are $2na-a^2$ coordinates $z$.
Nothing in this argument used $\det\gamma(X)\ne0$; independence is the
necessary hypothesis.
\end{proof}
% END GMR3

% BEGIN GMR4
\paragraph{GMR4. Construction of the full continuous fiber family.}
\label{claim:GMR4}
With the self-dual measures of GMR2 there is a family
$\omega_{v,T}$ satisfying all of (GMR.3), with continuity on the whole
Hermitian space.  On the independent part of every fiber it is a Radon
measure, positive on every nonempty relatively open set.  It is zero
when that independent fiber is empty.  The same assertion holds at
infinity.  More generally, the assertion holds for any prescribed
pair of additive Haar measures on tuple space and Hermitian space;
self-duality is not needed for this local construction.
In particular the source's chosen \texttt{d\allowbreak X}, \texttt{d\allowbreak T}
and \texttt{om\allowbreak Fam} satisfy their defining alternatives rather than
their zero fallback alternatives.

\begin{proof}
Use a chart from GMR3 and write its inverse as $X=\chi(z,T)$.
Choose coordinate Haar measures $dz$ and $dT$ and express the ambient
measure in this chart as
\[
 dX=j(z,T)\,dz\,dT,\qquad j(z,T)>0.
 \tag{GMR.9}
\]
Here $j$ is the absolute Jacobian, including the positive constants
relating the coordinate measures to the prescribed additive Haar
measures.  It is continuous; over a nonarchimedean field the absolute
value of a nonvanishing analytic determinant is locally constant.
On the part of the fiber in this chart define its measure by
$j(z,T)\,dz$.  On overlapping charts the coordinate change fixes $T$.
Its derivative is block triangular, so its full Jacobian is its
vertical Jacobian.  The change-of-variables formula therefore says
that the two fiber measures agree on the overlap, for each $T$.

Take a countable locally finite chart cover of $U_{v,a}$ and glue these
agreeing measures.  This is an ordinary gluing of measures on open
sets: one can disjointify a countable cover and sum the restrictions,
and agreement on overlaps makes the result independent of the
disjointification.  Extend each resulting measure by zero to the
ambient tuple space.  The support assertion in (GMR.3) is immediate.
Fubini in each chart, followed by a countable disjointification and
monotone convergence, proves the disintegration identity for every
nonnegative measurable function.  This construction also shows that
the fiber integrals occurring there are measurable in $T$.

Let $f$ be continuous with compact support contained in $U_{v,a}$.
Cover its support by finitely many such charts and take a continuous
partition of unity subordinate to slightly smaller charts.  In a
nonarchimedean space a finite clopen refinement gives the same
partition.  Each resulting fiber integral is an integral of a
continuous function with a common compact vertical support, locally
in $T$.  Its absolute value is bounded by a fixed integrable function
because $j$ is bounded on that compact set.  Dominated convergence
proves continuity and finiteness.  The summands extend by zero past
the boundaries of their smaller charts, so this proves continuity on
all of $\mathcal H_{v,a}$, including parameters with empty fibers.
The construction is Radon on the fiber inside $U_{v,a}$ and its chart
density is strictly positive.  These observations prove the remaining
assertions.  If a singular fiber approaches a dependent tuple, no
unproved assertion of ambient local finiteness is being made: the
Radon assertion here is on the independent fiber, exactly where the
quotient argument below takes place.
\end{proof}
% END GMR4

% BEGIN GMR5
\paragraph{GMR5. Pointwise uniqueness and covariance of the fibers.}
\label{claim:GMR5}
For any fixed pair of additive Haar measures, the family in GMR4 is
the unique family satisfying the original \texttt{Is\allowbreak Fibre\allowbreak Family} condition.
In particular uniqueness holds at each singular Hermitian parameter,
not merely almost everywhere.  Every unitary left action preserves
$\omega_{v,T}$.  If $u$ is upper unitriangular, right multiplication
$X\mapsto Xu$ carries $\omega_{v,T}$ to $\omega_{v,u^*Tu}$.
The same assertions hold at infinity.

\begin{proof}
For a real continuous $f$ of compact support in the independent locus,
two proposed families have fiber integrals with equal integrals over
every measurable set of Hermitian parameters, by disintegration.
The functions are consequently equal almost everywhere.  Both are
continuous by the defining condition, and additive Haar measure is
positive on every nonempty open set.  They are therefore equal
everywhere.  For each fixed $T$, these test functions determine the
measure on the independent fiber: compact subsets have continuous
cutoffs with compact support in the independent locus, and the
defining integrability gives local finiteness there.  The measures are
Radon on that locally compact space, and a compact exhaustion followed
by monotone convergence determines their values on all Borel sets.
Both vanish outside that locus.  Hence the entire ambient measures
are equal at every $T$.

More generally, suppose invertible linear maps $A$ on tuple space and
$C$ on Hermitian space satisfy $\gamma A=C\gamma$ and preserve the
independent locus.  Pushing (GMR.3) through these changes of variables
and using uniqueness gives
\[
 A_*\omega_T=\frac{|\det C|}{|\det A|}\,\omega_{CT}.
 \tag{GMR.10}
\]
The determinant is over the base local field, with the product of
absolute values at infinity.  For $A(X)=gX$, $g\in U(J)$, its determinant
is $\operatorname{Nm}_{A_v/k_v}(\det g)^a=1$ and $C=1$.
For $A(X)=Xu$ with $u$ upper unitriangular, both $A$ and
$C(T)=u^*Tu$ are unipotent linear maps and have determinant one.
Thus (GMR.10) gives both claimed covariance assertions.  The existence
and uniqueness in GMR2 also show that these conclusions apply to
\emph{any} source witnesses for the self-dual measures and fiber
families, including the witnesses chosen by \texttt{om\allowbreak Fam}.
\end{proof}
% END GMR5

% BEGIN GMR6
\paragraph{GMR6. Exact counts of nondegenerate Hermitian frames.}
\label{claim:GMR6}
Let $k=\mathbb F_q$ have odd characteristic, let $A$ be its quadratic
etale algebra with its nontrivial involution, and let $J,T$ be
nondegenerate Hermitian forms of ranks $n,a$, where $0\leq a\leq n$.
Put $d=2na-a^2$, and put $\varepsilon=1$ in the split case and
$\varepsilon=-1$ in the field case.  Then
\[
 \#\{X\in\operatorname{Mat}_{n\times a}(A):X^*JX=T\}
  =q^d\prod_{j=n-a+1}^{n}(1-\varepsilon^j q^{-j}).
 \tag{GMR.11}
\]
This gives, in particular, the full original finite-field count bound
\texttt{FFCount} when $a+2\leq n$.

\begin{proof}
Over $\mathbb F_{q^2}$ every nondegenerate Hermitian form is equivalent
to the identity: diagonalize by orthogonal elimination and rescale
each nonzero diagonal entry, using the surjectivity of the norm to
$\mathbb F_q^\times$.  The unitary group acts transitively on the
ordered frames with Gram matrix $T$, since an isometry on their
nondegenerate spans extends by choosing orthonormal bases of their
orthogonal complements.  The stabilizer fixes the span pointwise and
is $U_{n-a}(q)$.

Here is the order calculation.  For a nontrivial additive character
$\eta$ of $\mathbb F_q$ and $s\ne0$,
\[
 \sum_{z\in\mathbb F_{q^2}}\eta(s z\bar z)
   =1+(q+1)\sum_{t\in\mathbb F_q^\times}\eta(st)=-q.
\]
Character orthogonality therefore gives
\[
 \#\{z\in(\mathbb F_{q^2})^m:\textstyle\sum z_i\bar z_i=1\}
  =q^{2m-1}-(-1)^m q^{m-1}.
\]
Choose the first unit vector of an orthonormal basis and then an
orthonormal basis of its complement.  Induction gives
\[
 |U_m(q)|=q^{m^2}\prod_{j=1}^{m}(1-(-1)^j q^{-j}).
 \tag{GMR.12}
\]
Taking $|U_n|/|U_{n-a}|$ proves (GMR.11) in the field case.

For $A=k\times k$, a Hermitian space is a perfect pairing between two
$k$-spaces.  Its unitary group is $\operatorname{GL}_n(k)$, and the
stabilizer of a nondegenerate ordered frame is
$\operatorname{GL}_{n-a}(k)$.  Counting successively independent
columns gives $|\operatorname{GL}_m(k)|=q^{m^2}\prod_{j=1}^m(1-q^{-j})$.
Its quotient gives the split formula.  Equivalently one may choose
the first-component frame in $\prod_{i=0}^{a-1}(q^n-q^i)$ ways and
then choose its paired second-component frame in $q^{a(n-a)}$ ways;
this is exactly the source's \texttt{FFSplit} formula.

To obtain the uniform bound, set $s=n-a+1$.  When $a+2\leq n$, $s\geq3$.
For $b_j=-\varepsilon^j q^{-j}$,
$\sum_{j=s}^n|b_j|\leq2q^{-s}$ and
$|\prod(1+b_j)-1|\leq\exp(\sum|b_j|)-1\leq4q^{-s}$.
Thus (GMR.11) differs from $q^d$ by at most $4q^{d-2}$.
The constant $4$, including the empty product $a=0$, supplies the
existential constant in the full original \texttt{FFCount}.
\end{proof}
% END GMR6

% BEGIN GMR7
\paragraph{GMR7. Exact good-place mass of the canonical Gram fiber.}
\label{claim:GMR7}
Fix $1\leq a$, $2a<n$, a rational tuple $X_0$ with nondegenerate
$T=X_0^*JX_0$, an adelic character satisfying \texttt{Is\allowbreak Char\allowbreak AF}, and
any rational $F$-basis used to define the source lattices.
Outside a finite set of places, the compact fiber
\[
 C_v=\{X\in\operatorname{Mat}_{n\times a}(\mathcal O_{E,v}):X^*JX=T\}
\]
is nonempty and has canonical fiber mass
\[
 a_v:=\omega_{v,T}(C_v)
       =\prod_{j=n-a+1}^{n}(1-\varepsilon_v^j q_v^{-j}),
 \tag{GMR.13}
\]
where $\varepsilon_v=1$ at split places and $-1$ at inert unramified
places.  Formula (GMR.13) holds for every permitted self-dual measure
and fiber family.

\begin{proof}
Exclude places above two, places ramified in $E/F$ or in the different
of $F$, nonunit places of the character parameter in GMR1, denominators
of $J,T,X_0$ and their nonzero determinants, and the finitely many
places where the chosen rational basis lattice is not
$\mathcal O_{E,v}$.  At a remaining place $A_v$ is unramified etale,
$J,T$ are integral and unimodular, and the character has conductor
$\mathcal O_v$.  The two trace-pairing lattices on tuple space and
Hermitian space are perfect integral duals.  Thus their self-dual
masses are one by GMR2.  The specified tuple $X_0$ belongs to $C_v$.

Reduction modulo $\varpi$ of any solution has independent columns,
because its Gram matrix $\bar T$ is invertible.  At such a solution
the differential (GMR.8) is onto over the residue field.  An arbitrary
lift of a solution modulo $\varpi^m$ has a Gram error at order
$\varpi^m$; correcting it by $\varpi^mY$ changes that error by
$D\gamma(Y)$ modulo $\varpi^{m+1}$.  Surjectivity shows that each
solution has exactly $q_v^{2na-a^2}$ lifts to the next level.
If $N_m$ is the number of solutions modulo $\varpi^m$, then
$N_m=N_1q_v^{(2na-a^2)(m-1)}$.

The ambient volume of the integral tuples satisfying
$\gamma(X)\equiv T\pmod{\varpi^m}$ is $N_mq_v^{-2nam}$.
The volume of the corresponding Hermitian ball is $q_v^{-a^2m}$.
Their quotient is consequently $N_1q_v^{-(2na-a^2)}$, independently
of $m\geq1$.  By (GMR.3) this quotient is the average of the fiber
mass of the integral tuples over that ball.  To apply the continuity
clause literally, multiply the lattice indicator by the indicator
of $\gamma(X)\equiv T\pmod\varpi$: its support is a compact subset
of the independent locus, and on all these balls it has the same
fiber integral.  Continuity from GMR4 lets $m$ tend to infinity and
identifies the average with $\omega_{v,T}(C_v)$.
Now GMR6 gives exactly (GMR.13).  GMR2 and GMR5 remove all dependence
on the chosen witnesses.  Enlarging the exceptional set handles any
other rational basis, which is precisely the basis quantifier in
\texttt{Is\allowbreak Pure\allowbreak Fin}.
\end{proof}
% END GMR7

% BEGIN GMR8
\paragraph{GMR8. The density product is convergent and strictly positive.}
\label{claim:GMR8}
In the range of GMR7, put $s=n-a+1$.  For every good place,
\[
 a_v>0,\qquad |a_v-1|\leq4q_v^{-s},\qquad
 |\log a_v|\leq4q_v^{-s},\qquad s\geq a+2\geq3.
 \tag{GMR.14}
\]
Hence every product of the $a_v$ outside a fixed finite exceptional
set converges to a finite strictly positive real number.

\begin{proof}
Each factor in (GMR.13) is positive.  The bound for $a_v-1$ was proved
in GMR6.  Since $q_v^{-j}\leq1/8$ for $j\geq s\geq3$, the elementary
inequality $|\log(1+t)|\leq2|t|$ for $|t|\leq1/2$ gives
\[
 |\log a_v|\leq2\sum_{j=s}^nq_v^{-j}\leq4q_v^{-s}.
\]
There are at most $[F:\mathbb Q]$ places above a rational prime $p$,
and each has $q_v\geq p$.  Therefore
\[
 \sum_vq_v^{-s}\leq[F:\mathbb Q]\sum_{m=2}^{\infty}m^{-s}<\infty.
\]
The sum of the real logarithms converges absolutely; its exponential
is the desired finite positive product.  This proves positivity of
the normalization, rather than inferring it from convergence of a
possibly zero Euler product.
\end{proof}
% END GMR8

% BEGIN GMR9
\paragraph{GMR9. The global fiber measure satisfies the full \texttt{Is\allowbreak Om\allowbreak T}.}
\label{claim:GMR9}
Under the original range $1\leq a$, $2a<n$ and for the actual rational
nondegenerate Gram matrix $T$ of $X_0$, there is a nonzero Radon measure
$\nu_T$ on adelic tuple space satisfying the full original
\texttt{Is\allowbreak Om\allowbreak T}, including all basis, test-function, self-dual-measure
and fiber-family quantifiers in (GMR.5).  It is positive on every
nonempty relatively open subset of the adelic Gram fiber.

\begin{proof}
Write $Y_v=\gamma_v^{-1}(T)$ and $Y_\infty=\gamma_\infty^{-1}(T)$.
Nondegeneracy of $T$ implies that every tuple in these fibers is
independent.  Thus each $Y_v$ is closed in tuple space and its measure
$\omega_{v,T}$ is Radon by GMR4.  Choose an exceptional finite set
$S_0$ as in GMR7.  For $v\notin S_0$, the set $C_v\subset Y_v$ is
compact open of positive finite mass $a_v$.  Let
$\bar\omega_v=a_v^{-1}\omega_{v,T}$, so
$\bar\omega_v(C_v)=1$.

Here is the restricted-product measure construction with its scalar
displayed.  On
\[
 Y(S)=\prod_{v\in S}Y_v\times\prod_{v\notin S}C_v,
       \qquad S_0\subset S\text{ finite},
\]
use the product of $\omega_{v,T}$ for $v\in S_0$ and
$\bar\omega_v$ for $v\notin S_0$.  The infinitely many remaining
factors are probability measures on compact metrizable spaces, so
their product is the countably additive product measure.  When $S$
is enlarged, restriction to the smaller $Y(S)$ agrees, because the
new compact factors have normalized mass one.  These measures
therefore glue on the union $Y_f=\bigcup_SY(S)$.  Multiply the result
by the single positive number
\[
 A_0=\prod_{v\notin S_0}a_v.
 \tag{GMR.15}
\]
It exists by GMR8.  Call the resulting measure $\nu_{T,f}$.
Equivalently its restriction to $Y(S)$ is
\[
 \left(\prod_{v\notin S}a_v\right)
 \left(\bigotimes_{v\in S}\omega_{v,T}\right)
 \otimes\left(\bigotimes_{v\notin S}\bar\omega_v|_{C_v}\right).
 \tag{GMR.16}
\]
This formula proves directly that changing $S_0$ has no effect.

The usual restricted-product topology makes every $Y(S)$ open.
A compact set is covered by finitely many of these opens and hence
is contained in one of them.  Its projection at every place in $S$
is compact, so (GMR.16) proves finite mass on compact sets.  The
measure is Radon, either by regularity of these countable products
of compact metric probability spaces and local Radon factors, or by
compact exhaustion of each of their countably many charts.
It is nonzero and positive on every nonempty basic open set: the
finitely many specified local open sets have positive mass by GMR4,
and the remaining compact factors have mass one before multiplication
by a strictly positive scalar.

AMD1 identifies finite-adelic tensor coordinates with their
restricted-product local coordinates.  Under this identification
$Y_f$ is exactly the finite-adelic Gram fiber.  Set
$\nu_T=\omega_{\infty,T}\otimes\nu_{T,f}$ and push it into adelic
tuple space.  It is supported on $\gamma=T$, is Radon, and is nonzero
because the rational tuple $X_0$ supplies a point in every local
fiber and hence a nonempty product neighborhood.

Now take any test data in (GMR.4), using any rational basis.  Enlarge
$S$ to include $S_0$, the exceptional basis places, and the finitely
many places where $\Phi_v$ is not the standard lattice indicator.
On the fiber, $\Phi$ vanishes outside $Y_\infty\times Y(S)$, and its
support is compact in the infinite coordinate and the finitely many
unrestricted coordinates.  Formula (GMR.16) and Fubini therefore
prove absolute integrability and
\[
 \int\Phi\,d\nu_T
 =\left(\int\Phi_\infty\,d\omega_{\infty,T}\right)
   \left(\prod_{v\in S}\int\Phi_v\,d\omega_{v,T}\right)
   \prod_{v\notin S}a_v.
 \tag{GMR.17}
\]
The infinite product on the right is convergent by GMR8.  This also
proves multipliability when a finite local integral is zero or
complex; outside $S$ the factors are still the positive real $a_v$.
It is (GMR.5).  Finally, GMR2 and GMR5 prove equality with every
self-dual measure and every fiber family appearing in the source
\texttt{All\allowbreak Some} quantifiers.  They also provide their existence.
Thus the full original contract, rather than just a factorization
for one selected set of witnesses, has been proved.
\end{proof}
% END GMR9

% BEGIN GMR10
\paragraph{GMR10. The adelic Gram fiber is the actual unitary orbit.}
\label{claim:GMR10}
For the data of GMR9, the full adelic Gram fiber equals
$U(J)(\mathbb A_F)X_0$.  The orbit map identifies the quotient by the
actual stabilizer with this fiber, has continuous local sections, and
is open.  The measure $\nu_T$ of GMR9 is invariant under the actual
adelic unitary group.

\begin{proof}
We give the needed extension argument rather than assume an orbit
identification as a measure input.  Over a quadratic field extension,
two nonisotropic vectors $u,v$ of the same norm $t$ can be matched by
a unitary transformation.  Set $w=u-v$, $d=h(w,u)$.  If $d\ne0$, the
map
\[
 R(x)=x-wd^{-1}h(w,x)
 \tag{GMR.18}
\]
sends $u$ to $v$ and is unitary, since
$h(w,w)=d+\bar d$ makes the two cross terms in
$h(Rx,Ry)-h(x,y)$ cancel its quadratic term.
If $d=0$, apply the same construction to $u,-v$, for which the
corresponding denominator is $2t\ne0$, and then negate on the current
orthogonal complement.  Working successively in the orthogonal
complements of already matched nonisotropic vectors extends every
isometry of nondegenerate subspaces.

The same argument gives extension on degenerate subspaces when it is
needed below.  Split a subspace into its radical and a nondegenerate
complement and first match the latter.  For a basis $e_i$ of the
remaining totally isotropic subspace, choose dual vectors $f_i$ with
$h(e_i,f_j)=\delta_{ij}$.  Replace the column $f$ by
$f-e\,h(f,f)/2$; its Gram matrix becomes zero and the dual pairings
are unchanged.  Thus $e,f$ span a hyperbolic space.  Do this on both
sides and match the two hyperbolic spaces by the nondegenerate
extension just proved.  This extends the original subspace isometry.
For $A=k\times k$, the Hermitian space is a perfect pairing of two
$k$-spaces; extending dual bases proves the same statement directly.
These arguments include the real unitary groups.

It follows locally that every tuple with Gram matrix $T$ is $g_vX_0$.
At almost all finite places the tuple and $X_0$ are integral and have
unimodular Gram matrix.  The isometry can then be chosen integral.
Here are the lifting details.  Over the residue field, GMR6's
orthogonal-complement argument supplies an isometry of the two frames.
The equation $g^*Jg=J$ is smooth over the valuation ring: its derivative
is onto the Hermitian matrices because two and $J$ are units.  Hence
the residue isometry lifts successively modulo $\varpi^m$ to an
integral unitary element.  If its frame is correct only modulo
$\varpi^m$, the next error $Z$ satisfies
$X_0^*JZ+Z^*JX_0=0$ modulo $\varpi$.  Relative to the nondegenerate
span of $X_0$ and its orthogonal complement, extend
$X_0\mapsto Z$ to a skew-adjoint endomorphism by choosing its other
off-diagonal block to be the negative adjoint and its final diagonal
block to be zero.  The correction $1+\varpi^m A$ is unitary to the
next order and corrects that frame error.  Smooth lifting corrects its
remaining unitary error at higher order.  Iteration converges in the
complete valuation ring.  Thus the local choices lie in
$U(J)(\mathcal O_v)$ almost everywhere and assemble to an actual
adelic unitary element.  This proves the orbit equality.

We also check its topology.  On the independent locus the tangent
space to the Gram fiber consists of the $Z$ with
$X^*JZ+Z^*JX=0$.  Any such partial map $X\mapsto Z$ extends to a
skew-adjoint endomorphism: extend the columns of $X$ to a basis, form
the sesquilinear matrix whose prescribed columns are
$h(-,Z)$, fill the prescribed rows by skew-adjointness, and put zero
in the unspecified diagonal block.  The condition on $Z$ is exactly
the compatibility on the overlapping prescribed block.  Multiplication
by $J^{-1}$ recovers the desired endomorphism.  The orbit differential
is therefore onto.  The inverse function theorem gives analytic local
sections and makes the local orbit map open.

At a good integral place the same argument works over the valuation
ring.  The compact fiber $C_v$ has a finite clopen cover on which there
are continuous sections with values in $U(J)(\mathcal O_v)$; a finite
disjoint refinement patches these sections to a continuous section
on $C_v$.  Products of these sections, with the finitely many other
local sections, give a continuous section on an adelic neighborhood.
Consequently the continuous bijection from the group quotient to
the fiber is an open map, hence a homeomorphism.

Finally every local $g_v$ preserves $\omega_{v,T}$ by GMR5, and at
almost every place it also preserves $C_v$.  Formula (GMR.16) on
cylinder sets proves invariance of $\nu_T$; those sets generate the
Borel sigma algebra.  No normalization scalar occurs in this action.
\end{proof}
% END GMR10

% BEGIN GMR11
\paragraph{GMR11. Quotient normalization for every section.}
\label{claim:GMR11}
Let a second countable locally compact group $G$ act transitively on
a locally compact space $Y$, let $y_0\in Y$ have closed stabilizer $H$,
and suppose $G/H\simeq Y$ with continuous local sections of the orbit
map.  Let $\rho$ be a nonzero invariant Radon measure on $Y$, positive
on nonempty open sets.  For every supplied left Haar measure $\mu_G$
there is a unique left Haar measure $\mu_H$ for which
\[
 \int_G f\,d\mu_G
   =\int_Y\int_H f(s(y)h)\,d\mu_H(h)\,d\rho(y)
 \tag{GMR.19}
\]
for all nonnegative measurable $f$ and all sections $s$, with no
measurability requirement on the section.  Measures supported on an
orbit in an ambient space satisfy the identical formula on that
ambient space by extension by zero.

\begin{proof}
Start with any left Haar measure $\eta$ on $H$.  For $f\in C_c(G)$,
the integral $q_f(y)=\int_H f(s(y)h)\,d\eta(h)$ is independent of $s$:
two representatives of the same coset $gH$ differ by right
multiplication of $s(y)$ by an element of $H$, and that change is a
\emph{left} change of variable in the inner integral.
Using a continuous local section shows that $q_f$ is continuous.
Indeed, on a compact neighborhood in its domain, the contributing
$h$ lie in the compact set formed from that section's image and
$\operatorname{supp}f$, intersected with $H$; dominated convergence
applies on this fixed compact set.  Its support is contained in the
compact image of $\operatorname{supp}f$ in $Y$.

Choose a countable cover by section domains and disjointify it to
obtain a Borel section.  For arbitrary nonnegative measurable $f$, its
inner integral is then measurable, and integrating it against $\rho$
defines a measure $L_\eta$ on $G$.  Countable additivity follows by
monotone convergence.  The preceding compact-support argument proves
local finiteness.  Positivity of both measures on neighborhoods gives
$L_\eta\ne0$.
For $g\in G$, write $g s(y)=s(gy)c(g,y)$ with $c(g,y)\in H$.
Left invariance of $\eta$ and invariance of $\rho$ give
$L_\eta(f\circ L_g)=L_\eta(f)$.  Thus $L_\eta$ is a left Haar measure
and $L_\eta=c\mu_G$ for some finite $c>0$.
Set $\mu_H=c^{-1}\eta$.  This proves (GMR.19).

Every other Haar measure on $H$ is a positive multiple of $\eta$;
a nonzero compactly supported test function in (GMR.19) determines
that multiple.  This proves uniqueness.  For any other section,
including a nonmeasurable one, the inner integral is pointwise the
same function of $y$ as for the Borel section.  Thus the original
all-section formula follows literally.  Values away from the orbit
are irrelevant when its complement has measure zero.
\end{proof}
% END GMR11

% BEGIN GMR12
\paragraph{GMR12. The full original global stabilizer-measure supplier.}
\label{claim:GMR12}
Let $E/F,c,n,J$ satisfy the original CM and nondegeneracy hypotheses.
Retain the original premises \texttt{SWa1} and
\texttt{Stabilizer\allowbreak First\allowbreak Chern\allowbreak Compatibility}.  For \emph{every} character
$\psi$ satisfying \texttt{Is\allowbreak Char\allowbreak AF}, every $a$ with $1\leq a$ and
$2a<n$, every rational ordered $a$-tuple $X_0$ whose Gram matrix
$T$ is nondegenerate, every witness that $T$ is Hermitian, and every
supplied Haar measure $\mu_G$ on $U(J)(\mathbb A_F)$, there exist
$\nu_T$ and a Haar measure $\mu_{X_0}$ on the actual pointwise
stabilizer of $X_0$ such that
\[
 \texttt{Is\allowbreak Om\allowbreak T}(\psi,T,\nu_T),\qquad
 \texttt{Is\allowbreak Compat\allowbreak T}(X_0,\mu_G,\mu_{X_0},\nu_T).
 \tag{GMR.20}
\]
This is the full statement of
\texttt{exists\_\allowbreak adelic\allowbreak Torus\allowbreak Measure\_\allowbreak stabiliser\allowbreak Haar}.

\begin{proof}
Construct $\nu_T$ by GMR9.  GMR10 identifies its support with the
actual orbit, proves invariance, and provides the quotient topology
and local sections.  The group is a closed matrix group in the
adelic general linear group; it is locally compact and second
countable.  Its stabilizer is closed, since it is the zero set of
the continuous equations $gX_0=X_0$.  Apply GMR11 to the supplied
$\mu_G$, not to a newly chosen normalization.  Its conclusion is
exactly the all-section formula and support condition in
\texttt{Is\allowbreak Compat\allowbreak T}.  GMR9 supplies all the quantifiers in
\texttt{Is\allowbreak Om\allowbreak T}.  A witness of Hermitian symmetry merely specifies
the same element of the Hermitian subspace and changes none of these
measures.

This proves the original supplier under its exact premises.
The proof in fact did not use \texttt{SWa1} or
\texttt{Stabilizer\allowbreak First\allowbreak Chern\allowbreak Compatibility}: their role in the source
construction is replaced here by the explicit local construction and
the positive product in GMR6--GMR9.  This observation does not prove
either of those premises and does not assert any equality of
arithmetic quotient volumes.
\end{proof}
% END GMR12

\subsection{The local flag stabilizer and its product normalization}

Fix a standard rational frame
$P=(e_1,\ldots,e_r,f_1,\ldots,f_r,z)$, meaning
\[
 h(e_i,e_j)=h(f_i,f_j)=0,\quad h(e_i,f_j)=\delta_{ij},\quad
 h(e_i,z)=h(f_i,z)=0,\quad h(z,z)=-T_0,\quad T_0\in F^\times.
 \tag{GMR.21}
\]
Set $F_0=(e_1,\ldots,e_r,z)$ and
$T_F=\operatorname{diag}(0,\ldots,0,-T_0)$.
At a fixed finite place, let $S_0$ be the pointwise stabilizer of $F_0$
and let $S$ consist of the unitary elements acting upper
unitriangularly on this ordered tuple.  These are the literal
\texttt{S0} and \texttt{Ssub}, not conjugate replacement groups.
Let $\mathcal U=U_{r+1}(A_v)$ and identify its underlying space with
$A_v^{\{(i,j):i<j\}}$ by its strict upper entries.

% BEGIN GMR13
\paragraph{GMR13. The independent singular fiber is the local orbit.}
\label{claim:GMR13}
The tuple $F_0$ is independent and its Gram matrix is $T_F$.
The independent part of $\gamma^{-1}(T_F)$ is exactly $G_vF_0$.
Its canonical measure $\omega_{v,T_F}$ is a nonzero invariant Radon
measure on that orbit, and the orbit map has continuous local
sections.  For every supplied Haar measure $\mu_{G,v}$ there is a
Haar measure $\mu_0$ on the actual $S_0$ satisfying the full original
\texttt{Is\allowbreak Compat\allowbreak Loc} with $\omega_{v,T_F}$.
More generally, prescribe arbitrary additive Haar measures $dX,dT$
on the two spaces and any family $\mathrm{fib}$ satisfying their
original \texttt{Is\allowbreak Fibre\allowbreak Family} condition.  The same conclusion holds
with $\mathrm{fib}(T_F)$ in place of the canonical measure, for every
supplied $\mu_{G,v}$.  Thus the full original
\texttt{exists\_\allowbreak haar\_\allowbreak frame\allowbreak Stabiliser\_\allowbreak compatible\_\allowbreak fibre} is included;
its two additive measures need not be self-dual.

\begin{proof}
If $\sum_i a_i e_i+bz=0$, pairing with the $f_j$ gives $a_j=0$ and
then pairing with $z$ gives $b=0$.  This proves independence, also
after scalar extension.  Its Gram matrix is (GMR.21).
For any other independent tuple of that Gram matrix, the last
vector is nonisotropic and its first $r$ vectors span a totally
isotropic subspace of its orthogonal complement.  The degenerate
subspace extension proved in GMR10 sends the ordered tuple $F_0$
to this tuple.  Thus the whole independent fiber is one actual orbit.

GMR4 makes its fiber measure Radon on the orbit and strictly positive
near every point; $F_0$ supplies a point, so it is nonzero.
GMR5 proves invariance.  The skew-adjoint extension of a tangent
map in GMR10 applies to any independent subspace, and therefore
gives local sections here as well.  The orbit is locally closed:
it is the closed Gram level set intersected with the open
independent locus.  Its stabilizer is closed and its quotient has
this orbit topology.  GMR11 now supplies $\mu_0$ relative to the
\emph{supplied} $\mu_{G,v}$ and proves the formula for every section.
GMR5 identifies the fiber measure with the source's
\texttt{om\allowbreak At} at $T_F$.  This proves the original local
frame-stabilizer compatibility supplier even though $T_F$ is singular
when $r>0$.
For the more general assertion, perform the GMR4 chart construction
with exactly the given $dX,dT$.  The positive chart densities give
the same Radon and nonvanishing properties.  GMR5's pointwise
uniqueness identifies the given family with this constructed one,
and its determinant-one covariance gives unitary invariance for
these normalizations as well.  GMR11 now applies without any change.
\end{proof}
% END GMR13

% BEGIN GMR14
\paragraph{GMR14. Unitary Haar measure is unimodular.}
\label{claim:GMR14}
Every Haar measure on the actual local unitary group $G_v$ is both
left and right invariant.  In particular conjugation and right
translation by the elements of $S$ preserve the supplied Haar measure
on $G_v$.

\begin{proof}
The Lie algebra consists of $A_v$-linear endomorphisms $A$ with
$A^*J+JA=0$.  On it the trace pairing
$(A,B)\mapsto\operatorname{Tr}_{A_v/k_v}\operatorname{tr}(AB)$ is
nondegenerate.  To verify this, extend scalars from $k_v$ to the
quadratic etale algebra: the unitary Lie algebra becomes the full
matrix Lie algebra and this pairing becomes a nonzero scalar
multiple of its nondegenerate trace pairing.  Faithful scalar
extension then proves nondegeneracy over $k_v$ itself.  In the split
case this is also immediate from the description as
$\mathfrak{gl}_n(k_v)$.

Conjugation preserves this pairing, so
$(\det\operatorname{Ad}g)^2=1$.  Its absolute determinant is therefore
one.  In local analytic coordinates a Haar density may be formed
from a left-invariant top differential form.  Conjugation by $g$
multiplies that density by the absolute determinant of its
differential at the identity, namely
$|\det\operatorname{Ad}g|=1$.
Equivalently, comparing the translated coordinate densities proves
the familiar modular formula directly.  Hence conjugation preserves
Haar measure.  Composing conjugation with a left translation gives
a right translation, so right translations preserve it as well.
Haar uniqueness makes this conclusion independent of the initial
positive normalization.
\end{proof}
% END GMR14

% BEGIN GMR15
\paragraph{GMR15. Every element of the flag stabilizer preserves its kernel Haar measure.}
\label{claim:GMR15}
For every Haar measure $\eta$ on $S_0$ and every $s\in S$,
conjugation by $s$ preserves $\eta$.  In particular $S_0$ is
unimodular.  This proves the normalization assertion needed in the
original \texttt{similitude\_\allowbreak conj\_\allowbreak preserves\_\allowbreak unitary\_\allowbreak haar\_\allowbreak at\_\allowbreak finite\_\allowbreak place}
at the actual frame and for every Haar normalization.
In that original interface the additional prescribed additive Haar
measures and the prescribed fiber family may be any allowed ones;
the conclusion does not depend on them and our proof does not use
them.  Only the actual standard frame and the supplied Haar measure
on $S_0$ occur in the conclusion.

\begin{proof}
Let $u\in\mathcal U$ be the unique matrix with $sF_0=F_0u$.
The map $R_u:X\mapsto Xu$ commutes with the left $G_v$ action and
preserves the independent fiber of $T_F$, since $u^*T_Fu=T_F$.
It preserves $\omega_{v,T_F}$ by GMR5.  It also shows that $s$
normalizes $S_0$.

First use the measure $\mu_0$ fixed by GMR13's quotient formula.
If $C_s(h)=shs^{-1}$, its pushforward is another Haar measure,
so $(C_s)_*\mu_0=c\mu_0$ for a finite $c>0$.
Choose a section $q$ of the orbit map and let
$q'(y)=q(R_u^{-1}y)s$; this is again a section because
$sF_0=R_uF_0$ and $R_u$ commutes with the action.
Apply GMR13's formula to $f(\,\cdot\,s^{-1})$ using $q'$.
Its inner argument is
\[
 q'(y)h s^{-1}=q(R_u^{-1}y)\,s h s^{-1}.
\]
GMR14 says that the left integral is $\int f\,d\mu_{G,v}$.
Changing $y$ by $R_u$ and then using the Haar factor $c$ says that
the right integral is $c\int f\,d\mu_{G,v}$.
Choose a nonnegative compactly supported $f$ with finite positive
integral.  We conclude $c=1$.

Every other Haar measure $\eta$ is a positive scalar multiple of
$\mu_0$, so it is preserved as well.  Taking $s\in S_0$ proves
invariance under inner automorphisms of $S_0$.  Together with left
invariance this is right invariance, proving unimodularity.
Thus the product construction below uses an established exact
normalizer action, rather than an unspecified modular factor.
\end{proof}
% END GMR15

% BEGIN GMR16
\paragraph{GMR16. The actual upper-unitriangular quotient and a continuous section.}
\label{claim:GMR16}
The map $\pi:S\to\mathcal U$ determined by $sF_0=F_0\pi(s)$ is a
continuous surjective homomorphism with kernel $S_0$.  It has a
continuous section $t$.  Consequently
\[
 S_0\times\mathcal U\longrightarrow S,\qquad (s_0,u)\longmapsto s_0t(u)
 \tag{GMR.22}
\]
is a homeomorphism.  The source measure
$\texttt{sd\allowbreak E}(\psi,v,\texttt{UIdx}(r))$ is a left and right Haar measure
on $\mathcal U$ under its strict-upper-coordinate identification.

\begin{proof}
Independence of $F_0$ makes its coefficients unique; taking products
proves the homomorphism assertion and taking coordinates proves
continuity.  Its kernel is the pointwise stabilizer, exactly $S_0$.
To prove surjectivity with the needed topology, write
\[
 u=\begin{pmatrix}A&b\\0&1\end{pmatrix},
 \qquad A\in U_r(A_v),\quad b\in A_v^r.
\]
Use the frame's nondegenerate subspace spanned by $e,f,z$, and let the
following map act as the identity on its orthogonal complement:
\[
\begin{split}
 t(u)e&=eA,\\
 t(u)z&=z+eb,\\
 t(u)f&=\left(f+T_0^{-1}z b^*
                  +(2T_0)^{-1}e b b^*\right)A^{-*}.
\end{split}
 \tag{GMR.23}
\]
Here $e,f$ denote the matrices with those columns and
$A^{-*}=(A^*)^{-1}$.  Formula (GMR.21) gives
$h(t e,t f)=1$, $h(tz,t f)=0$ and $h(t f,t f)=0$:
in the last identity the two terms with coefficient $(2T_0)^{-1}$
cancel the term $-T_0\cdot T_0^{-2}$ from the $z$ coordinate.
The other pairings are unchanged.  Thus $t(u)$ is unitary and has
the prescribed action on $F_0$.  Its formulas are continuous in
$A,b$, including $r=0$, when the section is the identity.

Every $s\in S$ has the unique expression
$(s t(\pi(s))^{-1})t(\pi(s))$.  Its first factor belongs to $S_0$,
and these inverse formulas are continuous, proving (GMR.22).

Finally in the strict-upper coordinates the product law is
\[
 (xy)_{ij}=x_{ij}+y_{ij}+\sum_{i<k<j}x_{ik}y_{kj}.
 \tag{GMR.24}
\]
For fixed $x$, this is an affine triangular change of the $y$
coordinates with diagonal entries one, ordered by increasing
$j-i$; the same holds with the roles reversed.  It preserves
additive Haar measure by repeated one-coordinate translations and
Fubini, or by its determinant one.  GMR2 identifies the source
\texttt{sd\allowbreak E} with a nonzero additive Haar measure on this coordinate
space.  Therefore it is also left and right Haar for the group law
(GMR.24), with exactly that self-dual normalization.
\end{proof}
% END GMR16

% BEGIN GMR17
\paragraph{GMR17. The full all-section product formula on the flag stabilizer.}
\label{claim:GMR17}
Let $\mu_0$ be the measure on $S_0$ from GMR13, and let $\mu_U$ be
the literal \texttt{sd\allowbreak E} measure in GMR16.  There is a Haar measure
$\mu_S$ on the actual $S$ for which
\[
 \int_S f(s)\,d\mu_S(s)
 =\int_{\mathcal U}\int_{S_0}f(s_0t'(u))\,d\mu_0(s_0)\,d\mu_U(u)
 \tag{GMR.25}
\]
for every nonnegative measurable function on the ambient unitary
group and every map $t'$ inducing the prescribed upper-unitriangular
action on $F_0$.  In particular this is the exact original
\texttt{Is\allowbreak Prod\allowbreak U}, with no measurability assumption on $t'$.

\begin{proof}
Push forward $\mu_0\otimes\mu_U$ by (GMR.22).  This defines a
nonzero Radon measure on $S$, finite on compact sets and positive on
open sets.  To prove left invariance, fix $g\in S$ and put
$v=\pi(g)$.  There is a factor $c(g,u)\in S_0$ such that
$g t(u)=c(g,u)t(vu)$.  Thus
\[
 g s_0 t(u)=(g s_0 g^{-1})c(g,u)t(vu).
\]
GMR15 permits the first conjugation in the inner integral, and the
unimodularity of $S_0$ permits right multiplication by $c(g,u)$.
GMR16 permits the change $u\mapsto vu$ in the outer integral.
The resulting integral is the original one.  Hence the pushed-forward
measure is left invariant and is Haar.

If $t'$ has the same frame action as $t$, then
$t'(u)t(u)^{-1}\in S_0$.  Right invariance of $\mu_0$ makes the inner
integrals with $t'$ and $t$ equal pointwise in $u$.  The one with the
continuous section $t$ is measurable by Tonelli.  Thus the formula
also holds for an arbitrary $t'$; its inner integral is the same
measurable function.  Pushing $\mu_S$ to the ambient group gives
the source's \texttt{amb} convention in \texttt{Is\allowbreak Prod\allowbreak U} exactly.
\end{proof}
% END GMR17

% BEGIN GMR18
\paragraph{GMR18. The full original normalized local frame-stabilizer supplier.}
\label{claim:GMR18}
Assume the original CM, Hermitian, \texttt{SWa1} and
\texttt{Stabilizer\allowbreak First\allowbreak Chern\allowbreak Compatibility} premises.  Let $\psi$
satisfy \texttt{Is\allowbreak Char\allowbreak AF} and \texttt{Is\allowbreak Std\allowbreak Char}.  For every finite
place $v$, every $r$, every actual standard frame $P$ with
$2(r+1)<n$, and every supplied Haar measure $\mu_{G,v}$, there is a
Haar measure $\mu_{S,v}$ on the actual $\texttt{Ssub}(P)$ such that
\[
 \texttt{Loc\allowbreak Tie\allowbreak W}(c,J,\psi,v,P,\mu_{G,v},\mu_{S,v}).
 \tag{GMR.26}
\]
This is the full original
\texttt{exists\_\allowbreak frame\allowbreak Stabilizer\allowbreak Haar\_\allowbreak normalised\_\allowbreak fibre\allowbreak Measure}.

\begin{proof}
Use GMR2--GMR5 at $a=r+1$ to obtain the source's selected self-dual
measures and its selected entire fiber family.  Since $T_F$ is
Hermitian, \texttt{om\allowbreak At} is that family's value at $T_F$, not its
zero fallback.  GMR13 gives a Haar measure $\mu_0$ on the literal
pointwise stabilizer $S_0$ with \texttt{Is\allowbreak Compat\allowbreak Loc} relative to the
supplied $\mu_{G,v}$.  GMR2 and GMR16 show that the selected
\texttt{sd\allowbreak E} exists and is the required Haar measure in the
upper-unitriangular coordinates.  GMR15 and GMR17 construct $\mu_{S,v}$
and prove the all-section \texttt{Is\allowbreak Prod\allowbreak U} formula.

These are exactly the three entries of \texttt{Loc\allowbreak Tie\allowbreak W}: the witness
$\mu_0$ is Haar, is compatible with \texttt{om\allowbreak At}, and gives the
product formula with \texttt{sd\allowbreak E}.  All the ambient groups, the
character, the frame and the supplied $\mu_{G,v}$ have been kept.
The unused original volume premises are retained in the statement;
no local density or normalization supplier has been assumed in their
place.  The argument even allows $r=0$: $\mathcal U$ is the trivial
group with self-dual mass one, and (GMR.25) then gives
$\mu_S=\mu_0$.
\end{proof}
% END GMR18

\subsection{Simultaneous rescaling without changing the supplied bridge}

% BEGIN GMR19
\paragraph{GMR19. The full original \texttt{Rescale\allowbreak Step3}.}
\label{claim:GMR19}
Assume the original \texttt{Siegel\allowbreak Weil\allowbreak Unitary}\allowbreak
\texttt{Convergent} premise.
For every original $D$, every supplied $X:\texttt{Hida\allowbreak Setting}(D)$,
every supplied $h_b:\texttt{Ad\allowbreak Bridge\allowbreak Iso}(X.B.I)$, and every witness
of $X.\texttt{Meas\allowbreak Tie}(h_b)$, there are normalized measure records
$m_G$ and $m_W(u)$ such that
\[
 X'=\texttt{with\allowbreak Meas}(X,m_G,m_W)
 \tag{GMR.27}
\]
has all the following properties simultaneously:
\[
\begin{gathered}
 X'.\texttt{Meas\allowbreak Tie}(h_b),\qquad
 \exists\psi:\mathbb A_F\longrightarrow\mathbb C^\times,\quad
 \texttt{Is\allowbreak Char\allowbreak AF}(\psi)\ \wedge\ \texttt{Is\allowbreak Std\allowbreak Char}(\psi),\\
 \texttt{MeasCompat\ensuremath{\psi}}(X',\psi),\qquad
 \forall w\ \texttt{Loc\allowbreak Compat\allowbreak X}(X',\psi,w).
\end{gathered}
 \tag{GMR.28}
\]
Here $F=X.B.I.F$ and every group, frame and tensor coordinate is the
one belonging to $X$.  We can, and do, leave
$\mu_G,\mu_{G,\infty},\mu_{G,f},\mu_{G,v}$ unchanged and retain
\[
 m_W(u).\mu_{N,\infty}=X.\texttt{meas\allowbreak W}(u).\mu_{N,\infty},\qquad
 m_W(u).\mu_{H,\infty}=X.\texttt{meas\allowbreak W}(u).\mu_{H,\infty}
 \tag{GMR.29}
\]
for every word.  The same single $\psi$ works for all words, all
finite places, and all auxiliary tuples of length $2D.R.rS+2$.

\begin{proof}
Put $r=D.R.rS$, $n=X.B.I.n$.  The setting's CM and Hermitian fields
give the hypotheses used above.  Its rank equation is
$n=D.R.k+D.R.m-1$.  Write $D.R.m=2(r+1)$, using evenness and
$D.R.m\geq2$.  Since $D.R.k\geq D.R.m+3$, we have
$n\geq4r+6$, in particular the source bound $4r+5\leq n$.
Thus
\[
 2(2r+2)<n,\qquad 2(2r+1)<n,\qquad 2(r+1)<n.
 \tag{GMR.30}
\]

For every word $u$, let $P_u$ be the source's actual
$\texttt{frame\allowbreak Of}(X,u)$.  The setting's basis and its equality for
$J$ give
\[
 h_J(\operatorname{crd}x,\operatorname{crd}y)
   =X.\iota_E^{-1}(h_{VS}(x,y)).
 \tag{GMR.31}
\]
Indeed expand both vectors in that basis, use its $J$ entries, and
use that $X.\iota_E$ respects the involutions.  The stored
$ee,ff,ef,zu\_\allowbreak orth,zu\_\allowbreak norm$ equations therefore give (GMR.21)
for $P_u$, with $T_u$ the nonzero rational integer specified by the
word.  It is nonzero because that integer is prime in the original
word datum.  The ordered tuple
$\texttt{tup\allowbreak H}(P_u)=(e_1,\ldots,e_r,z,f_1,\ldots,f_r)$ has
nondegenerate Gram matrix: after moving $z$ last it is
\[
 \begin{pmatrix}0&1&0\\1&0&0\\0&0&-T_u\end{pmatrix},
 \tag{GMR.32}
\]
where the first two diagonal blocks have size $r$.  Its determinant
is $(-1)^r(-T_u)\ne0$.  This checks the original standard-frame
and rank suppliers directly on $X$.

Choose \emph{now}, once and for all, the standard character of GMR1
on $\mathbb A_F$.  No word, place, or auxiliary tuple has been chosen
at this stage.  Specializing the retained Siegel--Weil premise at
the actual $F,E,c,n,J$ gives its finite-volume and smooth-density
components.  By their literal definitions these are the source's
\texttt{SWa1} and \texttt{Stabilizer\allowbreak First\allowbreak Chern\allowbreak Compatibility} premises
in GMR12 and GMR18.  We use those original statements at this scope;
their constructive proofs above did not otherwise require these
volume components.

First handle all auxiliary tuples $X_0$ of length $2r+2$.
If $X_0$ has nondegenerate Gram matrix, apply GMR12 with the
unchanged global measure $X.\texttt{meas\allowbreak G}.\mu_G$.
Choose its stabilizer Haar measure $\mu'_B(X_0)$ and fiber measure
$\nu_{X_0}$.  If its Gram matrix is degenerate, retain the existing
$X.\texttt{meas\allowbreak G}.\mu_B(X_0)$; the source imposes no
\texttt{MeasCompat\ensuremath{\psi}} fiber condition on that tuple.
Thus a single choice function $\mu'_B$ is Haar for every tuple and
has the required compatibility for every nondegenerate one.

Next fix a word $u$ temporarily.  Apply GMR12 to the actual tuple
$\texttt{tup\allowbreak H}(P_u)$, of length $2r+1$, using (GMR.30)--(GMR.32).
This gives a global Haar measure $\mu_{H,u}$ and its global fiber
measure, compatible with the same $\mu_G$ and the same $\psi$.
For every finite place $v$, apply GMR18 to $P_u$ and the unchanged
supplied $X.\texttt{meas\allowbreak G}.\mu_{G,v}$; choose the resulting
$\mu_{S,u,v}$ with its full \texttt{Loc\allowbreak Tie\allowbreak W}.
Apply AMD17 with precisely the old $\mu_{N,\infty}$ and
$\mu_{H,\infty}$ for this word, this newly fixed $\mu_{H,u}$,
and this family $\mu_{S,u,v}$.  It produces a normalized
$m_W(u)$ retaining all these entries and with
$\texttt{covol}(m_W(u).\mu_N)=1$.

To make the normalization transparent, AMD16 first chooses the
global $N$-measure of literal covolume one.  AMD13 absorbs the
two required global-to-infinite Haar scalars into $\mu_{N,f}$ and
$\mu_{H,f}$ separately.  AMD7 supplies the finite $S$-measure as
the product pushforward of those two finite factors.  AMD14 verifies
all eleven \texttt{Meas\allowbreak W.\allowbreak Is\allowbreak Text\allowbreak Norm} clauses.  In particular neither
measure in (GMR.29), the prescribed global $H$-measure, nor any
prescribed local $S$-measure is rescaled.  The source record does
not require an additional equality identifying this finite
$S$-measure with a restricted product of the local $S$-measures;
no such extra equality is used here.
Choosing these witnesses for all words gives the required family
$m_W$.  All choices have occurred after the common $\psi$ was fixed.

Define $m_G$ by changing only $\mu_B$ to $\mu'_B$ in the old
\texttt{Meas\allowbreak G} record.  Its Haar clauses hold and its product
splitting clause is unchanged, so it satisfies the full
\texttt{Meas\allowbreak G.\allowbreak Is\allowbreak Text\allowbreak Norm}.  Form $X'$ by (GMR.27).
The first three entries of \texttt{Meas\allowbreak Tie} are unchanged: they
refer to the same $\mu_G,\mu_{G,\infty},\mu_{G,f}$, the same
block measure $X.B.\mu_f$, and the same maps through $h_b$.
Its last two entries, for each word, are the ambient pushforwards
of $\mu_{N,\infty}$ and $\mu_{H,\infty}$; (GMR.29) makes those
entries unchanged as well.  Thus the exact supplied bridge $h_b$
still witnesses $X'.\texttt{Meas\allowbreak Tie}$.

Finally unfold \texttt{MeasCompat\ensuremath{\psi}}.  Its first clause is,
for each word, the global $H$ fiber compatibility together with
unipotent covolume one; these were just supplied.  Its second clause
is global compatibility for every nondegenerate auxiliary tuple;
that is the choice of $\mu'_B$.  Unfolding
\texttt{Loc\allowbreak Compat\allowbreak X} gives \texttt{Loc\allowbreak Tie\allowbreak W} at every finite place
for each word, exactly the local choices above.  We have therefore
proved every entry of (GMR.28), with the original order and range
of quantifiers.  This is the full original \texttt{Rescale\allowbreak Step3}.
\end{proof}
% END GMR19

\subsection{Exact next use and retained global analytic inputs}

GMR19 applies directly to the supplied-form core setting of BQS23,
which already consumes ASG, AMD and the completed Bessel chain.
It adds \texttt{MeasCompat\ensuremath{\psi}} and all \texttt{Loc\allowbreak Compat\allowbreak X} to
that \emph{same} setting while preserving its bridge and its
infinite measures.  It does not construct the original supplied
\texttt{Forms0} datum.  The source's stronger all-family
\texttt{exists\_\allowbreak setting} subsequently uses the entire family
$F(u)$ to produce \texttt{Transport} and \texttt{Level\allowbreak Match};
that step has distinct genus, level-conjugacy and representation
intertwiner producers.  If such a transport is already supplied,
\texttt{tr\allowbreak Meas} keeps its Witt map, adelic algebra map, representatives
and representation intertwiner literally unchanged under
\texttt{with\allowbreak Meas}.  No new transport existence is inferred from a
measure normalization.

For clarity, the retained source premise
\texttt{Siegel\allowbreak Weil\allowbreak Unitary\allowbreak Convergent} says the following for every CM
extension and every nondegenerate Hermitian ambient space.  Haar
measures and arithmetic fundamental domains exist, and each such
fundamental domain has finite Haar mass.  For every tuple length
$a$ with $2a<n$, every nondegenerate rational Hermitian $T$, every
rational coefficient basis, every factorizable test function whose
infinite part is smooth, compactly supported and rapidly decreasing,
and every pure finite Schwartz--Bruhat factorization, all the source
Haar, fundamental-domain, character and self-dual-measure choices
admit continuous Gram-pushforward densities on the nondegenerate
Hermitian locus.  Their values at $T$ are multipliable, and their
infinite value times the finite product equals the normalized
arithmetic theta mean.  The source defines that mean as the complex
integral divided by the real mass of the chosen fundamental domain.
This exact universal premise is retained; the local disintegration
and positive product proved here are not an assertion of that
arithmetic theta identity.

In particular neither equality of the arithmetic stabilizer and
ambient volumes, nor the later orbital-counting identity, follows
just by replacing its Haar measures.  The next original consumer
\texttt{in\allowbreak Volumes\_\allowbreak of\_\allowbreak compat} still takes its named
\texttt{In\allowbreak Volumes\allowbreak At} input.  The actual
\texttt{Count\allowbreak Idempotents\allowbreak Class}, local Bessel uniqueness, the
all-family column transports, and the final global spectral
identifications remain separate source obligations.  This chapter
closes the measure-rescaling supplier and its two analytic measure
constructors, with their complete original scopes, rather than a
whole selected global obligation.


\section{Wach lattices from one prismatic realization}
\label{sec:wach-lattice-descent-v053}

\status{This chapter proves the algebraic descent, local-prism comparison,
and lattice assembly used in the original Wach-lattice producer. The
lattice is a submodule of the \emph{supplied} \'etale module, with its
original Frobenius, group action and maps of lattices. The final all-layer
return retains three precisely stated period-ring suppliers: the Wach
prism, injectivity and reflection of units for the coefficient embedding,
and descent of cyclotomic-kernel-fixed Frobenius matrix solutions.
Their existence is not inferred from the assembly proved here.}

The obstruction addressed here is more concrete than the instruction to
choose a coherent Wach lattice suggests. A crystalline lattice has
a prismatic realization. Its value on one local prism gives a lattice
inside a Witt-vector tensor module. We first show that this lattice
comes from the given \'etale module, then take its inverse image there.
Fullness of the crystalline realization transports every lattice
homomorphism through this construction. Thus the same inverse-image
definition provides all the required functoriality.

The finite-dimensional coefficient arguments are written out because
neither faithful flatness of an unspecified period-ring extension nor
a previously constructed Wach functor is an input. In particular,
the crucial base-change isomorphism below follows from a determinant
argument and a tensor normal form. All rings in the algebraic lemmas
are commutative and all sums over finite sets include the empty case.

\subsection{Coefficients, completion and denominator matrices}

% BEGIN WLD1
\paragraph{WLD1: tensor cancellation with its actual maps.}
\label{claim:WLD1}
Let $Z\to O$ and $Z\to W$ be ring maps and let $T$ be an $O$-module.
Put $B=O\otimes_ZW$. There is a natural additive isomorphism
\[
 \kappa_T:W\otimes_ZT\longrightarrow B\otimes_OT,\qquad
 w\otimes t\longmapsto(1\otimes w)\otimes t.                 \tag{WLD.1}
\]
With $c(w)=1\otimes w$, it satisfies
$\kappa_T(wy)=c(w)\kappa_T(y)$. It commutes with every map of
lattices and every coefficient operator whose actions agree on
$c(W)$. If $b_1,\ldots,b_q$ is a $W$-basis of $W\otimes_ZT$,
then every $y\in B\otimes_OT$ has a unique expression
\[
                 y=\sum_i c(w_i)\kappa_T(b_i),\qquad w_i\in W.
                                                               \tag{WLD.2}
\]

\begin{proof}
Define the inverse on pure tensors by
\[
             (o\otimes w)\otimes t\longmapsto w\otimes ot.
\]
The $Z$-balancing relation in the first tensor product and the
$O$-balancing relation in the second are both preserved, so this
defines an additive map. Its two composites with (WLD.1) are the
identity on pure tensors and hence on all tensors. The scalar
formula follows on pure tensors from multiplication in $W$.
For operators $\sigma:W\to W$, $\sigma':B\to B$ with
$\sigma'c=c\sigma$ and maps $f,f'$ agreeing on $T$, both possible
composites send $w\otimes t$ to $c(\sigma w)\otimes f(t)$.
Additivity proves the assertion for arbitrary tensors. Finally apply
$\kappa_T$ to the unique basis expansion of $\kappa_T^{-1}(y)$.
These checks apply to the Frobenius and Galois operators as well as
to homomorphisms of the original integral lattices.
\end{proof}
% END WLD1

% BEGIN WLD2
\paragraph{WLD2: finite coefficient coordinates and reflection of units.}
\label{claim:WLD2}
If $o_1,\ldots,o_d$ is a $Z$-basis of $O$, every element of
$O\otimes_ZW$ is uniquely $\sum_j o_jc(w_j)$. If this tensor ring
is nonzero, $c(w)$ being a unit implies that $w$ is a unit of $W$.
If $T$ is free of rank $r$ over $O$, its underlying $Z$-module and
$W\otimes_ZT$ are free of rank $dr$ over $Z$ and $W$, respectively.

\begin{proof}
Tensor the coordinate isomorphism $O\simeq Z^d$ with $W$.
Multiplication by $c(w)$ is multiplication by $w$ on each of the $d$
coordinates. If it is invertible, its determinant $w^d$ is a unit.
The nonzero tensor ring excludes $d=0$; a positive power is a unit
only if its base is, since $w\cdot w^{d-1}(w^d)^{-1}=1$.
For the last assertion, if $t_1,\ldots,t_r$ is an $O$-basis of $T$,
the products $o_jt_k$ form a $Z$-basis. Tensoring gives the stated
$W$-basis. This also covers $r=0$.
\end{proof}
% END WLD2

% BEGIN WLD3
\paragraph{WLD3: lifting coefficient decompositions through completion.}
\label{claim:WLD3}
Let $O,E,W,B$ be rings, let $\iota:E\to B$ be an $O$-algebra map
and $c:W\to B$ a ring map. Suppose $o_1,\ldots,o_d\in O$ give
unique coordinates
\[
                       B=\left\{\sum_j o_jc(w_j):w_j\in W\right\}.
\]
Assume $E$ and $W$ are separated and complete for the ideals generated
by an integer $p$. If every $e\in E$ has an expression modulo $pE$
\[
 e\equiv\sum_j o_je_j\pmod {pE},\qquad \iota(e_j)\in c(W),
\]
then it has an exact expression of the same form:
\[
                  e=\sum_j o_je_j,\qquad \iota(e_j)\in c(W).
                                                               \tag{WLD.3}
\]
No division by $p$, injectivity of $c$, or a topology on $B$ is assumed.

\begin{proof}
Choose, for every $u\in E$, elements $g_j(u)\in E$, $w_j(u)\in W$
and $h(u)\in E$ such that
\[
 \iota(g_j(u))=c(w_j(u)),\qquad
                  u=\sum_j o_jg_j(u)+ph(u).
\]
Starting with $u_0=e$, put $u_{k+1}=h(u_k)$. Iteration gives
\[
 e=\sum_j o_j\sum_{k<N}p^kg_j(u_k)+p^Nu_N.                 \tag{WLD.4}
\]
Completeness supplies sums
$e_j=\sum_{k\ge0}p^kg_j(u_k)$ in $E$ and
$v_j=\sum_{k\ge0}p^kw_j(u_k)$ in $W$: concretely, the difference
from the $N$th partial sum is divisible by $p^N$.

The coefficient hypothesis makes $B$ separated for divisibility by
$p$. Indeed, if $b\in\bigcap_Np^NB$, express $b$ and each witness
$b=p^Nb_N$ in the unique $W$-coordinates. Uniqueness shows each
coordinate of $b$ is divisible by every $p^N$, hence is zero by
separatedness of $W$. Now
$\iota(e_j)-c(v_j)$ is divisible by every $p^N$ in $B$, because its
partial sums agree. Therefore it is zero. Similarly (WLD.4) shows
$e-\sum_j o_je_j$ belongs to every $p^NE$, so it is zero.
This proves (WLD.3), including cases where $p$ is a zero divisor.
\end{proof}
% END WLD3

% BEGIN WLD4
\paragraph{WLD4: coefficient decompositions in the supplied period ring.}
\label{claim:WLD4}
Let $A\hookrightarrow E$ be the power-series coefficient ring and its
localization-completion in the original \texttt{WSer} data. Write
$\pi_A$ for \texttt{letpi}. Let $a:S\to A$, $\psi:S\to W$ and
$c:W\to B$ satisfy
\[
       \iota(a(t))=c(\psi(t)),\qquad a(\pi)=\pi_A,
                                                               \tag{WLD.5}
\]
where $A\to E$ is suppressed in notation. Suppose every $\alpha\in A$
is $\sum_j o_ja(t_j)$, $B$ has the unique $W$-coordinates of WLD3,
$W$ is $p$-adically complete, and $c$ reflects units. Then the
original ring $E$ satisfies (WLD.3).

\begin{proof}
The original density assertion, applied with precision one, gives
$\alpha\in A$ and $m\ge0$ such that
$e\pi_A^m-\alpha\in pE$. The element $\pi_A$ is a unit in $E$.
Its image is $c(\psi(\pi))$, so $\psi(\pi)$ is a unit in $W$.
Write $\alpha=\sum_j o_ja(t_j)$ and set
$e_j=\pi_A^{-m}a(t_j)$. Then
\[
 e-\sum_j o_je_j\in pE,\qquad
 \iota(e_j)=c\bigl(\psi(\pi)^{-m}\psi(t_j)\bigr).
\]
The first equality uses multiplication by the unit $\pi_A^{-m}$;
it does not cancel $p$. Apply WLD3, using the completeness of $E$
included in \texttt{WSer}. Thus the desired exact coefficients
are obtained in the \emph{same} $E$, rather than by replacing it with
a new completion.
\end{proof}
% END WLD4

% BEGIN WLD5
\paragraph{WLD5: independence after clearing one common denominator.}
\label{claim:WLD5}
Let $a:S\to E$, $\iota:E\hookrightarrow B$, $c:W\to B$,
$\psi:S\to W$ satisfy $\iota a=c\psi$, and let $a(Q)$ be a unit.
Let $y_1,\ldots,y_q$ in a $B$-module be independent for coefficients
from $W$, meaning
\[
             \sum_i c(w_i)y_i=0\ \Longrightarrow\ w_i=0.
\]
Define the subring
\[
 E_Q=\{e\in E:\ a(Q)^ne=a(t)
                    \text{ for some }n\ge0,\ t\in S\}.
\]
Then $\sum_i\iota(e_i)y_i=0$ with $e_i\in E_Q$ implies $e_i=0$.

\begin{proof}
The set $E_Q$ is closed under sums, products and negatives by putting
fractions over a common power of $a(Q)$. Choose presentations
$a(Q)^{n_i}e_i=a(t_i)$ and take $N\ge n_i$ for every $i$.
Multiplying the relation by $\iota(a(Q)^N)$ gives
\[
        \sum_i c\bigl(\psi(Q^{N-n_i}t_i)\bigr)y_i=0.
\]
Independence makes every displayed $W$-coefficient zero.
The identity $\iota a=c\psi$ and injectivity of $\iota$ imply
$a(Q)^Ne_i=0$. Since $a(Q)$ is a unit, $e_i=0$.
\end{proof}
% END WLD5

% BEGIN WLD6
\paragraph{WLD6: two Frobenius span inclusions give an invertible matrix.}
\label{claim:WLD6}
In WLD5, suppose another tuple $z_1,\ldots,z_q$ satisfies
\[
 c(\psi Q)^{k_j}z_j=\sum_i c(\psi(s_{ij}))y_i,\qquad
 c(\psi Q)^{h_i}y_i=\sum_j c(\psi(t_{ji}))z_j
\]
for suitable nonnegative exponents and coefficients in $S$.
Then there is $U\in\operatorname{GL}_q(E)$ with
\[
                         z_j=\sum_i\iota(U_{ij})y_i.       \tag{WLD.6}
\]

\begin{proof}
Invert $a(Q)$ in $E$ in the first family of equations. It gives a
matrix $U$ with entries in $E_Q$ and the displayed expression for
$z_j$. The second family similarly gives a matrix $V$ over $E_Q$
expressing the $y_i$ in terms of the $z_j$. Substitute to obtain
$\sum_m\iota((UV-I)_{mi})y_m=0$. WLD5 gives $UV=I$.
For square matrices over a commutative ring this implies
$\det(U)\det(V)=1$; the adjugate formula then makes $U$ invertible.
Neither the ambient module nor the coefficient extension is assumed
flat in this argument.
\end{proof}
% END WLD6

\subsection{The exact base-change and projectivity arguments}

% BEGIN WLD7
\paragraph{WLD7: spanning the supplied module by a descended tuple.}
\label{claim:WLD7}
Let $O\to A\to E$, let $\iota:E\hookrightarrow B$ be an $O$-algebra
map reflecting units, and let $a:S\to A$, $c:W\to B$ be ring maps.
Let $D$ be free of rank $r$ over $E$, $Y$ a $B$-module, and
$j:D\to Y$ additive with $j(ex)=\iota(e)j(x)$. Assume
$\iota(a(S))\subseteq c(W)$ and
\[
 \forall e\in E\quad
 e=\sum_{j=1}^d o_je_j\quad\text{with }\iota(e_j)\in c(W).
                                                               \tag{WLD.7}
\]
Suppose $q=dr$ and $x_1,\ldots,x_q\in D$ give unique $W$-coordinates
on $Y$ through their images:
\[
                    y=\sum_i c(w_i)j(x_i).
                                                               \tag{WLD.8}
\]
Then $x_1,\ldots,x_q$ span $D$ over $E$.

\begin{proof}
Choose an $E$-basis $b_1,\ldots,b_r$. Index a tuple of length $dr$
by pairs and put $u_{jk}=o_jb_k$. Decompose each coordinate of
each $x_i$ using (WLD.7). This produces a square matrix $C'$ over
$E$ and a square matrix $C$ over $W$ with
\[
             x_i=\sum_l C'_{il}u_l,\qquad
                         \iota(C'_{il})=c(C_{il}).
\]
Use (WLD.8) to express $j(u_l)=\sum_m c(B_{lm})j(x_m)$.
Substitution and uniqueness in (WLD.8) give $CB=I$.
Thus $\det C$ is a unit. Since
$\iota(\det C')=c(\det C)$ and $\iota$ reflects units,
$\det C'$ is a unit in $E$. The inverse matrix expresses every
$u_l$ as an $E$-linear combination of the $x_i$.

Finally decompose $1=\sum_j o_je_j$ using (WLD.7). For every $k$,
$b_k=\sum_j e_ju_{jk}$, so the $b_k$, and hence all of $D$, lie
in the span of the $x_i$. This proof does not claim that the
$q=dr$ vectors are an $E$-basis; generally they are not.
\end{proof}
% END WLD7

% BEGIN WLD8
\paragraph{WLD8: the tensor map is an isomorphism without an extra flatness premise.}
\label{claim:WLD8}
Keep the assumptions of WLD7. Let $N\subset D$ be an $A$-submodule
whose elements are precisely
\[
                        \sum_i a(s_i)x_i,\qquad s_i\in S.
                                                               \tag{WLD.9}
\]
Then the actual inclusion induces an isomorphism
\[
                         E\otimes_AN\xrightarrow{\ \sim\ }D.
                                                               \tag{WLD.10}
\]

\begin{proof}
Each $x_i$ lies in $N$, so WLD7 gives surjectivity.
For injectivity, first put every tensor into the normal form
\[
              t=\sum_i e_i\otimes x_i,\qquad
                                      \iota(e_i)\in c(W).
                                                               \tag{WLD.11}
\]
For a pure tensor $e\otimes n$, decompose $e=\sum_j o_je_j$.
Move $o_j\in A$ across the tensor sign. Since $N$ is an
$A$-module, $o_jn$ has an expression (WLD.9). Move its coefficients
$a(s_i)$ back across the tensor sign. The resulting coefficients
are sums of products of $e_j$ with $a(s_i)$; their images lie in
the subring $c(W)$. Additivity proves (WLD.11) for all tensors.

If (WLD.11) maps to zero, choose $w_i\in W$ with
$c(w_i)=\iota(e_i)$. Applying $j$ gives
$\sum_i c(w_i)j(x_i)=0$. Uniqueness of (WLD.8), compared with
the zero tuple, gives $w_i=0$. Injectivity of $\iota$ gives
$e_i=0$, whence $t=0$. The map in (WLD.10) is therefore bijective
and is the original base-change map on every $1\otimes n$.
This proves the full source statement
\texttt{is\allowbreak Base\allowbreak Change\_\allowbreak lattice\_\allowbreak of\_\allowbreak basis\_\allowbreak after\_\allowbreak extension}.
\end{proof}
% END WLD8

% BEGIN WLD9
\paragraph{WLD9: projectivity over the original power-series ring.}
\label{claim:WLD9}
Let $Z\to S\xrightarrow{a}A$ be ring maps, and let $N$ be an
$A$-submodule of an $A$-module. Suppose (WLD.9) holds and the
coefficients $s_i$ in that expression are unique. If
$S$ is finite free over $Z$ with basis $v_1,\ldots,v_h$, then
$N$, with the restricted $Z$-action, is free on
\[
                              a(v_j)x_i.
\]
In particular it is projective. In the Wach application
$Z=\mathbb Z_p[[X]]$, with its specified map into $S$.

\begin{proof}
Write each $s_i$ uniquely as $\sum_j z_{ij}v_j$.
This gives existence of the indicated expression for every element
of $N$. If that expression is zero, uniqueness of the $S$-coordinates
first gives $\sum_jz_{ij}v_j=0$ for every $i$; the $Z$-basis of $S$
then gives $z_{ij}=0$. The resulting explicit isomorphism
$Z^{qh}\simeq N$ proves freeness and projectivity with exactly the
restricted scalar action. We do not infer projectivity over $A$
from this argument.
\end{proof}
% END WLD9

\subsection{What a local prism gives inside the Witt tensor module}

The word \emph{crystal} below means the precise source data
\texttt{FCrys}: a finite projective module $M(P)$ on each prism,
semilinear transition maps inducing base-change isomorphisms, their
identity and composition laws, and a Frobenius isomorphism after
localization at the prism ideal. We use its stated fraction-naturality
law, so no injectivity of a transition map is silently assumed.
A realization identifies the base change at the distinguished
$A_{\inf}$ prism with $W\otimes_{\mathbb Z_p}T$, respecting Frobenius
fractions and Galois action. Here $W=W(C^\flat)$ is the source's
\texttt{WCb}; it is a $p$-adically complete domain. These are the
original period rings, not substitute abstract realizations.

% BEGIN WLD10
\paragraph{WLD10: bases and principal kernels under semilinear base change.}
\label{claim:WLD10}
Let $f:S\to R$ be a ring map and $u:M\to N$ an $f$-semilinear map
inducing an isomorphism $R\otimes_{S,f}M\simeq N$. If $M$ has
basis $b_1,\ldots,b_q$, then $u(b_1),\ldots,u(b_q)$ is an
$R$-basis of $N$. If $\ker f=(\pi)$, then
\[
                            \ker u=\pi M.                \tag{WLD.12}
\]
These conclusions require no flatness of $R$ over $S$.

\begin{proof}
Tensor the explicit coordinate isomorphism $M\simeq S^q$ with $R$
and compose it with the given base-change isomorphism.
The resulting map sends the standard basis to $u(b_i)$, proving
the first assertion. If $u(x)=0$ and $x=\sum_i s_ib_i$,
the independence of $u(b_i)$ gives $f(s_i)=0$. Write
$s_i=\pi t_i$; then $x=\pi\sum_i t_ib_i$. Conversely
$f(\pi)=0$ shows $\pi M\subseteq\ker u$.
The source's \texttt{Semi\allowbreak BC} is the equivalent universal property:
the inverse to the displayed coordinate map is the unique linear
map induced by $x\mapsto(f(s_i(x)))_i$. Thus the proof applies
directly to that formulation.
\end{proof}
% END WLD10

% BEGIN WLD11
\paragraph{WLD11: the two Frobenius inclusions at a local prism.}
\label{claim:WLD11}
Fix a local prism $P_S$ with coefficient ring $S$, an element
$Q$ of its prism ideal, and a map $i:P_S\to P_{\inf}$.
Let a crystalline lattice $T$ have a realizing crystal $\mathcal C$.
Write
\[
 M=\mathcal C(P_S),\quad
 Y_0=W\otimes_{\mathbb Z_p}T,\quad
 \nu_0=j\circ\mathcal C(i):M\longrightarrow Y_0,
\]
and let $\psi:S\to W$ be its scalar map. Suppose $\psi(Q)\ne0$.
There are bases $b_1,\ldots,b_q$ of $M$ over $S$ and
$c_1,\ldots,c_q$ of $Y_0$ over $W$, with $c_i=\nu_0(b_i)$.
For the actual Frobenius $F$ on $Y_0$ one has
\[
 \begin{split}
 &\forall m\in M\quad
   \exists k\ge0,\ m'\in M:\quad
           \psi(Q)^kF(\nu_0m)=\nu_0m',\\
 &\forall m\in M\quad
   \exists k\ge0,\ s_1,\ldots,s_q\in S:\quad
           \psi(Q)^k\nu_0m=\sum_i\psi(s_i)F(c_i).
 \end{split}                                                \tag{WLD.13}
\]

\begin{proof}
A finite projective module over a local ring is finite free.
For completeness, choose lifts of a residue-field basis, use
Nakayama to obtain a surjection from a finite free module, and split
it by projectivity. Its complementary summand has zero reduction
and is finite, hence vanishes by Nakayama. Choose the resulting
basis $b_i$. Apply WLD10 to the crystal transition and to the
realization. Their composite carries $b_i$ to a $W$-basis $c_i$.

We spell out the fraction comparison. Write the localized Frobenius
of $m$ at $P_S$ as $m'/Q^k$, and the corresponding Frobenius at
$P_{\inf}$ as $n/i(Q)^h$. Fraction naturality says that, after
multiplication by a power of $i(Q)^2$,
\[
               i(Q)^h\,\mathcal C(i)(m')
                       -i(Q)^k n
\]
is zero. Apply $j$. The realization identifies
$j(n)=\psi(Q)^hF(\nu_0m)$. Consequently a power of the nonzero
element $\psi(Q)$ annihilates
$\nu_0m'-\psi(Q)^kF(\nu_0m)$. The ambient module is free over the
domain $W$, so this difference is zero. This proves the first
inclusion without cancelling an element inside the possibly
non-domain ring $S$.

For the other inclusion, write the localized Frobenius of $b_i$
as $m_i/Q^{k_i}$. Its semilinear base-change property makes those
fractions a basis of $M[1/Q]$, by WLD10. Thus the localized class
of any $m$ lies in the span of the classes of the $m_i$.
Clear the finitely many coefficient denominators. If the resulting
equality initially holds only in the localization, multiply by
one further power of $Q$ killing its discrepancy in $M$.
We obtain an actual equality
\[
                         Q^k m=\sum_i t_i m_i
\]
in $M$. Applying $\nu_0$ and using the first part for each $b_i$
gives
\[
 \psi(Q)^k\nu_0m
       =\sum_i\psi(t_iQ^{k_i})F(c_i),
\]
which is the second assertion. The case $q=0$ gives zero modules
and both assertions with exponent zero.
\end{proof}
% END WLD11

% BEGIN WLD12
\paragraph{WLD12: group action and divisibility at the specialization.}
\label{claim:WLD12}
In WLD11, suppose the source supplies maps
$g_S:P_S\to P_S$ and $e:P_S\to P_W$ such that
\[
 i\circ g_S=g_{\inf}\circ i,\qquad e\circ g_S=e,\qquad
                         \ker(e:S\to P_W.A)=(\pi).
\]
Then the realized action satisfies
\[
 g(\nu_0m)=\nu_0(\mathcal C(g_S)m),\qquad
                 \mathcal C(g_S)m-m=\pi m_g
                         \quad\text{for some }m_g\in M.    \tag{WLD.14}
\]
Thus WLD11--WLD12 give every conclusion of the original
\texttt{prismatic\allowbreak FCrystal\_\allowbreak value\_\allowbreak local\allowbreak Prism\_\allowbreak ainf\_\allowbreak witt\allowbreak Tensor}
on its supplied local prism and realization.

\begin{proof}
The first equality follows by applying the transition composition
law to $i\circ g_S=g_{\inf}\circ i$, then the Galois compatibility
of $j$. For the second, transition along $e$ annihilates
$\mathcal C(g_S)m-m$ because $e\circ g_S=e$.
That transition has the semilinear base-change property.
Apply (WLD.12) with the basis of $M$ from WLD11.
This is a divisibility statement in the original $M$, with the
original $\pi$; it is stronger than mere equality after localization.
\end{proof}
% END WLD12

% BEGIN WLD13
\paragraph{WLD13: preservation of realized images under every lattice map.}
\label{claim:WLD13}
Assume the original \texttt{Bhatt\allowbreak Scholze\allowbreak Crystalline\allowbreak Lattices}
statement. Choose any realizing crystals for crystalline lattices
$T,T'$. If $f:T\to T'$ is a $\mathbb Z_p$-linear equivariant map,
then
\[
 (1\otimes f)\bigl(\nu_{0,T}(M_T)\bigr)
                          \subseteq \nu_{0,T'}(M_{T'}).    \tag{WLD.15}
\]
In particular this holds for every original $O$-linear lattice
homomorphism, and for multiplication by every $o\in O$ on $T$.

\begin{proof}
The full-faithfulness clause of the stated crystalline-lattice input
gives a morphism $u:\mathcal C_T\to\mathcal C_{T'}$ realizing $f$.
At the local prism use $u_{P_S}(m)$. Naturality along $i$ and the
identity defining the realization of $u$ give
\[
 \nu_{0,T'}(u_{P_S}(m))
       =(1\otimes f)\nu_{0,T}(m).
\]
An $O$-linear lattice map restricts to a $\mathbb Z_p$-linear map
and remains equivariant, since its underlying function is unchanged.
The same is true for scalar multiplication by $o$.
The proof applies to the crystals already chosen for $T$ and $T'$;
it does not require replacing either crystal by a newly compatible
choice. The uniqueness clause is available, but existence is enough
for this inclusion.
\end{proof}
% END WLD13

% BEGIN WLD14
\paragraph{WLD14: the exact image for a trivial lattice.}
\label{claim:WLD14}
Under WLD13, let $T$ have trivial action and a
$\mathbb Z_p$-basis $\tau_1,\ldots,\tau_d$. Then
\[
 \nu_{0,T}(M_T)
       =\left\{\sum_j\psi(s_j)\otimes\tau_j:s_j\in S\right\}.
                                                               \tag{WLD.16}
\]
For the source's specified trivial rank-one $O$-lattice,
choose its generator $e_0$ and a $\mathbb Z_p$-basis $o_j$ of $O$.
If every $\alpha\in A$ is $\sum_j o_ja(s_j)$, WLD1 carries
(WLD.16) to the exact subset
\[
      \left\{\iota(\alpha)(1\otimes e_0):\alpha\in A\right\}
                   \subseteq B\otimes_OT.                 \tag{WLD.17}
\]

\begin{proof}
The unit crystal has value the coefficient ring on each prism,
transition the ring map, and localized Frobenius the Frobenius
ring map. These maps have the required semilinear base-change
property because the module has basis $1$. Its realization sends
$s\in S$ to $\psi(s)\otimes1$.

Each coordinate projection $T\to\mathbb Z_p$ is equivariant.
Apply WLD13 to it and to the unit crystal. Hence the corresponding
coordinate of any $\nu_{0,T}(m)$ lies in $\psi(S)$. Reassembling
the coordinates proves inclusion in the right-hand side of
(WLD.16). Conversely the map $1\mapsto\tau_j$ is equivariant;
WLD13 applied in this direction shows
$\psi(s_j)\otimes\tau_j$ belongs to the image. Additivity proves
the reverse inclusion.

For the specified rank-one lattice, its nonzero line generator and
injectivity $O\hookrightarrow LL$ identify its integral lattice
with $Oe_0$. Consequently $\tau_j=o_je_0$ is a
$\mathbb Z_p$-basis. Under WLD1,
\[
 \sum_j\psi(s_j)\otimes o_je_0
       \longmapsto
       \left(\sum_j o_jc(\psi(s_j))\right)(1\otimes e_0)
       =\iota\left(\sum_j o_ja(s_j)\right)(1\otimes e_0).
\]
The assumed coefficient decomposition of $A$ proves (WLD.17)
in both directions. This fixes the actual trivial lattice, not
only its isomorphism class.
\end{proof}
% END WLD14

\subsection{Returning the prismatic basis to the given \'etale module}

Fix the original group, series, embedding and \'etale data
$\Gamma,R,I,\mathcal D$. In ordinary notation put
\[
 A=\operatorname{LetA}(O,n),\quad E=R.\mathrm{EE},\quad
 B=R_W,\quad D_T=\mathcal D(T),\quad Y_T=B\otimes_OT.
\]
The map $j_T:D_T\to Y_T$ is injective and
$j_T(ex)=\iota(e)j_T(x)$. Its image spans $Y_T$ over $B$.
The functor has its supplied semilinear Frobenius $\varphi_D$,
$\Gamma$-action $\gamma_D$, and maps $D(f)$; $j_T$ intertwines
these with the actual tensor operators. The modules $D_T$ are
free and have the stated rank of $T$. The maps $D(f)$ already
satisfy identity, composition, addition and scalar laws.
These are precisely the original \texttt{WEt} fields and are
inputs here, not new existence assumptions.

% BEGIN WLD15
\paragraph{WLD15: fixed Frobenius tuples lie in the given \'etale module.}
\label{claim:WLD15}
Suppose $\iota$ reflects units, and $H\subset G$ maps trivially to
$\Gamma$. Fix $x_1,\ldots,x_q\in Y_T$ and a matrix $U\in M_q(E)$
with
\[
 h(x_i)=x_i\ (h\in H),\qquad
            \varphi_Y(x_j)=\sum_i\iota(U_{ij})x_i.          \tag{WLD.18}
\]
Assume the following precise coefficient-descent assertion for this
$U$: for every $r$, $F\in\operatorname{GL}_r(E)$ and
$V\in M_{r,q}(B)$, if its entries are $H$-fixed and
\[
                    \iota(F)\varphi_B(V)=V\iota(U),        \tag{WLD.19}
\]
then all its entries belong to $\iota(E)$. Then every $x_i$ lies
in $j_T(D_T)$.

\begin{proof}
If $B=0$ or $T$ has rank zero, $Y_T=0$, and zero supplies the
required preimage. Otherwise the original rank and freeness
fields give a finite $E$-basis $b_1,\ldots,b_r$ of $D_T$.
The $j_T(b_i)$ span $Y_T$, which is free of rank $r$ over $B$:
expand every $j_T(d)$ in those vectors, and use the spanning
field of $j_T$. A spanning $r$-tuple in a free module of rank $r$
is a basis. Indeed, its square coordinate matrix is surjective,
has a right inverse, and hence has invertible determinant.

Let $F$ be the matrix of $\varphi_D$ on this $E$-basis. Its
image in $B$ is the semilinear Frobenius matrix in the
$j_T(b_i)$ basis. That matrix is invertible. To see this without
assuming it, note that the fixed vectors $1\otimes t$ span $Y_T$.
For a fixed vector with coordinate column $v$ we have
$v=\iota(F)\varphi_B(v)$. Thus the columns of $\iota(F)$ span
$B^r$ and its determinant is a unit. Reflection of units makes
$\det F$ a unit in $E$.

Express $x_j=\sum_iV_{ij}j_T(b_i)$. Since $H$ acts trivially
on $\Gamma$, it fixes every $j_T(b_i)$; uniqueness of coordinates
and the first part of (WLD.18) make every $V_{ij}$ $H$-fixed.
The second part of (WLD.18), expanded in this basis, is exactly
(WLD.19). Choose $v_{ij}\in E$ with $\iota(v_{ij})=V_{ij}$.
Then $x_j=j_T(\sum_i v_{ij}b_i)$.

There is no hidden finite-rank inference when the rank field is
zero: the finite free module $Y_T$ is then zero, and injectivity
of $j_T$ forces $D_T=0$. This also covers the source convention
that natural-number finrank alone need not establish finite
generation of an arbitrary free module.
\end{proof}
% END WLD15

% BEGIN WLD16
\paragraph{WLD16: the supplied trivial object is crystalline.}
\label{claim:WLD16}
For the original period datum, every finite-dimensional trivial
representation is crystalline. In particular this applies to the
specified trivial line $\Gamma.\mathrm{one}$, including after any
allowed finite extension of coefficient fields.

\begin{proof}
Let $V$ be its underlying finite-dimensional $\mathbb Q_p$-space.
Choose a basis $v_i$. Write a tensor in
$V\otimes_{\mathbb Q_p}B_{\mathrm{cris}}$ uniquely as
$\sum_i v_i\otimes b_i$. Since the action on $V$ is trivial,
this tensor is invariant exactly when each $b_i$ is invariant.
The invariant-ring identity of the original period datum is
$B_{\mathrm{cris}}^G=K_0$. Hence
\[
            D_{\mathrm{cris}}(V)
              =V\otimes_{\mathbb Q_p}K_0
              \subseteq V\otimes_{\mathbb Q_p}B_{\mathrm{cris}}.
\]
The inclusion is injective because $K_0$ is a field and its map to
the nonzero period ring is injective. Taking dimensions gives
precisely the original crystalline dimension equality.
Finite dimension over $\mathbb Q_p$ follows from finite dimension
over the allowed finite coefficient field. This proves the entire
\texttt{is\allowbreak Crys\_\allowbreak lat\_\allowbreak of\_\allowbreak trivial} assertion on the supplied period data.
\end{proof}
% END WLD16

% BEGIN WLD17
\paragraph{WLD17: one prismatic basis inside the actual \'etale module.}
\label{claim:WLD17}
Assume also that $\iota:E\to B$ is injective, as in the original
\texttt{build\_\allowbreak prism} hypotheses.
Keep one local prism and all its maps as in WLD11--WLD14. Suppose
$a:S\hookrightarrow A$ and $\psi:S\to W$ satisfy (WLD.5),
$a(Q)=\mathcal Q$, and the image of $\mathcal Q$ is a unit in $E$.
Assume the source's fixed-tuple descent property: every $H$-fixed
finite tuple in $Y_T$ whose Frobenius is given by an invertible
matrix over $E$ belongs to $j_T(D_T)$. Here $H$ maps trivially
to $\Gamma$, and $a(g_S(t))=\gamma_A(\operatorname{res}_\Gamma g)a(t)$.
Then, for each crystalline $T$, one can choose $x_1,\ldots,x_q\in D_T$
with $q=dr_T$, where $d$ is the $\mathbb Z_p$-rank of $O$ and
$r_T$ the $O$-rank of $T$, such that
\[
 \begin{split}
 &\text{every }y\in Y_T
          \text{ is uniquely }\sum_i c(w_i)j_T(x_i),\\
 &L_T:=\kappa_T(\nu_{0,T}(M_T))
       =\left\{\sum_i c(\psi(s_i))j_T(x_i):s_i\in S\right\}.
 \end{split}                                                \tag{WLD.20}
\]
The actual fixed-tuple descent can be supplied by WLD15; it is not
a consequence of nonvanishing or of equality of ranks alone.

\begin{proof}
Take the bases $b_i,c_i$ of WLD11 and set
$y_i=\kappa_T(c_i)$. WLD1 gives their unique $W$-coordinates on
$Y_T$. Since $T$ is free of rank $r_T$ over $O$ and $O$ free
of rank $d$ over $\mathbb Z_p$, WLD2 and invariance of finite free
rank over the nonzero ring $W$ give $q=dr_T$.

The two Frobenius inclusions (WLD.13), transported by WLD1,
give precisely the hypotheses of WLD6 for the tuples
$y_i$ and $\varphi_Y(y_i)$. Thus the Frobenius matrix on this
tuple is invertible over $E$.

If $g\in H$, equivariance of $a$ gives $a(g_S(t))=a(t)$
for every $t\in S$. Its injectivity implies that $g_S$ is the
identity prism map. WLD12 and the identity law for crystal
transition maps then give $g(y_i)=y_i$. The supplied fixed-tuple
descent provides $x_i$ with $j_T(x_i)=y_i$. These preimages are
unique because $j_T$ is injective. Finally every element of $M_T$
has unique $S$-coordinates in the basis $b_i$; semilinearity
of $\nu_{0,T}$ and WLD1 give the second line of (WLD.20).
The chosen basis helps prove the statement, but the subset $L_T$
is the image of the chosen prismatic realization itself.
\end{proof}
% END WLD17

\subsection{Constructing and checking every field of the lattice}

We now prove the original lattice assembly at its full generic scope.
This also isolates exactly what the prismatic realization supplies.
Keep $A,E,B,D_T,j_T$ and all their original operators fixed.
Suppose $E\ne0$ and the following data have been supplied:
\begin{enumerate}
 \item Ring maps $Z_0[[X]]\xrightarrow{z}S\xrightarrow{a}A$,
 with $a$ injective and $S$ finite free over $Z_0[[X]]$;
 maps $c:W\to B$, $\psi:S\to W$ satisfying (WLD.5);
 $\iota:E\hookrightarrow B$ reflecting units;
 and elements $\pi,Q\in S$ with
 $a(\pi)=\pi_A$ and $a(Q)=\mathcal Q$.
 \item Elements $o_1,\ldots,o_d\in O$ such that
 $A=\sum_j o_ja(S)$, and the exact coefficient decomposition
 (WLD.7) in $E$.
 \item For every crystalline $T$, a subset $L_T\subseteq Y_T$
 and a tuple $x_i$ satisfying both assertions (WLD.20), with
 $q=dr_T$.
 \item For every crystalline $T$ and $y\in L_T$, some power
 of $c(\psi Q)$ times $\varphi_Y(y)$ lies in $L_T$, and some
 power times $y$ is a finite sum of terms
 $c(\psi t)\varphi_Y(u)$ with $t\in S$, $u\in L_T$.
 For every $g\in G$ and $y\in L_T$ there is $u\in L_T$ with
 $g(y)-y=c(\psi\pi)u$.
 \item The subset $L_T$ is preserved by every tensor map
 associated to an original homomorphism between crystalline lattices.
 \item The specified trivial object $T_1$ is crystalline, has an
 actual generator $e_0$ over $O$, and
 $L_{T_1}=\{\iota(\alpha)(1\otimes e_0):\alpha\in A\}$.
\end{enumerate}
These are the explicit hypotheses of the source's
\texttt{Wach\allowbreak Lattice.\allowbreak build}. They concern subsets of the ambient tensor
modules and finite coefficient data, not an already existing lattice
in $D_T$. In the prismatic application, WLD11--WLD17 prove the
subset assertions on one fixed prism, while WLD4 supplies the
coefficient decomposition.

% BEGIN WLD18
\paragraph{WLD18: the inverse image is a submodule over the full coefficient ring.}
\label{claim:WLD18}
For crystalline $T$ define
\[
           N(T)=\{x\in D_T:j_T(x)\in L_T\};
\]
for noncrystalline $T$ put $N(T)=0$. For crystalline $T$ this is
an $A$-submodule, and
\[
 x\in N(T)\quad\Longleftrightarrow\quad
                       x=\sum_i a(s_i)x_i\quad(s_i\in S).
                                                               \tag{WLD.21}
\]
Every element of $L_T$ has a unique preimage under $j_T$.

\begin{proof}
The second assertion of (WLD.20), semilinearity of $j_T$ and its
injectivity give (WLD.21). They also give existence and uniqueness
of the stated preimages. Formula (WLD.21) proves closure under
addition, zero and multiplication by $a(S)$.

For $o\in O$, apply the supplied map-preservation assertion to the
original lattice endomorphism $t\mapsto ot$. The scalar law of
the given \'etale functor and naturality of $j_T$ show that
$j_T(ox)$ is in $L_T$ whenever $j_T(x)$ is. Thus the inverse
image is stable under $O$ as well. Now write an arbitrary
$\alpha\in A$ as $\sum_j o_ja(s_j)$. The preceding two stability
statements and addition prove stability under $\alpha$.
This last step is essential: stability only under the smaller
ring $S$ would not supply the original \texttt{WN.\allowbreak N} field.
\end{proof}
% END WLD18

% BEGIN WLD19
\paragraph{WLD19: finite generation, actual scalar extension and projectivity.}
\label{claim:WLD19}
The submodule $N(T)$ is finitely generated over $A$ for every $T$.
For crystalline $T$ the inclusion gives the original isomorphism
\[
                        E\otimes_A N(T)\simeq D_T.         \tag{WLD.22}
\]
For every $T$, its restriction along $a\circ z$ is projective
over $Z_0[[X]]$.

\begin{proof}
For crystalline $T$, every $x_i$ lies in $N(T)$ by (WLD.21),
and that formula expresses every element as an $A$-linear
combination of the finite tuple. Conversely their $A$-span is
in $N(T)$ by WLD18. Thus they generate exactly $N(T)$.

To apply WLD8, one needs a finite $E$-basis of $D_T$ of rank $r_T$.
If $r_T>0$, the original freeness and finrank equality provide it.
If $r_T=0$, then $q=0$ and the unique-coordinate assertion
(WLD.20) gives $Y_T=0$; injectivity of $j_T$ gives $D_T=0$.
The empty basis suffices. WLD7--WLD8 now prove (WLD.22) on
the literal inclusion.

For projectivity, suppose $\sum_i a(s_i)x_i=0$. Applying $j_T$
gives $\sum_i c(\psi(s_i))j_T(x_i)=0$. Uniqueness of the
$W$-coordinates makes $\psi(s_i)=0$ for every $i$.
Then $\iota(a(s_i))=0$. The successive injective maps
$S\hookrightarrow A\hookrightarrow E\hookrightarrow B$
give $s_i=0$. Hence the $S$-coordinates in (WLD.21) are unique.
WLD9 proves that $N(T)$ is finite free, in particular projective,
over $Z_0[[X]]$ with precisely the specified restricted action.
For noncrystalline $T$, the zero module is finitely generated
and projective. No scalar-extension assertion is required in
that branch.
\end{proof}
% END WLD19

% BEGIN WLD20
\paragraph{WLD20: both original Frobenius localization conditions.}
\label{claim:WLD20}
For every $T$ and $x\in N(T)$ there is $k\ge0$ with
\[
                   \mathcal Q^k\varphi_D(x)\in N(T).
\]
For every $y\in N(T)$ there is $k\ge0$ with
\[
           \mathcal Q^ky\in
                  \operatorname{span}_A\{\varphi_D(x):x\in N(T)\}.
                                                               \tag{WLD.23}
\]
Here $\mathcal Q$ is the original
$\varphi_A^{\,s-1}(q_0)$, not a replacement denominator.

\begin{proof}
For crystalline $T$, $j_T$ intertwines Frobenius and carries
$\mathcal Q^kx$ to $c(\psi Q)^kj_T(x)$. The first subset
condition in item 4 therefore proves the first assertion.

For the second, write the corresponding ambient relation as
\[
 c(\psi Q)^kj_T(y)=
                 \sum_i c(\psi(t_i))\varphi_Y(u_i),
                         \qquad u_i\in L_T.
\]
By WLD18 lift each $u_i$ to $v_i\in N(T)$. Intertwining and
injectivity of $j_T$ give the actual equality
\[
                      \mathcal Q^ky=\sum_i a(t_i)\varphi_D(v_i)
\]
in $D_T$. This lies in the span in (WLD.23).
For noncrystalline $T$, $x=y=0$, and exponent zero works.
The exponents are permitted to depend on the vector, exactly
as in the original two \texttt{WN} fields.
\end{proof}
% END WLD20

% BEGIN WLD21
\paragraph{WLD21: group stability and trivial action modulo the fixed parameter.}
\label{claim:WLD21}
For every $\gamma\in\Gamma$ and $x\in N(T)$ there exists
$u\in N(T)$ with
\[
                        \gamma_D(x)-x=\pi_Au.             \tag{WLD.24}
\]
Consequently $N(T)$ is $\Gamma$-stable and its induced action
on $N(T)/\pi_A N(T)$ is trivial.

\begin{proof}
For crystalline $T$, use the surjectivity $G\to\Gamma$ in the
original group data to lift $\gamma$ to $g\in G$. Item 4 gives
$g(j_Tx)-j_Tx=c(\psi\pi)y$ for some $y\in L_T$.
Lift $y=j_Tu$ by WLD18. Compatibility of the action, (WLD.5)
and injectivity of $j_T$ give (WLD.24).
Since $N(T)$ is an $A$-submodule, the equality also gives
$\gamma_Dx=x+\pi_Au\in N(T)$.
For noncrystalline $T$ choose $u=0$. Notice that the divisibility
is inside $N(T)$, rather than merely inside the ambient $D_T$.
\end{proof}
% END WLD21

% BEGIN WLD22
\paragraph{WLD22: all original lattice homomorphisms preserve the construction.}
\label{claim:WLD22}
For any original lattice homomorphism $f:T\to T'$ with $T,T'$
crystalline, the supplied map $D(f)$ takes $N(T)$ into $N(T')$.
Its restrictions preserve the original identity, composition,
addition, scalar, Frobenius and $\Gamma$ relations.

\begin{proof}
If $x\in N(T)$, item 5 puts
$(1\otimes f)j_Tx$ in $L_{T'}$. The actual naturality identity
\[
                       j_{T'}D(f)x=(1\otimes f)j_Tx
\]
and the definition of $N(T')$ give the inclusion.
All the asserted map identities are equalities of restrictions
of the given maps on $D_T$; their images lie in the required
submodules by this inclusion and WLD20--WLD21. The original
Frobenius relation may be read after the prescribed localization,
where WLD20 makes the restrictions defined.
No new choices of matrices or bases are made for different $f$.
The original \texttt{map\_\allowbreak N} field only requires both lattices
to be crystalline, and we retain exactly that scope.
\end{proof}
% END WLD22

% BEGIN WLD23
\paragraph{WLD23: the specified trivial lattice and its generator.}
\label{claim:WLD23}
There is an element $e\in D_{T_1}$ with
\[
        j_{T_1}(e)=1\otimes e_0,\qquad N(T_1)=Ae,           \tag{WLD.25}
\]
where $e_0$ is the supplied generator of the original integral
trivial lattice.

\begin{proof}
Item 6 puts $1\otimes e_0$ in $L_{T_1}$, so WLD18 gives its
unique preimage $e$. An element $x$ belongs to $N(T_1)$ exactly
when $j_{T_1}x=\iota(\alpha)(1\otimes e_0)$ for some
$\alpha\in A$. Semilinearity and injectivity say exactly that
$x=\alpha e$. Thus (WLD.25) identifies the actual chosen
trivial object and its actual tensor generator.
\end{proof}
% END WLD23

% BEGIN WLD24
\paragraph{WLD24: the full original generic Wach-lattice constructor.}
\label{claim:WLD24}
Under the six displayed input conditions, the construction
$T\mapsto N(T)$ supplies \emph{every} field of
\[
                       \mathrm{WN}(\mathcal D,a\circ z).
\]
In particular it proves the original
\texttt{Global\allowbreak Theory.\allowbreak Wach\allowbreak Lattice.\allowbreak build} implication, with its
original input hypotheses.

\begin{proof}
WLD18 supplies the actual $A$-submodule. WLD19 supplies
\texttt{N\_\allowbreak base}, \texttt{N\_\allowbreak fg} and \texttt{N\_\allowbreak proj};
WLD20 supplies \texttt{N\_\allowbreak phi} and \texttt{N\_\allowbreak phiinv};
WLD21 supplies \texttt{N\_\allowbreak Gamma} and \texttt{N\_\allowbreak pi};
WLD22 supplies \texttt{map\_\allowbreak N}; and WLD23 supplies
\texttt{one\_\allowbreak e}. By definition $N(T)=0$ for noncrystalline
$T$, giving \texttt{N\_\allowbreak junk}. These assertions use the same
$N(T)$, $\pi_A$, $\mathcal Q$ and the same given $D_T$ throughout.

The source's denominator really is invertible in $E$. Indeed
$\varphi_A(\pi_A)=\pi_Aq_0$, and both $\pi_A$ and
$\varphi_E(\pi_A)$ are units there. Therefore $q_0$ is a unit;
applying $\varphi_E$ another $s-1$ times proves that
$\mathcal Q=\varphi_A^{s-1}(q_0)$ is a unit. No additional
height bound or change of parameter is used.
\end{proof}
% END WLD24

% BEGIN WLD25
\paragraph{WLD25: the full prismatic lattice assembly.}
\label{claim:WLD25}
Assume the exact original
\texttt{Bhatt\allowbreak Scholze\allowbreak Crystalline\allowbreak Lattices} input, together with the
specified local prism, its maps, and the scalar data in
\texttt{Wach\allowbreak Lattice.\allowbreak build\_\allowbreak prism}. Thus the prism ring $S$ is local
and finite free over the specified $\mathbb Z_p[[X]]$; its
specialization has kernel $(\pi)$; its ideal is $(Q)$; its
embedding $a:S\hookrightarrow A$ has $a(\pi)=\pi_A$ and
$a(Q)=\mathcal Q$; its Galois maps commute with the distinguished
embedding and act trivially on the specialization. The supplied
basis of $O$ gives the stated $S$-coordinate decomposition of
$A$, the stated $W$-coordinates of $B$, and the tensor
identifications (WLD.1). The map $\iota$ is injective and reflects
units, $c$ reflects units, and the original fixed-tuple descent
hypothesis holds for the stated $H$ with trivial image in $\Gamma$.
For every crystalline $T$, also retain the original realization
datum \texttt{hreal}: a crystal and its Witt-tensor realization,
with its $\mathbb Z_p$-linear action equal to the specified action
on $T$. This datum is an explicit premise of the generic theorem;
the all-layer specialization will supply it from the published input.
Then the original output
\[
                         \mathrm{Nonempty}\bigl(
                         \mathrm{WN}(\mathcal D,a\circ z)\bigr)
                                                               \tag{WLD.26}
\]
is proved.

\begin{proof}
Choose a realizing crystal for each crystalline $T$, as allowed
by the supplied realization data. In the all-layer application
these data come from the existence clause of the crystalline-lattice
input. Keep that choice fixed for this $T$ and set
\[
             L_T=\kappa_T\bigl(j_T^{\mathrm{crys}}
                                  (\mathcal C_T(i)M_T)\bigr).
\]
Here $j_T^{\mathrm{crys}}$ is the crystalline realization, distinct
from the given \'etale embedding $j_T$; this notation makes the
two maps explicit.

By WLD24's last paragraph, $\mathcal Q$ is a unit in $E$.
The identity $\iota(a(Q))=c(\psi Q)$ and nonzero $B$
give $\psi(Q)\ne0$. Thus WLD11--WLD12 apply.
WLD1 transfers their equations to $Y_T$, and WLD17 gives
the descended basis and (WLD.20). The two Frobenius inclusions
and the group divisibility in item 4 are exactly (WLD.13)--(WLD.14)
on the image, with the same $\pi$ and $Q$.

WLD13 proves item 5 for \emph{all} the original crystalline
lattice maps, including the scalar endomorphisms needed for
$A$-stability. WLD16 shows that the specified trivial object
is crystalline. WLD14 gives precisely item 6 for its generator.
WLD4 supplies the exact coefficient decomposition in $E$;
its $W$-completeness is the $p$-adic completeness of the original
Witt ring. The remaining scalar conditions are the displayed
input data of this theorem. All six conditions are therefore
proved, and WLD24 gives (WLD.26).

For the source's arbitrary generic coefficient ring $O$, the given
basis and injective map $O\to LL$ supply the rank-one generator
used in WLD14; no completeness or discrete valuation hypothesis
on that generic $O$ has been added. Naturality WLD13 works by
restriction to $\mathbb Z_p$ and the supplied tensor maps.
Thus this is the full original \texttt{build\_\allowbreak prism} implication,
not merely its specialization to one conveniently chosen \'etale model.
\end{proof}
% END WLD25

\subsection{The remaining all-layer input, stated without extra coherence}

% BEGIN WLD26
\paragraph{WLD26: a precise three-supplier return to all-layer Wach existence.}
\label{claim:WLD26}
Fix an original prime $p$, local period datum $P$, tied integral
datum $A_{\inf}$, and retain
\texttt{Bhatt\allowbreak Scholze\allowbreak Crystalline\allowbreak Lattices}. Suppose the following
three source suppliers have been supplied for every original
$\mathrm{Good}(K,s,LL)$ and every permitted $\Gamma,R,I$:
\begin{enumerate}
 \item The actual Wach prism and its maps in
 \texttt{exists\_\allowbreak wach\_\allowbreak prism\_\allowbreak over\_\allowbreak integers}, including the
 scalar identity $a\circ z=\mathrm{zlA}$ and both finite
 coefficient-coordinate assertions.
 \item Injectivity and reflection of units for
 $\iota=I.\mathrm{ioE}$, the corresponding two fields of
 \texttt{wach\allowbreak Inclusion\_\allowbreak injective\_\allowbreak saturated\_\allowbreak inter}.
 \item The following homogeneous case of
 \texttt{fixed\allowbreak By\allowbreak Cyclotomic\allowbreak Kernel\_\allowbreak etale\allowbreak Phi\_\allowbreak solution\_\allowbreak mem\_\allowbreak wach}:
 for all $r,q$, $F\in\operatorname{GL}_r(E)$,
 $U\in\operatorname{GL}_q(E)$ and $V\in M_{r,q}(B)$ with
 cyclotomic-kernel-fixed entries,
 \[
                 \iota(F)\varphi_B(V)=V\iota(U)
                     \quad\Longrightarrow\quad V_{ij}\in\iota(E).
 \]
\end{enumerate}
Then the full original \texttt{Wach\allowbreak Ex} holds: for every such
$\Gamma,R,I$ and \emph{every supplied} \'etale module $\mathcal D$,
there is an $\mathrm{NS}(\mathcal D)$ with every original field.

\begin{proof}
Take any permitted layer and any of its allowed \'etale modules.
The first supplier gives exactly the local prism data used by
WLD25. Choose a finite $\mathbb Z_p$-basis of the ring of integers
$O$ of $LL$. The standard finite-free structure of the integers
in a finite $p$-adic field extension supplies this basis.
WLD1 supplies the original tensor cancellation maps; WLD2 supplies
the $W$-coordinates of $B=O\otimes_{\mathbb Z_p}W$ and reflection
of units by $c$. Its Frobenius and Galois compatibilities hold on
pure tensors, as checked in WLD1.

Let $H$ be the original cyclotomic kernel. Its image in $\Gamma$
is trivial: apply the injective character $\chi_\Gamma$, use
its specified cyclotomic-character equality, and then the
injectivity $O\hookrightarrow LL$.
WLD15, using suppliers 2 and 3, gives the fixed-tuple descent
required by WLD25. The existence clause of the crystalline-lattice
input provides the realizing crystal for every crystalline $T$;
its fullness clause is exactly the one used in WLD13--WLD14.
The tied local data give the required standard-period hypothesis.
WLD25 now yields \texttt{WN} with scalar map $a\circ z$.
The identity in supplier 1 identifies that map with the original
$\mathrm{zlA}$, so the output is literally \texttt{NS}.

All choices were made after the supplied layer and \'etale module
were fixed. The universal quantifiers over those data have not
been reduced. The existing DFD12 coefficient-extension argument
may supply the crystalline-lattice input from the original
integral Bhatt--Scholze input, but it does not supply any of the
three period-ring results retained above.
\end{proof}
% END WLD26

\paragraph{What this chapter supplies, and what remains.}
The completed original mathematical producers are the coefficient
lifting theorem, the denominator-matrix theorem, the base-change
lattice theorem, the finite-free restriction theorem, the local-prism
crystalline value theorem, and both generic Wach-lattice constructors.
They now have one detailed human argument with their actual common
objects. The matrix descent in WLD15 is proved from its explicitly
stated coefficient-descent premise; WLD26 is a conditional all-layer
return with exactly the three displayed period-ring suppliers.
Those suppliers remain mathematical tasks, not declarations proved
by this chapter.

The next useful period-ring target is the original injectivity and
unit-reflection theorem, or the homogeneous fixed-Frobenius matrix
descent at the exact scope in WLD26. The prism construction is
independent of those two and can proceed in parallel. After full
\texttt{Wach\allowbreak Ex}, the faithful-invariant and \'etale-basis work,
the Herr class maps, and the quantitative de Rham comparison are
still needed for the complete \texttt{HW/\allowbreak Fon/\allowbreak Cmp} triple.
The \texttt{Layers} assembly permits a whole complete record to be
chosen at each allowed outer index; it does not add an unspecified
inter-layer transition condition. SSC's completed signature comparison
and the later local family bounds retain their original scopes.
No whole \texttt{GLOBAL\_\allowbreak DENOMINATOR} closure is claimed.

\paragraph{Source correspondence.}
Every source below is at the accepted commit:
\par\noindent\texttt{7cbf7a0c32fca1badf05931879a37c3f8a2b03b4}.
\par\noindent
The exact source ranges and their hashes are recorded with each claim.
Under \texttt{Proofs/\allowbreak Inertia\allowbreak Eigenvalues}, the primary producers are
\texttt{Exists\allowbreak Witt\allowbreak Tilt\allowbreak Lattice\allowbreak Tensor\allowbreak Cancel},
\texttt{Exists\allowbreak Adic\allowbreak Coordinates\allowbreak Of\allowbreak Basis\allowbreak Mod\allowbreak P},
\texttt{Exists\allowbreak Is\allowbreak Unit\allowbreak Matrix\allowbreak Of\allowbreak Pow\allowbreak Smul\allowbreak Span\allowbreak Invert},
\texttt{Is\allowbreak Base\allowbreak Change\allowbreak Lattice\allowbreak Of\allowbreak Basis\allowbreak After\allowbreak Extension},
\texttt{Projective\allowbreak Of\allowbreak Free\allowbreak Lattice\allowbreak Over\allowbreak Finite\allowbreak Free\allowbreak Ring},
\texttt{Prismatic\allowbreak FCrystal\allowbreak Value\allowbreak Local\allowbreak Prism\allowbreak Ainf\allowbreak Witt\allowbreak Tensor},
\texttt{Witt\allowbreak Twist\allowbreak Invariant\allowbreak Basis\allowbreak Change},
\texttt{Is\allowbreak Crys\allowbreak Lat\allowbreak Of\allowbreak Trivial}, \texttt{Is\allowbreak Unit\allowbreak Of\allowbreak One\allowbreak Tmul}, and
\texttt{Wach\allowbreak Lattices\allowbreak Exist\allowbreak All\allowbreak Layers\allowbreak Of\allowbreak Bhatt\allowbreak Scholze}.
The unchanged \texttt{WN}, \texttt{WEt} and \texttt{Layers}
contracts are in \texttt{Definitions/\allowbreak Non\allowbreak Galois\allowbreak Descent/\allowbreak Exists\allowbreak Of\allowbreak Layers};
the all-layer quantifiers are in \texttt{Conv\allowbreak Zero}.






\endgroup
\end{document}
